
\Data\ID{1003019801}
\Updated{2022-08-25 10:08:02}
\Level{A}
\SubArea{Rational equations and inequalities}
\LANGen{
\Question{
Find the domain of the expression.
\[
\frac{x^2+4x-5}{-3x^2-19x+14}
\]}
{\(\mathbb{R}\setminus\left\{-7;\frac23\right\}\)}{\(\mathbb{R}\setminus\left\{7;-\frac23\right\}\)}{\(\left(7;-\frac23\right)\)}{\(\left(-7;\frac23\right)\)}{}{}}
\LANGpl{
\Question{
Znajdź dziedzinę podanego wyrażenia.
\[
\frac{x^2+4x-5}{-3x^2-19x+14}
\]}
{\(x\in\mathbb{R}\setminus\left\{-7;\frac23\right\}\)}{\(x\in\mathbb{R}\setminus\left\{7;-\frac23\right\}\)}{\(x\in\left(7;-\frac23\right)\)}{\(x\in\left(-7;\frac23\right)\)}{}{}}
\LANGcs{
\Question{
Určete definiční obor výrazu.
\[
\frac{x^2+4x-5}{-3x^2-19x+14}
\]}
{\(\mathbb{R}\setminus\left\{-7;\frac23\right\}\)}{\(\mathbb{R}\setminus\left\{7;-\frac23\right\}\)}{\(\left(7;-\frac23\right)\)}{\(\left(-7;\frac23\right)\)}{}{}}
\LANGsk{
\Question{
Určte definičný obor výrazu.
\[
\frac{x^2+4x-5}{-3x^2-19x+14}
\]}
{\(\mathbb{R}\setminus\left\{-7;\frac23\right\}\)}{\(\mathbb{R}\setminus\left\{7;-\frac23\right\}\)}{\(\left(7;-\frac23\right)\)}{\(\left(-7;\frac23\right)\)}{}{}}
\LANGes{
\Question{
Halla el dominio de la expresión.
\[
\frac{x^2+4x-5}{-3x^2-19x+14}
\]}
{\(\mathbb{R}\setminus\left\{-7;\frac23\right\}\)}{\(\mathbb{R}\setminus\left\{7;-\frac23\right\}\)}{\(\left(7;-\frac23\right)\)}{\(\left(-7;\frac23\right)\)}{}{}}
\End

\Data\ID{1003044806}
\Updated{2022-05-23 12:05:07}
\Level{A}
\SubArea{Rational equations and inequalities}
\LANGen{
\Question{
Find the domain of the equation \( \frac{(x-1)(x+2)}{(x-1)(x+3)}=\frac1x \).}
{\( \mathbb{R}\setminus\{-3;0;1\} \)}{\( \mathbb{R}\setminus\{-3;0\} \)}{\( \mathbb{R}\setminus\{-3;-2;0\} \)}{\( \mathbb{R}\setminus\{-3;1\} \)}{}{}}
\LANGpl{
\Question{
Wyznacz dziedzinę równania \( \frac{(x-1)(x+2)}{(x-1)(x+3)}=\frac1x \).}
{\( \mathbb{R}\setminus\{-3;0;1\} \)}{\( \mathbb{R}\setminus\{-3;0\} \)}{\( \mathbb{R}\setminus\{-3;-2;0\} \)}{\( \mathbb{R}\setminus\{-3;1\} \)}{}{}}
\LANGcs{
\Question{
Určete množinu, na které má rovnice \( \frac{(x-1)(x+2)}{(x-1)(x+3)}=\frac1x \) smysl.}
{\( \mathbb{R}\setminus\{-3;0;1\} \)}{\( \mathbb{R}\setminus\{-3;0\} \)}{\( \mathbb{R}\setminus\{-3;-2;0\} \)}{\( \mathbb{R}\setminus\{-3;1\} \)}{}{}}
\LANGsk{
\Question{
Určte definičný obor rovnice \( \frac{(x-1)(x+2)}{(x-1)(x+3)}=\frac1x \).}
{\( \mathbb{R}\setminus\{-3;0;1\} \)}{\( \mathbb{R}\setminus\{-3;0\} \)}{\( \mathbb{R}\setminus\{-3;-2;0\} \)}{\( \mathbb{R}\setminus\{-3;1\} \)}{}{}}
\LANGes{
\Question{
Halla el dominio de la ecuación \( \frac{(x-1)(x+2)}{(x-1)(x+3)}=\frac1x \).}
{\( \mathbb{R}\setminus\{-3;0;1\} \)}{\( \mathbb{R}\setminus\{-3;0\} \)}{\( \mathbb{R}\setminus\{-3;-2;0\} \)}{\( \mathbb{R}\setminus\{-3;1\} \)}{}{}}
\End

\Data\ID{1003044807}
\Updated{2022-05-23 12:05:07}
\Level{A}
\SubArea{Rational equations and inequalities}
\LANGen{
\Question{
Find the domain of the equation \( \frac{x+1}{x^2-3}=\frac{x^2+x-2}{2x+8} \).}
{\( \mathbb{R}\setminus\{-4;-\sqrt3;\sqrt3\} \)}{\( \mathbb{R}\setminus\{-2;-1;1\} \)}{\( \mathbb{R}\setminus\{-4;\sqrt3\} \)}{\( \mathbb{R}\setminus\{-4\} \)}{}{}}
\LANGpl{
\Question{
Wyznacz dziedzinę równania \( \frac{x+1}{x^2-3}=\frac{x^2+x-2}{2x+8} \).}
{\( \mathbb{R}\setminus\{-4;-\sqrt3;\sqrt3\} \)}{\( \mathbb{R}\setminus\{-2;-1;1\} \)}{\( \mathbb{R}\setminus\{-4;\sqrt3\} \)}{\( \mathbb{R}\setminus\{-4\} \)}{}{}}
\LANGcs{
\Question{
Určete množinu, na které má rovnice \( \frac{x+1}{x^2-3}=\frac{x^2+x-2}{2x+8} \) smysl.}
{\( \mathbb{R}\setminus\{-4;-\sqrt3;\sqrt3\} \)}{\( \mathbb{R}\setminus\{-2;-1;1\} \)}{\( \mathbb{R}\setminus\{-4;\sqrt3\} \)}{\( \mathbb{R}\setminus\{-4\} \)}{}{}}
\LANGsk{
\Question{
Určte definičný obor rovnice \( \frac{x+1}{x^2-3}=\frac{x^2+x-2}{2x+8} \).}
{\( \mathbb{R}\setminus\{-4;-\sqrt3;\sqrt3\} \)}{\( \mathbb{R}\setminus\{-2;-1;1\} \)}{\( \mathbb{R}\setminus\{-4;\sqrt3\} \)}{\( \mathbb{R}\setminus\{-4\} \)}{}{}}
\LANGes{
\Question{
Halla el dominio de la ecuación \( \frac{x+1}{x^2-3}=\frac{x^2+x-2}{2x+8} \).}
{\( \mathbb{R}\setminus\{-4;-\sqrt3;\sqrt3\} \)}{\( \mathbb{R}\setminus\{-2;-1;1\} \)}{\( \mathbb{R}\setminus\{-4;\sqrt3\} \)}{\( \mathbb{R}\setminus\{-4\} \)}{}{}}
\End

\Data\ID{1003044808}
\Updated{2022-05-23 12:05:07}
\Level{A}
\SubArea{Rational equations and inequalities}
\LANGen{
\Question{
Find the domain of the equation \( \frac{3x+2}{\frac x3}=1 \).}
{\( \mathbb{R}\setminus \{0\} \)}{\( \mathbb{R} \)}{\( \mathbb{R}\setminus \{-\frac23\} \)}{\( \mathbb{R}\setminus \{-\frac23;0\} \)}{}{}}
\LANGpl{
\Question{
Wyznacz dziedzinę równania \( \frac{3x+2}{\frac x3}=1 \).}
{\( \mathbb{R}\setminus \{0\} \)}{\( \mathbb{R} \)}{\( \mathbb{R}\setminus \{-\frac23\} \)}{\( \mathbb{R}\setminus \{-\frac23;0\} \)}{}{}}
\LANGcs{
\Question{
Určete množinu, na které má rovnice \( \frac{3x+2}{\frac x3}=1 \) smysl.}
{\( \mathbb{R}\setminus \{0\} \)}{\( \mathbb{R} \)}{\( \mathbb{R}\setminus \{-\frac23\} \)}{\( \mathbb{R}\setminus \{-\frac23;0\} \)}{}{}}
\LANGsk{
\Question{
Určte definičný obor rovnice \( \frac{3x+2}{\frac x3}=1 \).}
{\( \mathbb{R}\setminus \{0\} \)}{\( \mathbb{R} \)}{\( \mathbb{R}\setminus \{-\frac23\} \)}{\( \mathbb{R}\setminus \{-\frac23;0\} \)}{}{}}
\LANGes{
\Question{
Halla el dominio de la ecuación \( \frac{3x+2}{\frac x3}=1 \).}
{\( \mathbb{R}\setminus \{0\} \)}{\( \mathbb{R} \)}{\( \mathbb{R}\setminus \{-\frac23\} \)}{\( \mathbb{R}\setminus \{-\frac23;0\} \)}{}{}}
\End

\Data\ID{1003044809}
\Updated{2022-05-23 12:05:07}
\Level{A}
\SubArea{Rational equations and inequalities}
\LANGen{
\Question{
Find the domain of the equation \( \frac{2x+5}{x-2}=\frac1{x^2-1} \).}
{\( \mathbb{R}\setminus\{-1;1;2\} \)}{\( \mathbb{R}\setminus\{1;2\} \)}{\( \mathbb{R}\setminus\{-\frac52;-1;1\} \)}{\( \mathbb{R}\setminus\{2\} \)}{}{}}
\LANGpl{
\Question{
Wyznacz dziedzinę równania \( \frac{2x+5}{x-2}=\frac1{x^2-1} \).}
{\( \mathbb{R}\setminus\{-1;1;2\} \)}{\( \mathbb{R}\setminus\{1;2\} \)}{\( \mathbb{R}\setminus\{-\frac52;-1;1\} \)}{\( \mathbb{R}\setminus\{2\} \)}{}{}}
\LANGcs{
\Question{
Určete množinu, na které má rovnice \( \frac{2x+5}{x-2}=\frac1{x^2-1} \) smysl.}
{\( \mathbb{R}\setminus\{-1;1;2\} \)}{\( \mathbb{R}\setminus\{1;2\} \)}{\( \mathbb{R}\setminus\{-\frac52;-1;1\} \)}{\( \mathbb{R}\setminus\{2\} \)}{}{}}
\LANGsk{
\Question{
Určte definičný obor rovnice \( \frac{2x+5}{x-2}=\frac1{x^2-1} \).}
{\( \mathbb{R}\setminus\{-1;1;2\} \)}{\( \mathbb{R}\setminus\{1;2\} \)}{\( \mathbb{R}\setminus\{-\frac52;-1;1\} \)}{\( \mathbb{R}\setminus\{2\} \)}{}{}}
\LANGes{
\Question{
Halla el dominio de la ecuación \( \frac{2x+5}{x-2}=\frac1{x^2-1} \).}
{\( \mathbb{R}\setminus\{-1;1;2\} \)}{\( \mathbb{R}\setminus\{1;2\} \)}{\( \mathbb{R}\setminus\{-\frac52;-1;1\} \)}{\( \mathbb{R}\setminus\{2\} \)}{}{}}
\End

\Data\ID{1003044810}
\Updated{2022-05-23 12:05:07}
\Level{A}
\SubArea{Rational equations and inequalities}
\LANGen{
\Question{
Find the domain of the equation  \( \frac{x+3}{x^2-8}=\frac1x \).}
{\( \mathbb{R}\setminus\{-2\sqrt2;0;2\sqrt2\} \)}{\( \mathbb{R}\setminus\{0\} \)}{\( \mathbb{R}\setminus\{-3\} \)}{\( \mathbb{R}\setminus\{0;4\} \)}{}{}}
\LANGpl{
\Question{
Wyznacz dziedzinę równania \( \frac{x+3}{x^2-8}=\frac1x \).}
{\( \mathbb{R}\setminus\{-2\sqrt2;0;2\sqrt2\} \)}{\( \mathbb{R}\setminus\{0\} \)}{\( \mathbb{R}\setminus\{-3\} \)}{\( \mathbb{R}\setminus\{0;4\} \)}{}{}}
\LANGcs{
\Question{
Určete množinu, na které má rovnice \( \frac{x+3}{x^2-8}=\frac1x \) smysl.}
{\( \mathbb{R}\setminus\{-2\sqrt2;0;2\sqrt2\} \)}{\( \mathbb{R}\setminus\{0\} \)}{\( \mathbb{R}\setminus\{-3\} \)}{\( \mathbb{R}\setminus\{0;4\} \)}{}{}}
\LANGsk{
\Question{
Určte definičný obor rovnice  \( \frac{x+3}{x^2-8}=\frac1x \).}
{\( \mathbb{R}\setminus\{-2\sqrt2;0;2\sqrt2\} \)}{\( \mathbb{R}\setminus\{0\} \)}{\( \mathbb{R}\setminus\{-3\} \)}{\( \mathbb{R}\setminus\{0;4\} \)}{}{}}
\LANGes{
\Question{
Halla el dominio de la ecuación \( \frac{x+3}{x^2-8}=\frac1x \).}
{\( \mathbb{R}\setminus\{-2\sqrt2;0;2\sqrt2\} \)}{\( \mathbb{R}\setminus\{0\} \)}{\( \mathbb{R}\setminus\{-3\} \)}{\( \mathbb{R}\setminus\{0;4\} \)}{}{}}
\End

\Data\ID{1003119001}
\Updated{2020-08-22 06:08:14}
\Level{A}
\SubArea{Rational equations and inequalities}
\LANGen{
\Question{
Find the solution set of the following equation.
\[ \frac{x+1}{x-1}=0 \]}
{\( \{-1\} \)}{\( \{-1;1\} \)}{\( \{-1;0;1\} \)}{\( \emptyset \)}{}{}}
\LANGpl{
\Question{
Wyznacz zbiór rozwiązań podanego równania.
\[ \frac{x+1}{x-1}=0 \]}
{\( \{-1\} \)}{\( \{-1;1\} \)}{\( \{-1;0;1\} \)}{\( \emptyset \)}{}{}}
\LANGcs{
\Question{
Určete množinu řešení dané rovnice.
\[ \frac{x+1}{x-1}=0 \]}
{\( \{-1\} \)}{\( \{-1;1\} \)}{\( \{-1;0;1\} \)}{\( \emptyset \)}{}{}}
\LANGsk{
\Question{
Určte množinu riešení danej rovnice.
\[ \frac{x+1}{x-1}=0 \]}
{\( \{-1\} \)}{\( \{-1;1\} \)}{\( \{-1;0;1\} \)}{\( \emptyset \)}{}{}}
\LANGes{
\Question{
Encuentra el conjunto de soluciones de la siguiente ecuación.
\[ \frac{x+1}{x-1}=0 \]}
{\( \{-1\} \)}{\( \{-1;1\} \)}{\( \{-1;0;1\} \)}{\( \emptyset \)}{}{}}
\End

\Data\ID{1003119002}
\Updated{2022-04-23 05:04:48}
\Level{A}
\SubArea{Rational equations and inequalities}
\LANGen{
\Question{
Find the solution set of the following equation.
\[ \frac{x^2+4x+4}{x+2}=0 \]}
{\( \emptyset \)}{\( \{-2\} \)}{\( \{-2;2\} \)}{\( \{2\} \)}{}{}}
\LANGpl{
\Question{
Wyznacz zbiór rozwiązań podanego równania.
\[ \frac{x^2+4x+4}{x+2}=0 \]}
{\( \emptyset \)}{\( \{-2\} \)}{\( \{-2;2\} \)}{\( \{2\} \)}{}{}}
\LANGcs{
\Question{
Určete množinu řešení dané rovnice.
\[ \frac{x^2+4x+4}{x+2}=0 \]}
{\( \emptyset \)}{\( \{-2\} \)}{\( \{-2;2\} \)}{\( \{2\} \)}{}{}}
\LANGsk{
\Question{
Určte množinu riešení danej rovnice.
\[ \frac{x^2+4x+4}{x+2}=0 \]}
{\( \emptyset \)}{\( \{-2\} \)}{\( \{-2;2\} \)}{\( \{2\} \)}{}{}}
\LANGes{
\Question{
Encuentra el conjunto de soluciones de la siguiente ecuación.
\[ \frac{x^2+4x+4}{x+2}=0 \]}
{\( \emptyset \)}{\( \{-2\} \)}{\( \{-2;2\} \)}{\( \{2\} \)}{}{}}
\End

\Data\ID{1003119003}
\Updated{2022-04-23 05:04:48}
\Level{A}
\SubArea{Rational equations and inequalities}
\LANGen{
\Question{
Find the solution set of the following equation.
\[ \frac{3x^2-27}{x+3}=0 \]}
{\( \{ 3 \} \)}{\( \{ -3;3 \} \)}{\( \{ 9 \} \)}{\( \emptyset \)}{}{}}
\LANGpl{
\Question{
Wyznacz zbiór rozwiązań podanego równania.
\[ \frac{3x^2-27}{x+3}=0 \]}
{\( \{ 3 \} \)}{\( \{ -3;3 \} \)}{\( \{ 9 \} \)}{\( \emptyset \)}{}{}}
\LANGcs{
\Question{
Určete množinu řešení dané rovnice.
\[ \frac{3x^2-27}{x+3}=0 \]}
{\( \{ 3 \} \)}{\( \{ -3;3 \} \)}{\( \{ 9 \} \)}{\( \emptyset \)}{}{}}
\LANGsk{
\Question{
Určte množinu riešení danej rovnice.
\[ \frac{3x^2-27}{x+3}=0 \]}
{\( \{ 3 \} \)}{\( \{ -3;3 \} \)}{\( \{ 9 \} \)}{\( \emptyset \)}{}{}}
\LANGes{
\Question{
Encuentra el conjunto de soluciones de la siguiente ecuación.
\[ \frac{3x^2-27}{x+3}=0 \]}
{\( \{ 3 \} \)}{\( \{ -3;3 \} \)}{\( \{ 9 \} \)}{\( \emptyset \)}{}{}}
\End

\Data\ID{1003119004}
\Updated{2022-04-23 05:04:48}
\Level{A}
\SubArea{Rational equations and inequalities}
\LANGen{
\Question{
Find the solution set of the following equation.
\[ \frac1{4x^2-4}=0 \]}
{\( \emptyset \)}{\( \{-1;1\} \)}{\( \{1\} \)}{\( \{-1\} \)}{}{}}
\LANGpl{
\Question{
Wyznacz zbiór rozwiązań podanego równania.
\[ \frac1{4x^2-4}=0 \]}
{\( \emptyset \)}{\( \{-1;1\} \)}{\( \{1\} \)}{\( \{-1\} \)}{}{}}
\LANGcs{
\Question{
Určete množinu řešení dané rovnice.
\[ \frac1{4x^2-4}=0 \]}
{\( \emptyset \)}{\( \{-1;1\} \)}{\( \{1\} \)}{\( \{-1\} \)}{}{}}
\LANGsk{
\Question{
Určte množinu riešení danej rovnice.
\[ \frac1{4x^2-4}=0 \]}
{\( \emptyset \)}{\( \{-1;1\} \)}{\( \{1\} \)}{\( \{-1\} \)}{}{}}
\LANGes{
\Question{
Encuentra el conjunto de soluciones de la siguiente ecuación.
\[ \frac1{4x^2-4}=0 \]}
{\( \emptyset \)}{\( \{-1;1\} \)}{\( \{1\} \)}{\( \{-1\} \)}{}{}}
\End

\Data\ID{1003119005}
\Updated{2020-08-22 06:08:03}
\Level{A}
\SubArea{Rational equations and inequalities}
\LANGen{
\Question{
Find the sum of all solutions of the following equation.
\[ \frac{3(x-1)(x+4)}{(x-3)}=0 \]}
{\( -3 \)}{\( 3 \)}{\( 0 \)}{\( -1 \)}{}{}}
\LANGpl{
\Question{
Wyznacz sumę wszystkich rozwiązań podanego równania.
\[ \frac{3(x-1)(x+4)}{(x-3)}=0 \]}
{\( -3 \)}{\( 3 \)}{\( 0 \)}{\( -1 \)}{}{}}
\LANGcs{
\Question{
Určete součet všech řešení dané rovnice.
\[ \frac{3(x-1)(x+4)}{(x-3)}=0 \]}
{\( -3 \)}{\( 3 \)}{\( 0 \)}{\( -1 \)}{}{}}
\LANGsk{
\Question{
Určte súčet všetkých riešení danej rovnice.
\[ \frac{3(x-1)(x+4)}{(x-3)}=0 \]}
{\( -3 \)}{\( 3 \)}{\( 0 \)}{\( -1 \)}{}{}}
\LANGes{
\Question{
Encuentra la suma de todas las soluciones de la siguiente ecuación. \[ \frac{3(x-1)(x+4)}{(x-3)}=0 \]}
{\( -3 \)}{\( 3 \)}{\( 0 \)}{\( -1 \)}{}{}}
\End

\Data\ID{1003119006}
\Updated{2020-08-22 06:08:28}
\Level{A}
\SubArea{Rational equations and inequalities}
\LANGen{
\Question{
Find the solution set of the following equation.
\[ \frac3{x-1}+\frac2{x+2}=0 \]}
{\( \left\{-\frac45\right\} \)}{\( \emptyset \)}{\( \left\{-2;1\right\} \)}{\( \left\{-\frac54\right\} \)}{}{}}
\LANGpl{
\Question{
Wyznacz zbiór rozwiązań podanego równania.
\[ \frac3{x-1}+\frac2{x+2}=0 \]}
{\( \left\{-\frac45\right\} \)}{\( \emptyset \)}{\( \left\{-2;1\right\} \)}{\( \left\{-\frac54\right\} \)}{}{}}
\LANGcs{
\Question{
Určete množinu řešení dané rovnice.
\[ \frac3{x-1}+\frac2{x+2}=0 \]}
{\( \left\{-\frac45\right\} \)}{\( \emptyset \)}{\( \left\{-2;1\right\} \)}{\( \left\{-\frac54\right\} \)}{}{}}
\LANGsk{
\Question{
Určte množinu riešení danej rovnice.
\[ \frac3{x-1}+\frac2{x+2}=0 \]}
{\( \left\{-\frac45\right\} \)}{\( \emptyset \)}{\( \left\{-2;1\right\} \)}{\( \left\{-\frac54\right\} \)}{}{}}
\LANGes{
\Question{
Encuentra el conjunto de soluciones de la siguiente ecuación.
\[ \frac3{x-1}+\frac2{x+2}=0 \]}
{\( \left\{-\frac45\right\} \)}{\( \emptyset \)}{\( \left\{-2;1\right\} \)}{\( \left\{-\frac54\right\} \)}{}{}}
\End

\Data\ID{1003119007}
\Updated{2020-08-22 09:08:14}
\Level{A}
\SubArea{Rational equations and inequalities}
\LANGen{
\Question{
Find the sum of all solutions of the following equation.
\[ \frac{(x-4)+(x+2)}{(x+2)}=0 \]}
{\( 1 \)}{\( 4 \)}{\( 5 \)}{\( 3 \)}{}{}}
\LANGpl{
\Question{
Wyznacz sumę wszystkich rozwiązań podanego równania.
\[ \frac{(x-4)+(x+2)}{(x+2)}=0 \]}
{\( 1 \)}{\( 4 \)}{\( 5 \)}{\( 3 \)}{}{}}
\LANGcs{
\Question{
Určete součet všech řešení dané rovnice.
\[ \frac{(x-4)+(x+2)}{(x+2)}=0 \]}
{\( 1 \)}{\( 4 \)}{\( 5 \)}{\( 3 \)}{}{}}
\LANGsk{
\Question{
Určte súčet všetkých riešení danej rovnice.
\[ \frac{(x-4)+(x+2)}{(x+2)}=0 \]}
{\( 1 \)}{\( 4 \)}{\( 5 \)}{\( 3 \)}{}{}}
\LANGes{
\Question{
Encuentra la suma de todas las soluciones de la siguiente ecuación. 
\[ \frac{(x-4)+(x+2)}{(x+2)}=0 \]}
{\( 1 \)}{\( 4 \)}{\( 5 \)}{\( 3 \)}{}{}}
\End

\Data\ID{1003119008}
\Updated{2020-08-22 09:08:34}
\Level{A}
\SubArea{Rational equations and inequalities}
\LANGen{
\Question{
Find the solution set of the following equation.
\[ \frac{x-1}{x+3}\cdot\frac{x+3}{x-1}=0 \]}
{\( \emptyset \)}{\( \mathbb{R} \)}{\( \{-3;1\} \)}{\( \{-1;3\} \)}{}{}}
\LANGpl{
\Question{
Wyznacz zbiór rozwiązań podanego równania.
\[ \frac{x-1}{x+3}\cdot\frac{x+3}{x-1}=0 \]}
{\( \emptyset \)}{\( \mathbb{R} \)}{\( \{-3;1\} \)}{\( \{-1;3\} \)}{}{}}
\LANGcs{
\Question{
Určete množinu řešení dané rovnice.
\[ \frac{x-1}{x+3}\cdot\frac{x+3}{x-1}=0 \]}
{\( \emptyset \)}{\( \mathbb{R} \)}{\( \{-3;1\} \)}{\( \{-1;3\} \)}{}{}}
\LANGsk{
\Question{
Určte množinu riešení danej rovnice.
\[ \frac{x-1}{x+3}\cdot\frac{x+3}{x-1}=0 \]}
{\( \emptyset \)}{\( \mathbb{R} \)}{\( \{-3;1\} \)}{\( \{-1;3\} \)}{}{}}
\LANGes{
\Question{
Encuentra el conjunto de soluciones de la siguiente ecuación.
\[ \frac{x-1}{x+3}\cdot\frac{x+3}{x-1}=0 \]}
{\( \emptyset \)}{\( \mathbb{R} \)}{\( \{-3;1\} \)}{\( \{-1;3\} \)}{}{}}
\End

\Data\ID{1003119101}
\Updated{2020-08-22 09:08:57}
\Level{A}
\SubArea{Rational equations and inequalities}
\LANGen{
\Question{
Find the sum of all solutions of the following equation.
\[ \frac{3x^2+3x-18}{2x+6}=0 \]}
{\( 2 \)}{\( -1 \)}{\( 5 \)}{\( 1 \)}{}{}}
\LANGpl{
\Question{
Wyznacz sumę wszystkich rozwiązań podanego równania.
\[ \frac{3x^2+3x-18}{2x+6}=0 \]}
{\( 2 \)}{\( -1 \)}{\( 5 \)}{\( 1 \)}{}{}}
\LANGcs{
\Question{
Určete součet všech řešení dané rovnice.
\[ \frac{3x^2+3x-18}{2x+6}=0 \]}
{\( 2 \)}{\( -1 \)}{\( 5 \)}{\( 1 \)}{}{}}
\LANGsk{
\Question{
Určte súčet všetkých riešení danej rovnice.
\[ \frac{3x^2+3x-18}{2x+6}=0 \]}
{\( 2 \)}{\( -1 \)}{\( 5 \)}{\( 1 \)}{}{}}
\LANGes{
\Question{
Encuentra la suma de todas las soluciones de la siguiente ecuación. 
\[ \frac{3x^2+3x-18}{2x+6}=0 \]}
{\( 2 \)}{\( -1 \)}{\( 5 \)}{\( 1 \)}{}{}}
\End

\Data\ID{1003119102}
\Updated{2020-08-22 09:08:38}
\Level{A}
\SubArea{Rational equations and inequalities}
\LANGen{
\Question{
Find the sum of all solutions of the following equation.
\[ \frac{x^2+3x+2}{4x^2-4}=1 \]}
{\( 2 \)}{\( 1 \)}{\( 0 \)}{\( -1 \)}{}{}}
\LANGpl{
\Question{
Wyznacz sumę wszystkich rozwiązań podanego równania.
\[ \frac{x^2+3x+2}{4x^2-4}=1 \]}
{\( 2 \)}{\( 1 \)}{\( 0 \)}{\( -1 \)}{}{}}
\LANGcs{
\Question{
Určete součet všech řešení dané rovnice.
\[ \frac{x^2+3x+2}{4x^2-4}=1 \]}
{\( 2 \)}{\( 1 \)}{\( 0 \)}{\( -1 \)}{}{}}
\LANGsk{
\Question{
Určte súčet všetkých riešení danej rovnice.
\[ \frac{x^2+3x+2}{4x^2-4}=1 \]}
{\( 2 \)}{\( 1 \)}{\( 0 \)}{\( -1 \)}{}{}}
\LANGes{
\Question{
Encuentra la suma de todas las soluciones de la siguiente ecuación. 
\[ \frac{x^2+3x+2}{4x^2-4}=1 \]}
{\( 2 \)}{\( 1 \)}{\( 0 \)}{\( -1 \)}{}{}}
\End

\Data\ID{1003119103}
\Updated{2020-08-22 09:08:19}
\Level{A}
\SubArea{Rational equations and inequalities}
\LANGen{
\Question{
Find the sum of all solutions of the following equation.
\[ \frac{\left(x^2-1\right)\left(x^2+2x+1\right)}{(x+1)^2\cdot(x-1)^2}=x+1 \]}
{\( 2 \)}{\( 0 \)}{\( -2 \)}{\( -3 \)}{}{}}
\LANGpl{
\Question{
Wyznacz sumę wszystkich rozwiązań podanego równania.
\[ \frac{\left(x^2-1\right)\left(x^2+2x+1\right)}{(x+1)^2\cdot(x-1)^2}=x+1 \]}
{\( 2 \)}{\( 0 \)}{\( -2 \)}{\( -3 \)}{}{}}
\LANGcs{
\Question{
Určete součet všech řešení dané rovnice.
\[ \frac{\left(x^2-1\right)\left(x^2+2x+1\right)}{(x+1)^2\cdot(x-1)^2}=x+1 \]}
{\( 2 \)}{\( 0 \)}{\( -2 \)}{\( -3 \)}{}{}}
\LANGsk{
\Question{
Určte súčet všetkých riešení danej rovnice.
\[ \frac{\left(x^2-1\right)\left(x^2+2x+1\right)}{(x+1)^2\cdot(x-1)^2}=x+1 \]}
{\( 2 \)}{\( 0 \)}{\( -2 \)}{\( -3 \)}{}{}}
\LANGes{
\Question{
Encuentra la suma de todas las soluciones de la siguiente ecuación. 
\[ \frac{\left(x^2-1\right)\left(x^2+2x+1\right)}{(x+1)^2\cdot(x-1)^2}=x+1 \]}
{\( 2 \)}{\( 0 \)}{\( -2 \)}{\( -3 \)}{}{}}
\End

\Data\ID{1003119104}
\Updated{2020-08-22 09:08:54}
\Level{A}
\SubArea{Rational equations and inequalities}
\LANGen{
\Question{
Find the sum of all solutions of the following equation.
\[ \frac{x^2+5x+6}{x^2+2x-3}=0 \]}
{\( -2 \)}{\( -1 \)}{\( 2 \)}{\( 1 \)}{}{}}
\LANGpl{
\Question{
Wyznacz sumę wszystkich rozwiązań podanego równania.
\[ \frac{x^2+5x+6}{x^2+2x-3}=0 \]}
{\( -2 \)}{\( -1 \)}{\( 2 \)}{\( 1 \)}{}{}}
\LANGcs{
\Question{
Určete součet všech řešení dané rovnice.
\[ \frac{x^2+5x+6}{x^2+2x-3}=0 \]}
{\( -2 \)}{\( -1 \)}{\( 2 \)}{\( 1 \)}{}{}}
\LANGsk{
\Question{
Určte súčet všetkých riešení danej rovnice.
\[ \frac{x^2+5x+6}{x^2+2x-3}=0 \]}
{\( -2 \)}{\( -1 \)}{\( 2 \)}{\( 1 \)}{}{}}
\LANGes{
\Question{
Encuentra la suma de todas las soluciones de la siguiente ecuación. 
\[ \frac{x^2+5x+6}{x^2+2x-3}=0 \]}
{\( -2 \)}{\( -1 \)}{\( 2 \)}{\( 1 \)}{}{}}
\End

\Data\ID{1003119105}
\Updated{2020-08-22 09:08:19}
\Level{A}
\SubArea{Rational equations and inequalities}
\LANGen{
\Question{
Find the sum of all solutions of the following equation.
\[ \frac{3(x-1)(x+4)}{(x+4)}=0 \]}
{\( 1 \)}{\( -3 \)}{\( -5 \)}{\( -1 \)}{}{}}
\LANGpl{
\Question{
Wyznacz sumę wszystkich rozwiązań następującego równania. 
\[ \frac{3(x-1)(x+4)}{(x+4)}=0 \]}
{\( 1 \)}{\( -3 \)}{\( -5 \)}{\( -1 \)}{}{}}
\LANGcs{
\Question{
Určete součet všech řešení dané rovnice.
\[ \frac{3(x-1)(x+4)}{(x+4)}=0 \]}
{\( 1 \)}{\( -3 \)}{\( -5 \)}{\( -1 \)}{}{}}
\LANGsk{
\Question{
Určte súčet všetkých riešení danej rovnice.
\[ \frac{3(x-1)(x+4)}{(x+4)}=0 \]}
{\( 1 \)}{\( -3 \)}{\( -5 \)}{\( -1 \)}{}{}}
\LANGes{
\Question{
Encuentra la suma de todas las soluciones de la siguiente ecuación. 
\[ \frac{3(x-1)(x+4)}{(x+4)}=0 \]}
{\( 1 \)}{\( -3 \)}{\( -5 \)}{\( -1 \)}{}{}}
\End

\Data\ID{1003136401}
\Updated{2020-08-22 09:08:24}
\Level{A}
\SubArea{Rational equations and inequalities}
\LANGen{
\Question{
Choose the operation, which most effectively eliminates fractions from the equation. 
\[ 3+\frac2{x+4}=\frac1{3x+12} \]}
{multiplying both sides by \( 3x+12 \)}{multiplying both sides by \( (x+4)(3x+12) \)}{subtracting \( \frac2{x+4} \) from both sides}{multiplying both sides by \( 12x \)}{}{}}
\LANGpl{
\Question{
Wybierz działanie, które najbardziej efektywnie usunie ułamki z równania.
\[ 3+\frac2{x+4}=\frac1{3x+12} \]}
{pomnożyć obydwie strony przez \( 3x+12 \)}{pomnożyć obydwie strony przez \( (x+4)(3x+12) \)}{odejmując \( \frac2{x+4} \)  z obydwu stron}{pomnożyć obydwie strony \( 12x \)}{}{}}
\LANGcs{
\Question{
Vyberte nejvhodnější ekvivalentní úpravu, pomocí které odstraníte z rovnice lomené výrazy.
\[ 3+\frac2{x+4}=\frac1{3x+12} \]}
{vynásobení obou stran rovnice výrazem  \( 3x+12 \)}{vynásobení obou stran rovnice výrazem  \( (x+4)(3x+12) \)}{odečtení výrazu  \( \frac2{x+4} \) od obou stran rovnice}{vynásobení obou stran rovnice výrazem  \( 12x \)}{}{}}
\LANGsk{
\Question{
Vyberte najvhodnejšiu ekvivalentnú úpravu, pomocou ktorej odstránite z rovnice lomené výrazy. 
\[ 3+\frac2{x+4}=\frac1{3x+12} \]}
{vynásobenie obidvoch strán rovnice výrazom \( 3x+12 \)}{vynásobenie obidvoch strán rovnice výrazom \( (x+4)(3x+12) \)}{odčítanie výrazu \( \frac2{x+4} \) od obidvoch strán rovnice}{vynásobenie obidvoch strán rovnice výrazom \( 12x \)}{}{}}
\LANGes{
\Question{
Elige la operación matemática más apropiada para eliminar la fracción de la ecuación.  
\[ 3+\frac2{x+4}=\frac1{3x+12} \]}
{multiplicar los dos lados por \( 3x+12 \)}{multiplicar los dos lados por \( (x+4)(3x+12) \)}{restar \( \frac2{x+4} \) de los dos lados}{multiplicar los dos lados por \( 12x \)}{}{}}
\End

\Data\ID{1003136402}
\Updated{2020-08-22 10:08:08}
\Level{A}
\SubArea{Rational equations and inequalities}
\LANGen{
\Question{
Choose the operation, which most effectively eliminates fractions from the equation.
\[ \frac2{x^2-9}+\frac3{3-x}=\frac{x+1}{2x} \]}
{multiplying both sides by \( 2x\left(x^2-9\right) \)}{multiplying both sides by \( 2x\left(x^2-9\right)(3-x) \)}{multiplying both sides by \( 2x^2-9 \)}{multiplying both sides by \( 18x^2 \)}{}{}}
\LANGpl{
\Question{
Wybierz działanie, które najbardziej efektywnie usunie ułamki z równania.
\[ \frac2{x^2-9}+\frac3{3-x}=\frac{x+1}{2x} \]}
{pomnożyć obydwie strony przez \( 2x\left(x^2-9\right) \)}{pomnożyć obydwie sprawy przez \( 2x\left(x^2-9\right)(3-x) \)}{pomnożyć obydwie strony przez \( 2x^2-9 \)}{pomnożyć obydwie strony przez \( 18x^2 \)}{}{}}
\LANGcs{
\Question{
Vyberte nejvhodnější ekvivalentní úpravu, pomocí které odstraníte z rovnice lomené výrazy.
\[ \frac2{x^2-9}+\frac3{3-x}=\frac{x+1}{2x} \]}
{vynásobení obou stran rovnice výrazem \( 2x\left(x^2-9\right) \)}{vynásobení obou stran rovnice výrazem \( 2x\left(x^2-9\right)(3-x) \)}{vynásobení obou stran rovnice výrazem \( 2x^2-9 \)}{vynásobení obou stran rovnice výrazem \( 18x^2 \)}{}{}}
\LANGsk{
\Question{
Vyberte najvhodnejšiu ekvivalentnú úpravu, pomocou ktorej odstránite z rovnice lomené výrazy.
\[ \frac2{x^2-9}+\frac3{3-x}=\frac{x+1}{2x} \]}
{vynásobenie obidvoch strán rovnice výrazom \( 2x\left(x^2-9\right) \)}{vynásobenie obidvoch strán rovnice výrazom \( 2x\left(x^2-9\right)(3-x) \)}{vynásobenie obidvoch strán rovnice výrazom \( 2x^2-9 \)}{vynásobenie obidvoch strán rovnice výrazom \( 18x^2 \)}{}{}}
\LANGes{
\Question{
Elige la operación matemática más apropiada para eliminar la fracción de la ecuación. 
\[ \frac2{x^2-9}+\frac3{3-x}=\frac{x+1}{2x} \]}
{multiplicar los dos lados por \( 2x\left(x^2-9\right) \)}{multiplicar los dos lados por \( 2x\left(x^2-9\right)(3-x) \)}{multiplicar los dos lados por \( 2x^2-9 \)}{multiplicar los dos lados por \( 18x^2 \)}{}{}}
\End

\Data\ID{1003136403}
\Updated{2020-08-22 10:08:19}
\Level{A}
\SubArea{Rational equations and inequalities}
\LANGen{
\Question{
Choose the operation, which most effectively eliminates fractions from the equation.
\[ \frac{2x}{x^2-25}+\frac{3+x}{5-x}=\frac{x+1}{x+5} \]}
{multiplying both sides by \( x^2-25 \)}{multiplying both sides by \( (5-x)\left(x^2-25\right) \)}{multiplying both sides by \( x^2+25 \)}{multiplying both sides by \( (5-x)(x+5)\left(x^2-25\right) \)}{}{}}
\LANGpl{
\Question{
Wybierz działanie, które najbardziej efektywnie usunie ułamki z równania.
\[ \frac{2x}{x^2-25}+\frac{3+x}{5-x}=\frac{x+1}{x+5} \]}
{pomnożyć obydwie strony przez \( x^2-25 \)}{pomnożyć obydwie strony przez \( (5-x)\left(x^2-25\right) \)}{pomnożyć obydwie strony przez \( x^2+25 \)}{pomnożyć obydwie strony przez \( (5-x)(x+5)\left(x^2-25\right) \)}{}{}}
\LANGcs{
\Question{
Vyberte nejvhodnější ekvivalentní úpravu, pomocí které odstraníte z rovnice lomené výrazy.
\[ \frac{2x}{x^2-25}+\frac{3+x}{5-x}=\frac{x+1}{x+5} \]}
{vynásobení obou stran rovnice výrazem \( x^2-25 \)}{vynásobení obou stran rovnice výrazem \( (5-x)\left(x^2-25\right) \)}{vynásobení obou stran rovnice výrazem \( x^2+25 \)}{vynásobení obou stran rovnice výrazem \( (5-x)(x+5)\left(x^2-25\right) \)}{}{}}
\LANGsk{
\Question{
Vyberte najvhodnejšiu ekvivalentnú úpravu, pomocou ktorej odstránite z rovnice lomené výrazy.
\[ \frac{2x}{x^2-25}+\frac{3+x}{5-x}=\frac{x+1}{x+5} \]}
{vynásobenie obidvoch strán rovnice výrazom \( x^2-25 \)}{vynásobenie obidvoch strán rovnice výrazom \( (5-x)\left(x^2-25\right) \)}{vynásobenie obidvoch strán rovnice výrazom \( x^2+25 \)}{vynásobenie obidvoch strán rovnice výrazom \( (5-x)(x+5)\left(x^2-25\right) \)}{}{}}
\LANGes{
\Question{
Elige la operación matemática más apropiada para eliminar la fracción de la ecuación. 
\[ \frac{2x}{x^2-25}+\frac{3+x}{5-x}=\frac{x+1}{x+5} \]}
{multiplicar los dos lados por \( x^2-25 \)}{multiplicar los dos lados por \( (5-x)\left(x^2-25\right) \)}{multiplicar los dos lados por \( x^2+25 \)}{multiplicar los dos lados por \( (5-x)(x+5)\left(x^2-25\right) \)}{}{}}
\End

\Data\ID{1003136404}
\Updated{2020-08-22 10:08:21}
\Level{A}
\SubArea{Rational equations and inequalities}
\LANGen{
\Question{
Choose the resulting form of the given equation after multiplying both sides by  \( x^2-16 \).
\[ \frac x{x-4}+x+2=\frac{3+x}{x+4} \]}
{\( x(x+4)+(x^2-16)(x+2)=(3+x)(x-4) \)}{\( x(x-4)+(x^2-16)(x+2)=(3+x)(x+4) \)}{\( x(x+4)+x+2=(3+x)(x-4) \)}{\( x(x-4)+x+2=(3+x)(x+4) \)}{}{}}
\LANGpl{
\Question{
Wybierz wzór danego równania, który powstał w wyniku  pomnożeniu obu stron przez   \( x^2-16 \).
\[ \frac x{x-4}+x+2=\frac{3+x}{x+4} \]}
{\( x(x+4)+(x^2-16)(x+2)=(3+x)(x-4) \)}{\( x(x-4)+(x^2-16)(x+2)=(3+x)(x+4) \)}{\( x(x+4)+x+2=(3+x)(x-4) \)}{\( x(x-4)+x+2=(3+x)(x+4) \)}{}{}}
\LANGcs{
\Question{
Vyberte tvar rovnice, který dostaneme po vynásobení obou stran dané rovnice výrazem  \( x^2-16 \).
\[ \frac x{x-4}+x+2=\frac{3+x}{x+4} \]}
{\( x(x+4)+(x^2-16)(x+2)=(3+x)(x-4) \)}{\( x(x-4)+(x^2-16)(x+2)=(3+x)(x+4) \)}{\( x(x+4)+x+2=(3+x)(x-4) \)}{\( x(x-4)+x+2=(3+x)(x+4) \)}{}{}}
\LANGsk{
\Question{
Vyberte tvar rovnice, ktorý dostaneme po vynásobení obidvoch strán danej rovnice výrazom \( x^2-16 \).
\[ \frac x{x-4}+x+2=\frac{3+x}{x+4} \]}
{\( x(x+4)+(x^2-16)(x+2)=(3+x)(x-4) \)}{\( x(x-4)+(x^2-16)(x+2)=(3+x)(x+4) \)}{\( x(x+4)+x+2=(3+x)(x-4) \)}{\( x(x-4)+x+2=(3+x)(x+4) \)}{}{}}
\LANGes{
\Question{
Elige la forma final de la ecuación dada después de multiplicar los dos lados por \( x^2-16 \).
\[ \frac x{x-4}+x+2=\frac{3+x}{x+4} \]}
{\( x(x+4)+(x^2-16)(x+2)=(3+x)(x-4) \)}{\( x(x-4)+(x^2-16)(x+2)=(3+x)(x+4) \)}{\( x(x+4)+x+2=(3+x)(x-4) \)}{\( x(x-4)+x+2=(3+x)(x+4) \)}{}{}}
\End

\Data\ID{1003136405}
\Updated{2020-08-22 10:08:55}
\Level{A}
\SubArea{Rational equations and inequalities}
\LANGen{
\Question{
Choose the resulting form of the given equation after multiplying both sides by \( x^2-25 \).
\[ 1+\frac x{5-x}=\frac{3+x}{x+5}+\frac x{x^2-25} \]}
{\( x^2-25-x(x+5)=(3+x)(x-5)+x \)}{\( x^2-25+x(x+5)=(3+x)(x-5)+x \)}{\( x^2-25-x(x-5)=(3+x)(x-5)+x \)}{\( x^2-25+x(x-5)=(3+x)(x+5)+x \)}{}{}}
\LANGpl{
\Question{
Wybierz wzór danego równania, który powstał w wyniku  pomnożeniu obu stron przez \( x^2-25 \).
\[ 1+\frac x{5-x}=\frac{3+x}{x+5}+\frac x{x^2-25} \]}
{\( x^2-25-x(x+5)=(3+x)(x-5)+x \)}{\( x^2-25+x(x+5)=(3+x)(x-5)+x \)}{\( x^2-25-x(x-5)=(3+x)(x-5)+x \)}{\( x^2-25+x(x-5)=(3+x)(x+5)+x \)}{}{}}
\LANGcs{
\Question{
Vyberte tvar rovnice, který dostaneme po vynásobení obou stran dané rovnice výrazem \( x^2-25 \).
\[ 1+\frac x{5-x}=\frac{3+x}{x+5}+\frac x{x^2-25} \]}
{\( x^2-25-x(x+5)=(3+x)(x-5)+x \)}{\( x^2-25+x(x+5)=(3+x)(x-5)+x \)}{\( x^2-25-x(x-5)=(3+x)(x-5)+x \)}{\( x^2-25+x(x-5)=(3+x)(x+5)+x \)}{}{}}
\LANGsk{
\Question{
Vyberte tvar rovnice, ktorý dostaneme po vynásobení obidvoch strán rovnice výrazom \( x^2-25 \).
\[ 1+\frac x{5-x}=\frac{3+x}{x+5}+\frac x{x^2-25} \]}
{\( x^2-25-x(x+5)=(3+x)(x-5)+x \)}{\( x^2-25+x(x+5)=(3+x)(x-5)+x \)}{\( x^2-25-x(x-5)=(3+x)(x-5)+x \)}{\( x^2-25+x(x-5)=(3+x)(x+5)+x \)}{}{}}
\LANGes{
\Question{
Elige la forma final de la ecuación dada después de multiplicar los dos lados por \( x^2-25 \).
\[ 1+\frac x{5-x}=\frac{3+x}{x+5}+\frac x{x^2-25} \]}
{\( x^2-25-x(x+5)=(3+x)(x-5)+x \)}{\( x^2-25+x(x+5)=(3+x)(x-5)+x \)}{\( x^2-25-x(x-5)=(3+x)(x-5)+x \)}{\( x^2-25+x(x-5)=(3+x)(x+5)+x \)}{}{}}
\End

\Data\ID{1003136406}
\Updated{2020-08-22 11:08:11}
\Level{A}
\SubArea{Rational equations and inequalities}
\LANGen{
\Question{
Choose the resulting form of the given equation after multiplying both sides by \( x^2+5x+6 \).
\[ -1+\frac{x-2}{x+2}=\frac{1+x}{x^2+5x+6}-\frac x{x+3} \]}
{\( -x^2-5x-6+(x-2)(x+3)=1+x-x(x+2) \)}{\( -1\left(x^2+5x+6\right)+(x-2)(x-3)=1+x-x(x-2)  \)}{\( -1\left(x^2+5x+6\right)+(x-2)(x+3)=1+x+x(x+2) \)}{\( -x^2-5x-6+(x-2)(x+2)=1+x-x(x+3) \)}{}{}}
\LANGpl{
\Question{
Wybierz wzór danego równania, który powstał w wyniku  pomnożeniu obu stron przez \( x^2+5x+6 \).
\[ -1+\frac{x-2}{x+2}=\frac{1+x}{x^2+5x+6}-\frac x{x+3} \]}
{\( -x^2-5x-6+(x-2)(x+3)=1+x-x(x+2) \)}{\( -1\left(x^2+5x+6\right)+(x-2)(x-3)=1+x-x(x-2)  \)}{\( -1\left(x^2+5x+6\right)+(x-2)(x+3)=1+x+x(x+2) \)}{\( -x^2-5x-6+(x-2)(x+2)=1+x-x(x+3) \)}{}{}}
\LANGcs{
\Question{
Vyberte tvar rovnice, který dostaneme po vynásobení obou stran dané rovnice výrazem \( x^2+5x+6 \).
\[ -1+\frac{x-2}{x+2}=\frac{1+x}{x^2+5x+6}-\frac x{x+3} \]}
{\( -x^2-5x-6+(x-2)(x+3)=1+x-x(x+2) \)}{\( -1\left(x^2+5x+6\right)+(x-2)(x-3)=1+x-x(x-2)  \)}{\( -1\left(x^2+5x+6\right)+(x-2)(x+3)=1+x+x(x+2) \)}{\( -x^2-5x-6+(x-2)(x+2)=1+x-x(x+3) \)}{}{}}
\LANGsk{
\Question{
Vyberte tvar rovnice, ktorý dostaneme po vynásobení obidvoch strán danej rovnice výrazom \( x^2+5x+6 \).
\[ -1+\frac{x-2}{x+2}=\frac{1+x}{x^2+5x+6}-\frac x{x+3} \]}
{\( -x^2-5x-6+(x-2)(x+3)=1+x-x(x+2) \)}{\( -1\left(x^2+5x+6\right)+(x-2)(x-3)=1+x-x(x-2)  \)}{\( -1\left(x^2+5x+6\right)+(x-2)(x+3)=1+x+x(x+2) \)}{\( -x^2-5x-6+(x-2)(x+2)=1+x-x(x+3) \)}{}{}}
\LANGes{
\Question{
Elige la forma final de la ecuación dada después de multiplicar los dos lados por \( x^2+5x+6 \).
\[ -1+\frac{x-2}{x+2}=\frac{1+x}{x^2+5x+6}-\frac x{x+3} \]}
{\( -x^2-5x-6+(x-2)(x+3)=1+x-x(x+2) \)}{\( -1\left(x^2+5x+6\right)+(x-2)(x-3)=1+x-x(x-2)  \)}{\( -1\left(x^2+5x+6\right)+(x-2)(x+3)=1+x+x(x+2) \)}{\( -x^2-5x-6+(x-2)(x+2)=1+x-x(x+3) \)}{}{}}
\End

\Data\ID{1003136407}
\Updated{2020-08-22 11:08:43}
\Level{A}
\SubArea{Rational equations and inequalities}
\LANGen{
\Question{
Choose the resulting form of the given equation after multiplying both sides by \( x^2-9 \).
\[ -\frac{x+1}{9-x^2} =\frac{1+x}{3-x}-\frac x{3+x} \]}
{\( x+1=(1+x)(-x-3)-x(x-3)  \)}{\( -x-1=(1+x)(-x-3)-x(x-3)  \)}{\( x+1=(1+x)(x+3)+x(x-3) \)}{\( -x-1=(1+x)(-x-3)+x(x-3)  \)}{}{}}
\LANGpl{
\Question{
Wybierz wzór danego równania, który powstał w wyniku  pomnożeniu obu stron przez  \( x^2-9 \).
\[ -\frac{x+1}{9-x^2} =\frac{1+x}{3-x}-\frac x{3+x} \]}
{\( x+1=(1+x)(-x-3)-x(x-3)  \)}{\( -x-1=(1+x)(-x-3)-x(x-3)  \)}{\( x+1=(1+x)(x+3)+x(x-3) \)}{\( -x-1=(1+x)(-x-3)+x(x-3)  \)}{}{}}
\LANGcs{
\Question{
Vyberte tvar rovnice, který dostaneme po vynásobení obou stran dané rovnice výrazem \( x^2-9 \).
\[ -\frac{x+1}{9-x^2} =\frac{1+x}{3-x}-\frac x{3+x} \]}
{\( x+1=(1+x)(-x-3)-x(x-3)  \)}{\( -x-1=(1+x)(-x-3)-x(x-3)  \)}{\( x+1=(1+x)(x+3)+x(x-3) \)}{\( -x-1=(1+x)(-x-3)+x(x-3)  \)}{}{}}
\LANGsk{
\Question{
Vyberte tvar rovnice, ktorý dostaneme po vynásobení obidvoch strán danej rovnice výrazom \( x^2-9 \).
\[ -\frac{x+1}{9-x^2} =\frac{1+x}{3-x}-\frac x{3+x} \]}
{\( x+1=(1+x)(-x-3)-x(x-3)  \)}{\( -x-1=(1+x)(-x-3)-x(x-3)  \)}{\( x+1=(1+x)(x+3)+x(x-3) \)}{\( -x-1=(1+x)(-x-3)+x(x-3)  \)}{}{}}
\LANGes{
\Question{
Elige la forma final de la ecuación dada después de multiplicar los dos lados por \( x^2-9 \).
\[ -\frac{x+1}{9-x^2} =\frac{1+x}{3-x}-\frac x{3+x} \]}
{\( x+1=(1+x)(-x-3)-x(x-3)  \)}{\( -x-1=(1+x)(-x-3)-x(x-3)  \)}{\( x+1=(1+x)(x+3)+x(x-3) \)}{\( -x-1=(1+x)(-x-3)+x(x-3)  \)}{}{}}
\End

\Data\ID{1103044801}
\Updated{2020-08-22 11:08:38}
\Level{A}
\SubArea{Rational equations and inequalities}
\IMG{1103044801_0.pdf}
\LANGen{
\Question{
Given graphs of the functions \( f(x) =2x^2-2x-4 \) and \( g(x) = 2x+2 \), find the domain of the equation\( \frac{2x^2-2x-4}{2x+2} = 10 \).}
{\( \mathbb{R}\setminus\{-1\} \)}{\( \mathbb{R}\setminus\{-1;2\} \)}{\( \mathbb{R}\setminus\{-1;2;3\} \)}{\( \mathbb{R}\setminus\{-1;3\} \)}{}{}}
\LANGpl{
\Question{
Dane są wykresy funkcji \( f(x) =2x^2-2x-4 \) i \( g(x) = 2x+2 \), wyznacz dziedzinę równania \( \frac{2x^2-2x-4}{2x+2} = 10 \).}
{\( \mathbb{R}\setminus\{-1\} \)}{\( \mathbb{R}\setminus\{-1;2\} \)}{\( \mathbb{R}\setminus\{-1;2;3\} \)}{\( \mathbb{R}\setminus\{-1;3\} \)}{}{}}
\LANGcs{
\Question{
Pomocí grafů funkcí \( f(x) =2x^2-2x-4 \) a \( g(x) = 2x+2 \) určete množinu, na které má rovnice \( \frac{2x^2-2x-4}{2x+2} = 10 \) smysl.}
{\( \mathbb{R}\setminus\{-1\} \)}{\( \mathbb{R}\setminus\{-1;2\} \)}{\( \mathbb{R}\setminus\{-1;2;3\} \)}{\( \mathbb{R}\setminus\{-1;3\} \)}{}{}}
\LANGsk{
\Question{
Pomocou grafov funkcií \( f(x) =2x^2-2x-4 \) a \( g(x) = 2x+2 \) určte definičný obor rovnice \( \frac{2x^2-2x-4}{2x+2} = 10 \).}
{\( \mathbb{R}\setminus\{-1\} \)}{\( \mathbb{R}\setminus\{-1;2\} \)}{\( \mathbb{R}\setminus\{-1;2;3\} \)}{\( \mathbb{R}\setminus\{-1;3\} \)}{}{}}
\LANGes{
\Question{
Dadas las gráficas de las funciones \( f(x) =2x^2-2x-4 \) y \( g(x) = 2x+2 \), halla el dominio de la ecuación \( \frac{2x^2-2x-4}{2x+2} = 10 \).}
{\( \mathbb{R}\setminus\{-1\} \)}{\( \mathbb{R}\setminus\{-1;2\} \)}{\( \mathbb{R}\setminus\{-1;2;3\} \)}{\( \mathbb{R}\setminus\{-1;3\} \)}{}{}}
\End

\Data\ID{1103044802}
\Updated{2020-08-22 11:08:34}
\Level{A}
\SubArea{Rational equations and inequalities}
\IMG{1103044802_0.pdf}
\LANGen{
\Question{
Given graphs of the functions \( f(x)=x^2-4x \) and \( g(x) = 4x^2-16x+12 \), find the domain of the equation \( \frac{4x^2-16x+12}{x^2-4x}=6 \).}
{\( \mathbb{R}\setminus\{0;4\} \)}{\( \mathbb{R}\setminus\{1;3\} \)}{\( \mathbb{R}\setminus\{0;1;3;4\} \)}{\( \mathbb{R}\setminus\{2\} \)}{}{}}
\LANGpl{
\Question{
Dane są wykresy funkcji \( f(x)=x^2-4x \) i \( g(x) = 4x^2-16x+12 \), wyznacz dziedzinę równania \( \frac{4x^2-16x+12}{x^2-4x}=6 \).}
{\( \mathbb{R}\setminus\{0;4\} \)}{\( \mathbb{R}\setminus\{1;3\} \)}{\( \mathbb{R}\setminus\{0;1;3;4\} \)}{\( \mathbb{R}\setminus\{2\} \)}{}{}}
\LANGcs{
\Question{
Pomocí grafů funkcí \( f(x)=x^2-4x \) a \( g(x) = 4x^2-16x+12 \) určete množinu, na které má rovnice \( \frac{4x^2-16x+12}{x^2-4x}=6 \) smysl.}
{\( \mathbb{R}\setminus\{0;4\} \)}{\( \mathbb{R}\setminus\{1;3\} \)}{\( \mathbb{R}\setminus\{0;1;3;4\} \)}{\( \mathbb{R}\setminus\{2\} \)}{}{}}
\LANGsk{
\Question{
Pomocou grafov funkcií \( f(x)=x^2-4x \) a \( g(x) = 4x^2-16x+12 \) určte definičný obor rovnice \( \frac{4x^2-16x+12}{x^2-4x}=6 \).}
{\( \mathbb{R}\setminus\{0;4\} \)}{\( \mathbb{R}\setminus\{1;3\} \)}{\( \mathbb{R}\setminus\{0;1;3;4\} \)}{\( \mathbb{R}\setminus\{2\} \)}{}{}}
\LANGes{
\Question{
Dadas las gráficas de las funciones \( f(x)=x^2-4x \) y \( g(x) = 4x^2-16x+12 \), halla el dominio de la ecuación \( \frac{4x^2-16x+12}{x^2-4x}=6 \).}
{\( \mathbb{R}\setminus\{0;4\} \)}{\( \mathbb{R}\setminus\{1;3\} \)}{\( \mathbb{R}\setminus\{0;1;3;4\} \)}{\( \mathbb{R}\setminus\{2\} \)}{}{}}
\End

\Data\ID{1103044803}
\Updated{2022-04-23 05:04:48}
\Level{A}
\SubArea{Rational equations and inequalities}
\IMG{1103044803_0.pdf}
\LANGen{
\Question{
Given graphs of the functions  \( f(x)= x^2-x-6 \) and \( g(x) = x+2 \), find the domain of the equation \( \frac{x+2}{x^2-x-6}=1 \).}
{\( \mathbb{R}\setminus\{-2;3\} \)}{\( \mathbb{R}\setminus\{-2;4\} \)}{\( \mathbb{R}\setminus\{0\} \)}{\( \mathbb{R}\setminus\{-2\} \)}{}{}}
\LANGpl{
\Question{
Dane są wykresy funkcji  \( f(x)= x^2-x-6 \) i \( g(x) = x+2 \), wyznacz dziedzinę funkcji \( \frac{x+2}{x^2-x-6}=1 \).}
{\( \mathbb{R}\setminus\{-2;3\} \)}{\( \mathbb{R}\setminus\{-2;4\} \)}{\( \mathbb{R}\setminus\{0\} \)}{\( \mathbb{R}\setminus\{-2\} \)}{}{}}
\LANGcs{
\Question{
Pomocí grafů funkcí \( f(x)= x^2-x-6 \) a \( g(x) = x+2 \) určete množinu, na které má rovnice \( \frac{x+2}{x^2-x-6}=1 \) smysl.}
{\( \mathbb{R}\setminus\{-2;3\} \)}{\( \mathbb{R}\setminus\{-2;4\} \)}{\( \mathbb{R}\setminus\{0\} \)}{\( \mathbb{R}\setminus\{-2\} \)}{}{}}
\LANGsk{
\Question{
Pomocou grafov funkcií  \( f(x)= x^2-x-6 \) a \( g(x) = x+2 \) určte definičný obor funkcie \( \frac{x+2}{x^2-x-6}=1 \).}
{\( \mathbb{R}\setminus\{-2;3\} \)}{\( \mathbb{R}\setminus\{-2;4\} \)}{\( \mathbb{R}\setminus\{0\} \)}{\( \mathbb{R}\setminus\{-2\} \)}{}{}}
\LANGes{
\Question{
Dadas las gráficas de las funciones \( f(x)= x^2-x-6 \) y \( g(x) = x+2 \), halla el dominio de la ecuación \( \frac{x+2}{x^2-x-6}=1 \).}
{\( \mathbb{R}\setminus\{-2;3\} \)}{\( \mathbb{R}\setminus\{-2;4\} \)}{\( \mathbb{R}\setminus\{0\} \)}{\( \mathbb{R}\setminus\{-2\} \)}{}{}}
\End

\Data\ID{1103044804}
\Updated{2020-08-22 12:08:33}
\Level{A}
\SubArea{Rational equations and inequalities}
\IMG{1103044804_0.pdf}
\LANGen{
\Question{
Given graphs of the functions \( f(x) = x^2-x-6 \) and \( g(x) = x+2 \), find the domain of the equation \( \frac{x+2}{x^2-x-6}=\frac{x^2-x-6}{x+2} \).}
{\( \mathbb{R}\setminus\{-2;3\} \)}{\( \mathbb{R}\setminus\{-2;3;4\} \)}{\( \mathbb{R}\setminus\{-2\} \)}{\( \mathbb{R}\setminus\{-2;4\} \)}{}{}}
\LANGpl{
\Question{
Dane są wykresy funkcji \( f(x) = x^2-x-6 \) i \( g(x) = x+2 \), wyznacz dziedzinę równania \( \frac{x+2}{x^2-x-6}=\frac{x^2-x-6}{x+2} \).}
{\( \mathbb{R}\setminus\{-2;3\} \)}{\( \mathbb{R}\setminus\{-2;3;4\} \)}{\( \mathbb{R}\setminus\{-2\} \)}{\( \mathbb{R}\setminus\{-2;4\} \)}{}{}}
\LANGcs{
\Question{
Pomocí grafů funkcí \( f(x) = x^2-x-6 \) a \( g(x) = x+2 \) určete množinu, na které má rovnice \( \frac{x+2}{x^2-x-6}=\frac{x^2-x-6}{x+2} \) smysl.}
{\( \mathbb{R}\setminus\{-2;3\} \)}{\( \mathbb{R}\setminus\{-2;3;4\} \)}{\( \mathbb{R}\setminus\{-2\} \)}{\( \mathbb{R}\setminus\{-2;4\} \)}{}{}}
\LANGsk{
\Question{
Pomocou grafov funkcií \( f(x) = x^2-x-6 \) a \( g(x) = x+2 \) určte definičný obor rovnice \( \frac{x+2}{x^2-x-6}=\frac{x^2-x-6}{x+2} \).}
{\( \mathbb{R}\setminus\{-2;3\} \)}{\( \mathbb{R}\setminus\{-2;3;4\} \)}{\( \mathbb{R}\setminus\{-2\} \)}{\( \mathbb{R}\setminus\{-2;4\} \)}{}{}}
\LANGes{
\Question{
Dadas las gráficas de las funciones \( f(x) = x^2-x-6 \) y \( g(x) = x+2 \), halla el dominio de la ecuación \( \frac{x+2}{x^2-x-6}=\frac{x^2-x-6}{x+2} \).}
{\( \mathbb{R}\setminus\{-2;3\} \)}{\( \mathbb{R}\setminus\{-2;3;4\} \)}{\( \mathbb{R}\setminus\{-2\} \)}{\( \mathbb{R}\setminus\{-2;4\} \)}{}{}}
\End

\Data\ID{1103044805}
\Updated{2020-08-22 12:08:28}
\Level{A}
\SubArea{Rational equations and inequalities}
\IMG{1103044805_0.pdf}
\LANGen{
\Question{
Given graphs of the functions \( f(x)=-x^2-x+6 \) and \( g(x) =x^2-4x+4 \), find the domain of the equation \( \frac{-x^2-x+6}{x^2-4x+4} =-2 \).}
{\( \mathbb{R}\setminus\{2\} \)}{\( \mathbb{R}\setminus\{-3;2\} \)}{\( \mathbb{R}\setminus\{-3;-0.5;2\} \)}{\( \mathbb{R}\setminus\{-2\} \)}{}{}}
\LANGpl{
\Question{
Dane są wykresy funkcji \( f(x)=-x^2-x+6 \) i \( g(x) =x^2-4x+4 \), wyznacz dziedzinę równania \( \frac{-x^2-x+6}{x^2-4x+4} =-2 \).}
{\( \mathbb{R}\setminus\{2\} \)}{\( \mathbb{R}\setminus\{-3;2\} \)}{\( \mathbb{R}\setminus\{-3;-0{,}5;2\} \)}{\( \mathbb{R}\setminus\{-2\} \)}{}{}}
\LANGcs{
\Question{
Pomocí grafů funkcí \( f(x)=-x^2-x+6 \) a \( g(x) =x^2-4x+4 \) určete množinu, na které má rovnice \( \frac{-x^2-x+6}{x^2-4x+4} =-2 \) smysl.}
{\( \mathbb{R}\setminus\{2\} \)}{\( \mathbb{R}\setminus\{-3;2\} \)}{\( \mathbb{R}\setminus\{-3;-0{,}5;2\} \)}{\( \mathbb{R}\setminus\{-2\} \)}{}{}}
\LANGsk{
\Question{
Pomocou grafov funkcií \( f(x)=-x^2-x+6 \) a \( g(x) =x^2-4x+4 \) určte definičný obor rovnice \( \frac{-x^2-x+6}{x^2-4x+4} =-2 \).}
{\( \mathbb{R}\setminus\{2\} \)}{\( \mathbb{R}\setminus\{-3;2\} \)}{\( \mathbb{R}\setminus\{-3;-0{,}5;2\} \)}{\( \mathbb{R}\setminus\{-2\} \)}{}{}}
\LANGes{
\Question{
Dadas las gráficas de las funciones \( f(x)=-x^2-x+6 \) y \( g(x) =x^2-4x+4 \), halla el dominio de la ecuación \( \frac{-x^2-x+6}{x^2-4x+4} =-2 \).}
{\( \mathbb{R}\setminus\{2\} \)}{\( \mathbb{R}\setminus\{-3;2\} \)}{\( \mathbb{R}\setminus\{-3;-0.5;2\} \)}{\( \mathbb{R}\setminus\{-2\} \)}{}{}}
\End

\Data\ID{2000005301}
\Updated{2022-08-25 10:08:50}
\Level{A}
\SubArea{Rational equations and inequalities}
\LANGen{
\Question{
In real numbers, find the solution set of the following equation.
\[\frac{x^2+1}{x^2-9}=0\]}
{\( \emptyset \)}{\( \pm 1\)}{\( \pm 3\)}{\( -1\)}{}{}}
\LANGpl{
\Question{
Wyznacz zbiór rozwiązań równania:
\[\frac{x^2+1}{x^2-9}=0\]}
{\( \emptyset \)}{\( \pm 1\)}{\( \pm 3\)}{\( -1\)}{}{}}
\LANGcs{
\Question{
Určete v \(\mathbb{R}\) množinu řešení dané rovnice.
\[\frac{x^2+1}{x^2-9}=0\]}
{\( \emptyset \)}{\( \pm 1\)}{\( \pm 3\)}{\( -1\)}{}{}}
\LANGsk{
\Question{
Určte v \(\mathbb{R}\) množinu riešení danej rovnice.
\[\frac{x^2+1}{x^2-9}=0\]}
{\( \emptyset \)}{\( \pm 1\)}{\( \pm 3\)}{\( -1\)}{}{}}
\LANGes{
\Question{
En el conjunto de números reales, encuentra el conjunto de soluciones de la ecuación.
\[\frac{x^2+1}{x^2-9}=0\]}
{\( \emptyset \)}{\( \pm 1\)}{\( \pm 3\)}{\( -1\)}{}{}}
\End

\Data\ID{2010012101}
\Updated{2022-11-02 09:11:20}
\Level{A}
\SubArea{Rational equations and inequalities}
\LANGen{
\Question{
Find the solution set of the following equation. 
\[ 
 \frac{2} 
 {9x^2-9} = 0
\]}
{\(\emptyset\)}{\(\{ -1;1\}\)}{\(\{ - 2\}\)}{\(\{ 1\}\)}{}{}}
\LANGpl{
\Question{
Znajdź zbiór rozwiązań następującego równania. 
\[ 
 \frac{2} 
 {9x^2-9} = 0
\]}
{\(\emptyset\)}{\(\{ -1;1\}\)}{\(\{ - 2\}\)}{\(\{ 1\}\)}{}{}}
\LANGcs{
\Question{
Určete množinu řešení dané rovnice. 
\[ 
 \frac{2} 
 {9x^2-9} = 0
\]}
{\(\emptyset\)}{\(\{ -1;1\}\)}{\(\{ - 2\}\)}{\(\{ 1\}\)}{}{}}
\LANGsk{
\Question{
Nájdite množinu riešení danej rovnice. 
\[ 
 \frac{2} 
 {9x^2-9} = 0
\]}
{\(\emptyset\)}{\(\{ -1;1\}\)}{\(\{ - 2\}\)}{\(\{ 1\}\)}{}{}}
\LANGes{
\Question{
Halla el conjunto de soluciones de la siguiente ecuación. 
\[ 
 \frac{2} 
 {9x^2-9} = 0
\]}
{\(\emptyset\)}{\(\{ -1;1\}\)}{\(\{ - 2\}\)}{\(\{ 1\}\)}{}{}}
\End

\Data\ID{2010012102}
\Updated{2022-11-02 09:11:20}
\Level{A}
\SubArea{Rational equations and inequalities}
\LANGen{
\Question{
Find the solution set of the following equation. 
\[ 
 \frac{9x +3} 
{3x + 1} = 3
\]}
{\(\mathbb{R}\setminus \left \{-\frac{1}
{3}\right \}\)}{\(\mathbb{R}\)}{\(\{ 3\}\)}{\(\emptyset\)}{}{}}
\LANGpl{
\Question{
Znajdź zbiór rozwiązań następującego równania. 
\[ 
 \frac{9x +3} 
{3x + 1} = 3
\]}
{\(\mathbb{R}\setminus \left \{-\frac{1}
{3}\right \}\)}{\(\mathbb{R}\)}{\(\{ 3\}\)}{\(\emptyset\)}{}{}}
\LANGcs{
\Question{
Určete množinu řešení dané rovnice. 
\[ 
 \frac{9x +3} 
{3x + 1} = 3
\]}
{\(\mathbb{R}\setminus \left \{-\frac{1}
{3}\right \}\)}{\(\mathbb{R}\)}{\(\{ 3\}\)}{\(\emptyset\)}{}{}}
\LANGsk{
\Question{
Nájdite množinu riešení danej rovnice
\[ 
 \frac{9x +3} 
{3x + 1} = 3
\]}
{\(\mathbb{R}\setminus \left \{-\frac{1}
{3}\right \}\)}{\(\mathbb{R}\)}{\(\{ 3\}\)}{\(\emptyset\)}{}{}}
\LANGes{
\Question{
Halla el conjunto de soluciones de la siguiente ecuación. 
\[ 
 \frac{9x +3} 
{3x + 1} = 3
\]}
{\(\mathbb{R}\setminus \left \{-\frac{1}
{3}\right \}\)}{\(\mathbb{R}\)}{\(\{ 3\}\)}{\(\emptyset\)}{}{}}
\End

\Data\ID{2010012103}
\Updated{2022-11-02 09:11:20}
\Level{A}
\SubArea{Rational equations and inequalities}
\LANGen{
\Question{
Find the domain of the expression.
\[
\frac{x^2-x-12}{3x^2+17x-6}
\]}
{\(\mathbb{R}\setminus \left \{-6;\frac{1}
{3}\right \}\)}{\(\mathbb{R}\setminus \left \{-\frac{1}
{3};6\right \}\)}{\(\left(-\frac13;6\right)\)}{\(\left(-6;\frac13\right)\)}{}{}}
\LANGpl{
\Question{
Określ dziedzinę wyrażenia.
\[
\frac{x^2-x-12}{3x^2+17x-6}
\]}
{\(\mathbb{R}\setminus \left \{-6;\frac{1}
{3}\right \}\)}{\(\mathbb{R}\setminus \left \{-\frac{1}
{3};6\right \}\)}{\(\left(-\frac13;6\right)\)}{\(\left(-6;\frac13\right)\)}{}{}}
\LANGcs{
\Question{
Určete definiční obor daného výrazu.
\[
\frac{x^2-x-12}{3x^2+17x-6}
\]}
{\(\mathbb{R}\setminus \left \{-6;\frac{1}
{3}\right \}\)}{\(\mathbb{R}\setminus \left \{-\frac{1}
{3};6\right \}\)}{\(\left(-\frac13;6\right)\)}{\(\left(-6;\frac13\right)\)}{}{}}
\LANGsk{
\Question{
Určte definičný obor daného výrazu.
\[
\frac{x^2-x-12}{3x^2+17x-6}
\]}
{\(\mathbb{R}\setminus \left \{-6;\frac{1}
{3}\right \}\)}{\(\mathbb{R}\setminus \left \{-\frac{1}
{3};6\right \}\)}{\(\left(-\frac13;6\right)\)}{\(\left(-6;\frac13\right)\)}{}{}}
\LANGes{
\Question{
Halla el dominio de la siguiente expresión.
\[
\frac{x^2-x-12}{3x^2+17x-6}
\]}
{\(\mathbb{R}\setminus \left \{-6;\frac{1}
{3}\right \}\)}{\(\mathbb{R}\setminus \left \{-\frac{1}
{3};6\right \}\)}{\(\left(-\frac13;6\right)\)}{\(\left(-6;\frac13\right)\)}{}{}}
\End

\Data\ID{2010012104}
\Updated{2022-11-02 09:11:20}
\Level{A}
\SubArea{Rational equations and inequalities}
\IMG{2010012101-figure0.pdf}
\LANGen{
\Question{
Given graphs of the functions  \( f(x)= x^2+x-6 \) and \( g(x) = x-2 \), find the domain of the equation \( \frac{x-2}{x^2+x-6}=1 \).}
{\(\mathbb{R}\setminus \left \{-3;2\right \}\)}{\(\mathbb{R}\setminus \left \{-2;2\right \}\)}{\(\mathbb{R}\setminus \left \{-3;-2;2\right \}\)}{\(\mathbb{R}\setminus \left \{0\right \}\)}{}{}}
\LANGpl{
\Question{
Korzystając z wykresów funkcji \( f(x)= x^2+x-6 \) i \( g(x) = x-2 \), znajdź dziedzinę równania\( \frac{x-2}{x^2+x-6}=1 \).}
{\(\mathbb{R}\setminus \left \{-3;2\right \}\)}{\(\mathbb{R}\setminus \left \{-2;2\right \}\)}{\(\mathbb{R}\setminus \left \{-3;-2;2\right \}\)}{\(\mathbb{R}\setminus \left \{0\right \}\)}{}{}}
\LANGcs{
\Question{
Pomocí grafů funkcí  \( f(x)= x^2+x-6 \) a \( g(x) = x-2 \), určete množinu, na které má rovnice \( \frac{x-2}{x^2+x-6}=1 \) smysl.}
{\(\mathbb{R}\setminus \left \{-3;2\right \}\)}{\(\mathbb{R}\setminus \left \{-2;2\right \}\)}{\(\mathbb{R}\setminus \left \{-3;-2;2\right \}\)}{\(\mathbb{R}\setminus \left \{0\right \}\)}{}{}}
\LANGsk{
\Question{
Pomocou grafov funkcií  \( f(x)= x^2+x-6 \) a \( g(x) = x-2 \), určte množinu, na ktorej má rovnica \( \frac{x-2}{x^2+x-6}=1 \) zmysel.}
{\(\mathbb{R}\setminus \left \{-3;2\right \}\)}{\(\mathbb{R}\setminus \left \{-2;2\right \}\)}{\(\mathbb{R}\setminus \left \{-3;-2;2\right \}\)}{\(\mathbb{R}\setminus \left \{0\right \}\)}{}{}}
\LANGes{
\Question{
Dadas las gráficas de las funciones \( f(x)= x^2+x-6 \) y \( g(x) = x-2 \), halla el dominio de la ecuación \( \frac{x-2}{x^2+x-6}=1 \).}
{\(\mathbb{R}\setminus \left \{-3;2\right \}\)}{\(\mathbb{R}\setminus \left \{-2;2\right \}\)}{\(\mathbb{R}\setminus \left \{-3;-2;2\right \}\)}{\(\mathbb{R}\setminus \left \{0\right \}\)}{}{}}
\End

\Data\ID{2010012105}
\Updated{2022-11-02 09:11:20}
\Level{A}
\SubArea{Rational equations and inequalities}
\LANGen{
\Question{
Find the solution set of the following equation.
\[ \frac{x^2-6x+9}{x-3}=0 \]}
{\(\emptyset\)}{\(\left \{3\right \}\)}{\( \left \{-3;3\right \}\)}{\(\left \{-3\right \}\)}{}{}}
\LANGpl{
\Question{
Znajdź zbiór rozwiązań następującego równania.
\[ \frac{x^2-6x+9}{x-3}=0 \]}
{\(\emptyset\)}{\(\left \{3\right \}\)}{\( \left \{-3;3\right \}\)}{\(\left \{-3\right \}\)}{}{}}
\LANGcs{
\Question{
Určete množinu řešení dané rovnice.
\[ \frac{x^2-6x+9}{x-3}=0 \]}
{\(\emptyset\)}{\(\left \{3\right \}\)}{\( \left \{-3;3\right \}\)}{\(\left \{-3\right \}\)}{}{}}
\LANGsk{
\Question{
Nájdite množinu riešení danej rovnice.
\[ \frac{x^2-6x+9}{x-3}=0 \]}
{\(\emptyset\)}{\(\left \{3\right \}\)}{\( \left \{-3;3\right \}\)}{\(\left \{-3\right \}\)}{}{}}
\LANGes{
\Question{
Halla el conjunto de soluciones de la siguiente ecuación.
\[ \frac{x^2-6x+9}{x-3}=0 \]}
{\(\emptyset\)}{\(\left \{3\right \}\)}{\( \left \{-3;3\right \}\)}{\(\left \{-3\right \}\)}{}{}}
\End

\Data\ID{2010012106}
\Updated{2022-11-02 09:11:20}
\Level{A}
\SubArea{Rational equations and inequalities}
\LANGen{
\Question{
Find the solution set of the following equation.
\[ \frac{4x^2-16}{x-2}=0 \]}
{\(\left \{-2\right \}\)}{\( \left \{-2;2\right \}\)}{\(\left \{2\right \}\)}{\(\emptyset\)}{}{}}
\LANGpl{
\Question{
Znajdź zbiór rozwiązań następującego równania.
\[ \frac{4x^2-16}{x-2}=0 \]}
{\(\left \{-2\right \}\)}{\( \left \{-2;2\right \}\)}{\(\left \{2\right \}\)}{\(\emptyset\)}{}{}}
\LANGcs{
\Question{
Určete množinu řešení dané rovnice.
\[ \frac{4x^2-16}{x-2}=0 \]}
{\(\left \{-2\right \}\)}{\( \left \{-2;2\right \}\)}{\(\left \{2\right \}\)}{\(\emptyset\)}{}{}}
\LANGsk{
\Question{
Nájdite množinu riešení danej rovnice.
\[ \frac{4x^2-16}{x-2}=0 \]}
{\(\left \{-2\right \}\)}{\( \left \{-2;2\right \}\)}{\(\left \{2\right \}\)}{\(\emptyset\)}{}{}}
\LANGes{
\Question{
Halla el conjunto de soluciones de la siguiente ecuación.
\[ \frac{4x^2-16}{x-2}=0 \]}
{\(\left \{-2\right \}\)}{\( \left \{-2;2\right \}\)}{\(\left \{2\right \}\)}{\(\emptyset\)}{}{}}
\End

\Data\ID{2010012107}
\Updated{2022-11-02 09:11:20}
\Level{A}
\SubArea{Rational equations and inequalities}
\LANGen{
\Question{
Find the solution set of the following equation.
\[ \frac2{5x^2-20}=0 \]}
{\(\emptyset\)}{\(\left \{2\right \}\)}{\( \left \{-2;2\right \}\)}{\(\left \{-2\right \}\)}{}{}}
\LANGpl{
\Question{
Znajdź zbiór rozwiązań następującego równania.
\[ \frac2{5x^2-20}=0 \]}
{\(\emptyset\)}{\(\left \{2\right \}\)}{\( \left \{-2;2\right \}\)}{\(\left \{-2\right \}\)}{}{}}
\LANGcs{
\Question{
Určete množinu řešení dané rovnice.
\[ \frac2{5x^2-20}=0 \]}
{\(\emptyset\)}{\(\left \{2\right \}\)}{\( \left \{-2;2\right \}\)}{\(\left \{-2\right \}\)}{}{}}
\LANGsk{
\Question{
Nájdite množinu riešení danej rovnice.
\[ \frac2{5x^2-20}=0 \]}
{\(\emptyset\)}{\(\left \{2\right \}\)}{\( \left \{-2;2\right \}\)}{\(\left \{-2\right \}\)}{}{}}
\LANGes{
\Question{
Halla el conjunto de soluciones de la siguiente ecuación.
\[ \frac2{5x^2-20}=0 \]}
{\(\emptyset\)}{\(\left \{2\right \}\)}{\( \left \{-2;2\right \}\)}{\(\left \{-2\right \}\)}{}{}}
\End

\Data\ID{9000024105}
\Updated{2020-08-22 12:08:51}
\Level{A}
\SubArea{Rational equations and inequalities}
\LANGen{
\Question{
Identify the optimal first step to solve the following equation. The operation is
intended to be used on both sides of the equation. 
\[ 
 \frac{4 + x} 
{x + 1} = \frac{x - 3} 
{x + 2}
\]}
{multiply by \((x + 2)\cdot (x + 1)\),
assuming \(x\neq - 2\) 
and \(x\neq - 1\)}{multiply by \((4 + x)\cdot (x - 3)\),
assuming \(x\neq - 4\) 
and \(x\neq 3\)}{multiply by \((4 + x)\cdot (x + 1)\),
assuming \(x\neq - 4\) 
and \(x\neq - 1\)}{multiply by \((x - 3)\cdot (x + 2)\),
assuming \(x\neq 3\) 
and \(x\neq - 2\)}{multiply by \((x - 3)\),
assuming \(x\neq 3\)}{multiply by \((4 + x)\),
assuming \(x\neq - 4\)}}
\LANGpl{
\Question{
Które działanie należy wykonać jako pierwsze, aby rozwiązać podane równanie? Działanie należy wykonać po obydwu stronach równania.
\[ 
 \frac{4 + x} 
{x + 1} = \frac{x - 3} 
{x + 2}
\]}
{pomnożyć przez \((x + 2)\cdot (x + 1)\),
zakładając, że \(x\neq - 2\) 
i \(x\neq - 1\)}{pomnożyć przez \((4 + x)\cdot (x - 3)\),
zakładając, że \(x\neq - 4\) 
i \(x\neq 3\)}{pomnożyć przez \((4 + x)\cdot (x + 1)\),
zakładając, że \(x\neq - 4\) 
i \(x\neq - 1\)}{pomnożyć przez \((x - 3)\cdot (x + 2)\),
zakładając, że \(x\neq 3\) 
i \(x\neq - 2\)}{pomnożyć przez \((x - 3)\),
zakładając, że \(x\neq 3\)}{pomnożyć przez  \((4 + x)\),
zakładając, że \(x\neq - 4\)}}
\LANGcs{
\Question{
Z nabídnutých možností vyberte nejvhodnější ekvivalentní úpravu,
pomocí které začneme řešit danou rovnici. Operace je zamýšlena pro aplikaci na obě strany rovnice.
\[ 
 \frac{4 + x} 
{x + 1} = \frac{x - 3} 
{x + 2}
\]}
{vynásobení výrazem \((x + 2)\cdot (x + 1)\)
za předpokladu \(x\neq - 2\) a \(x\neq - 1\)}{vynásobení výrazem \((4 + x)\cdot (x - 3)\)
za předpokladu \(x\neq - 4\) a \(x\neq 3\)}{vynásobení výrazem \((4 + x)\cdot (x + 1)\)
za předpokladu \(x\neq - 4\) a \(x\neq - 1\)}{vynásobení výrazem \((x - 3)\cdot (x + 2)\)
za předpokladu \(x\neq 3\) a \(x\neq - 2\)}{vynásobení výrazem \((x - 3)\) za předpokladu \(x\neq 3\)}{vynásobení výrazem \((4 + x)\) za předpokladu \(x\neq - 4\)}}
\LANGsk{
\Question{
Z ponúkaných možností vyberte najvhodnejšiu ekvivalentnú úpravu, pomocou ktorej začnete riešiť danú rovnicu. Táto úprava bude aplikovaná na obidve strany rovnice. 
\[ 
 \frac{4 + x} 
{x + 1} = \frac{x - 3} 
{x + 2}
\]}
{vynásobenie výrazom \((x + 2)\cdot (x + 1)\)
za predpokladu \(x\neq - 2\) 
a \(x\neq - 1\)}{vynásobenie výrazom \((4 + x)\cdot (x - 3)\)
za predpokladu \(x\neq - 4\) 
a \(x\neq 3\)}{vynásobenie výrazom \((4 + x)\cdot (x + 1)\)
za predpokladu \(x\neq - 4\) 
a \(x\neq - 1\)}{vynásobenie výrazom \((x - 3)\cdot (x + 2)\)
za predpokladu \(x\neq 3\) 
a \(x\neq - 2\)}{vynásobenie výrazom \((x - 3)\)
za predpokladu \(x\neq 3\)}{vynásobenie výrazom \((4 + x)\)
za predpokladu \(x\neq - 4\)}}
\LANGes{
\Question{
Elige la operación matemática más adecuada para resolver la siguiente ecuación. La operación vale para los dos lados de la ecuación.  
\[ 
 \frac{4 + x} 
{x + 1} = \frac{x - 3} 
{x + 2}
\]}
{multiplicar por \((x + 2)\cdot (x + 1)\),
a condición de que \(x\neq - 2\) 
y \(x\neq - 1\)}{multiplicar por \((4 + x)\cdot (x - 3)\),
a condición de que \(x\neq - 4\) 
y \(x\neq 3\)}{multiplicar por \((4 + x)\cdot (x + 1)\),
a condición de que \(x\neq - 4\) 
y \(x\neq - 1\)}{multiplicar por \((x - 3)\cdot (x + 2)\),
a condición de que \(x\neq 3\) 
y \(x\neq - 2\)}{multiplicar por \((x - 3)\),
a condición de que \(x\neq 3\)}{multiplicar por \((4 + x)\),
a condición de que \(x\neq - 4\)}}
\End

\Data\ID{9000024106}
\Updated{2020-08-22 12:08:35}
\Level{A}
\SubArea{Rational equations and inequalities}
\LANGen{
\Question{
Identify the optimal first step to solve the following equation. The
operation is intended to be used on both sides of the equation. Assume
\(x\neq 1\) and
\(x\neq 2\).
\[ 
 \frac{1} 
{x - 1} = \frac{2} 
{x - 2}
\]}
{multiply by \((x - 1)\cdot (x - 2)\)}{multiply by \((x - 1)\)}{multiply by \((x - 2)\)}{multiply by \((x + 1)\)}{multiply by \((x + 2)\)}{multiply by \((x - 1)\cdot (x + 2)\)}}
\LANGpl{
\Question{
Które działanie należy wykonać jako pierwsze, aby rozwiązać podane równanie? Działanie należy wykonać po obydwu stronach równania. Załóżmy, że 
\(x\neq 1\) i
\(x\neq 2\).
\[ 
 \frac{1} 
{x - 1} = \frac{2} 
{x - 2}
\]}
{pomnożyć przez \((x - 1)\cdot (x - 2)\)}{pomnożyć przez \((x - 1)\)}{pomnożyć przez  \((x - 2)\)}{pomnożyć przez  \((x + 1)\)}{pomnożyć przez  \((x + 2)\)}{pomnożyć przez \((x - 1)\cdot (x + 2)\)}}
\LANGcs{
\Question{
Z nabídnutých možností vyberte nejvhodnější ekvivalentní úpravu,
pomocí které začneme řešit danou rovnici. Operace je zamýšlena pro aplikaci na obě strany rovnice za podmínek
\(x\neq 1\) a \(x\neq 2\).
\[ 
 \frac{1} 
{x - 1} = \frac{2} 
{x - 2}
\]}
{vynásobení výrazem \((x - 1)\cdot (x - 2)\)}{vynásobení výrazem \((x - 1)\)}{vynásobení výrazem \((x - 2)\)}{vynásobení výrazem \((x + 1)\)}{vynásobení výrazem \((x + 2)\)}{vynásobení výrazem \((x - 1)\cdot (x + 2)\)}}
\LANGsk{
\Question{
Z ponúkaných možností vyberte najvhodnejšiu ekvivalentnú úpravu, pomocou ktorej začnete riešiť danú rovnicu. Táto úprava bude aplikovaná na obidve strany rovnice za podmienok
\(x\neq 1\) and
\(x\neq 2\).
\[ 
 \frac{1} 
{x - 1} = \frac{2} 
{x - 2}
\]}
{vynásobenie výrazom \((x - 1)\cdot (x - 2)\)}{vynásobenie výrazom \((x - 1)\)}{vynásobenie výrazom \((x - 2)\)}{vynásobenie výrazom \((x + 1)\)}{vynásobenie výrazom \((x + 2)\)}{vynásobenie výrazom \((x - 1)\cdot (x + 2)\)}}
\LANGes{
\Question{
Elige la operación matemática más adecuada para resolver la siguiente ecuación. La operación vale para los dos lados de la ecuación a condición de que \(x\neq 1\) y
\(x\neq 2\).
\[ 
 \frac{1} 
{x - 1} = \frac{2} 
{x - 2}
\]}
{multiplicar por \((x - 1)\cdot (x - 2)\)}{multiplicar por \((x - 1)\)}{multiplicar por \((x - 2)\)}{multiplicar por \((x + 1)\)}{multiplicar por \((x + 2)\)}{multiplicar por \((x - 1)\cdot (x + 2)\)}}
\End

\Data\ID{9000024109}
\Updated{2020-08-22 12:08:24}
\Level{A}
\SubArea{Rational equations and inequalities}
\LANGen{
\Question{
Identify the optimal first step to solve the following equation. The operation is
intended to be used on both sides of the equation. 
\[ 
 \frac{2x + 1} 
 {x - 1} + \frac{x + 1} 
{x - 1} = \frac{11} 
 {2}
\]}
{multiply by \(2(x - 1)\),
assuming \(x\neq 1\)}{multiply by \((2x + 1)\),
assuming \(x\neq -\frac{1} 
{2}\)}{multiply by \((x + 1)\),
assuming \(x\neq - 1\)}{multiply by \(\frac{1}
{2x+1}\),
assuming \(x\neq -\frac{1} 
{2}\)}{multiply by \(\frac{1}
{x+1}\),
assuming \(x\neq - 1\)}{multiply by \(2(2x + 1)(x + 1)\),
assuming \(x\neq -\frac{1} 
{2}\) 
and \(x\neq - 1\)}}
\LANGpl{
\Question{
Które działanie należy wykonać jako pierwsze, aby rozwiązać podane równanie? Działanie należy wykonać po obydwu stronach równania.
\[ 
 \frac{2x + 1} 
 {x - 1} + \frac{x + 1} 
{x - 1} = \frac{11} 
 {2}
\]}
{pomnożyć przez \(2(x - 1)\),
zakładając, że \(x\neq 1\)}{pomnożyć przez  \((2x + 1)\),
zakładając, że \(x\neq -\frac{1} 
{2}\)}{pomnożyć przez \((x + 1)\),
zakładając, że \(x\neq - 1\)}{pomnożyć przez  \(\frac{1}
{2x+1}\),
zakładając, że \(x\neq -\frac{1} 
{2}\)}{pomnożyć przez \(\frac{1}
{x+1}\),
zakładając, że \(x\neq - 1\)}{pomnożyć przez  \(2(2x + 1)(x + 1)\),
zakładając, że \(x\neq -\frac{1} 
{2}\) 
i \(x\neq - 1\)}}
\LANGcs{
\Question{
Z nabídnutých možností vyberte nejvhodnější ekvivalentní úpravu,
pomocí které začneme řešit danou rovnici. Operace je zamýšlena pro aplikaci na obě strany rovnice.
\[ 
 \frac{2x + 1} 
 {x - 1} + \frac{x + 1} 
{x - 1} = \frac{11} 
 {2}
\]}
{vynásobení výrazem \(2(x - 1)\) za předpokladu \(x\neq 1\)}{vynásobení výrazem \((2x + 1)\) za předpokladu \(x\neq -\frac{1} {2}\)}{vynásobení výrazem \((x + 1)\)
za předpokladu \(x\neq - 1\)}{vynásobení výrazem \(\frac{1} {2x+1}\) za předpokladu 
\(x\neq -\frac{1} {2}\)}{vynásobení výrazem \(\frac{1}{x+1}\) za předpokladu \(x\neq - 1\)}{vynásobení výrazem \(2(2x + 1)(x + 1)\) za předpokladu 
\(x\neq -\frac{1} {2}\) a \(x\neq - 1\)}}
\LANGsk{
\Question{
Z ponúkaných možností vyberte najvhodnejšiu ekvivalentnú úpravu, pomocou ktorej začnete riešiť danú rovnicu. Táto úprava bude aplikovaná na obidve strany rovnice. 
\[ 
 \frac{2x + 1} 
 {x - 1} + \frac{x + 1} 
{x - 1} = \frac{11} 
 {2}
\]}
{vynásobenie výrazom \(2(x - 1)\)
za predpokladu \(x\neq 1\)}{vynásobenie výrazom \((2x + 1)\)
za predpokladu \(x\neq -\frac{1} 
{2}\)}{vynásobenie výrazom \((x + 1)\)
za predpokladu \(x\neq - 1\)}{vynásobenie výrazom \(\frac{1}
{2x+1}\)
za predpokladu \(x\neq -\frac{1} 
{2}\)}{vynásobenie výrazom \(\frac{1}
{x+1}\)
za predpokladu \(x\neq - 1\)}{vynásobenie výrazom \(2(2x + 1)(x + 1)\)
za predpokladu \(x\neq -\frac{1} 
{2}\) 
a \(x\neq - 1\)}}
\LANGes{
\Question{
Elige la operación matemática más adecuada para resolver la siguiente ecuación. La operación vale para los dos lados de la ecuación.  
\[ 
 \frac{2x + 1} 
 {x - 1} + \frac{x + 1} 
{x - 1} = \frac{11} 
 {2}
\]}
{multiplicar por \(2(x - 1)\),
a condición de que \(x\neq 1\)}{multiplicar por \((2x + 1)\),
a condición de que \(x\neq -\frac{1} 
{2}\)}{multiplicar por \((x + 1)\),
a condición de que \(x\neq - 1\)}{multiplicar por \(\frac{1}
{2x+1}\),
a condición de que \(x\neq -\frac{1} 
{2}\)}{multiplicar por \(\frac{1}
{x+1}\), a condición de que \(x\neq - 1\)}{multiplicar por \(2(2x + 1)(x + 1)\),
a condición de que \(x\neq -\frac{1} 
{2}\) 
y \(x\neq - 1\)}}
\End

\Data\ID{9000033301}
\Updated{2022-04-23 05:04:48}
\Level{A}
\SubArea{Rational equations and inequalities}
\LANGen{
\Question{
Find the solution set of the following equation. 
\[ 
 \frac{3x + 6} 
 {2 - x} = 0
\]}
{\(\{ - 2\}\)}{\(\emptyset \)}{\(\{2\}\)}{\(\mathbb{R}\setminus \{2\}\)}{}{}}
\LANGpl{
\Question{
Wyznacz zbiór rozwiązań podanego równania 
\[ 
 \frac{3x + 6} 
 {2 - x} = 0
\]}
{\(\{ - 2\}\)}{\(\emptyset \)}{\(\{2\}\)}{\(\mathbb{R}\setminus \{2\}\)}{}{}}
\LANGcs{
\Question{
Určete množinu řešení dané rovnice.
\[ \frac{3x+6}{2-x} = 0\]}
{\(\{ - 2\}\)}{\(\emptyset \)}{\(\{2\}\)}{\(\mathbb{R}\setminus \{2\}\)}{}{}}
\LANGsk{
\Question{
Určte množinu riešení danej rovnice.
\[ 
 \frac{3x + 6} 
 {2 - x} = 0
\]}
{\(\{ - 2\}\)}{\(\emptyset \)}{\(\{2\}\)}{\(\mathbb{R}\setminus \{2\}\)}{}{}}
\LANGes{
\Question{
Halla el conjunto de soluciones de la siguiente ecuación. 
\[ 
 \frac{3x + 6} 
 {2 - x} = 0
\]}
{\(\{ - 2\}\)}{\(\emptyset \)}{\(\{2\}\)}{\(\mathbb{R}\setminus \{2\}\)}{}{}}
\End

\Data\ID{9000033302}
\Updated{2022-04-23 05:04:48}
\Level{A}
\SubArea{Rational equations and inequalities}
\LANGen{
\Question{
Find the solution set of the following equation. 
\[ 
 \frac{4x - 2} 
{2x - 1} = 2
\]}
{\(\mathbb{R}\setminus \left \{\frac{1}
{2}\right \}\)}{\(\mathbb{R}\)}{\(\{2\}\)}{\(\emptyset \)}{}{}}
\LANGpl{
\Question{
Wyznacz zbiór rozwiązań podanego równania 
\[ 
 \frac{4x - 2} 
{2x - 1} = 2
\]}
{\(\mathbb{R}\setminus \left \{\frac{1}
{2}\right \}\)}{\(\mathbb{R}\)}{\(\{2\}\)}{\(\emptyset \)}{}{}}
\LANGcs{
\Question{
Určete množinu řešení dané rovnice. \[\frac{4x-2}{2x-1} = 2\]}
{\(\mathbb{R}\setminus \left \{\frac{1}
{2}\right \}\)}{\(\mathbb{R}\)}{\(\{2\}\)}{\(\emptyset \)}{}{}}
\LANGsk{
\Question{
Určte množinu riešení danej rovnice. 
\[ 
 \frac{4x - 2} 
{2x - 1} = 2
\]}
{\(\mathbb{R}\setminus \left \{\frac{1}
{2}\right \}\)}{\(\mathbb{R}\)}{\(\{2\}\)}{\(\emptyset \)}{}{}}
\LANGes{
\Question{
Halla el conjunto de soluciones de la siguiente ecuación.  
\[ 
 \frac{4x - 2} 
{2x - 1} = 2
\]}
{\(\mathbb{R}\setminus \left \{\frac{1}
{2}\right \}\)}{\(\mathbb{R}\)}{\(\{2\}\)}{\(\emptyset \)}{}{}}
\End

\Data\ID{9000033303}
\Updated{2020-08-22 12:08:14}
\Level{A}
\SubArea{Rational equations and inequalities}
\LANGen{
\Question{
Find the solution set of the following equation. 
\[ 
 \frac{4x + 8} 
 {x + 2} = 0
\]}
{\(\emptyset \)}{\(\{- 2\}\)}{\(\{2\}\)}{\(\left \{-\frac{3}
{4}\right \}\)}{}{}}
\LANGpl{
\Question{
Wyznacz zbiór rozwiązań podanego równania 
\[ 
 \frac{4x + 8} 
 {x + 2} = 0
\]}
{\(\emptyset \)}{\(\{- 2\}\)}{\(\{2\}\)}{\(\left \{-\frac{3}
{4}\right \}\)}{}{}}
\LANGcs{
\Question{
Určete množinu řešení dané rovnice. \[\frac{4x+8} {x+2} = 0\]}
{\(\emptyset \)}{\(\{- 2\}\)}{\(\{2\}\)}{\(\left \{-\frac{3}
{4}\right \}\)}{}{}}
\LANGsk{
\Question{
Určte množinu riešení danej rovnice. 
\[ 
 \frac{4x + 8} 
 {x + 2} = 0
\]}
{\(\emptyset \)}{\(\{- 2\}\)}{\(\{2\}\)}{\(\left \{-\frac{3}
{4}\right \}\)}{}{}}
\LANGes{
\Question{
Halla el conjunto de soluciones de la siguiente ecuación.  
\[ 
 \frac{4x + 8} 
 {x + 2} = 0
\]}
{\(\emptyset \)}{\(\{- 2\}\)}{\(\{2\}\)}{\(\left \{-\frac{3}
{4}\right \}\)}{}{}}
\End

\Data\ID{9000079203}
\Updated{2020-08-22 12:08:05}
\Level{A}
\SubArea{Rational equations and inequalities}
\LANGen{
\Question{
Find all real \(x\) 
for which the following expression equals zero. 
\[ 
 1 -\frac{2x + 1} 
 {x - 1}
\]}
{\(x = -2\)}{\(x = -\frac{1}
{2}\)}{\(x = 0\)}{\(x = -1\)}{}{}}
\LANGpl{
\Question{
Znajdź wszystkie rzeczywiste wartości \(x\), dla których podane wyrażenie jest równe \( 0\).  
\[ 
 1 -\frac{2x + 1} 
 {x - 1}
\]}
{\(x = -2\)}{\(x = -\frac{1}
{2}\)}{\(x = 0\)}{\(x = -1\)}{}{}}
\LANGcs{
\Question{
Pro kterou hodnotu proměnné \(x\) 
je výraz \(1 -\frac{2x+1} 
 {x-1} \) 
roven nule?}
{\(x = -2\)}{\(x = -\frac{1}
{2}\)}{\(x = 0\)}{\(x = -1\)}{}{}}
\LANGsk{
\Question{
Pre ktorú hodnotu premennej \(x\) 
je výraz \(1 -\frac{2x+1} 
 {x-1} \) 
rovný nule?}
{\(x = -2\)}{\(x = -\frac{1}
{2}\)}{\(x = 0\)}{\(x = -1\)}{}{}}
\LANGes{
\Question{
Encuentra todos los valores de \(x\) para los que la siguiente expresión es igual a cero. 
\[ 
 1 -\frac{2x + 1} 
 {x - 1}
\]}
{\(x = -2\)}{\(x = -\frac{1}
{2}\)}{\(x = 0\)}{\(x = -1\)}{}{}}
\End

\Data\ID{9000083701}
\Updated{2020-11-09 07:11:09}
\Level{A}
\SubArea{Rational equations and inequalities}
\LANGen{
\Question{
Find all the values of \(x\in \mathbb{R}\) 
for which the given expression equals zero. 
\[ 
 \frac{x^{2} - 16} 
 {2x - 8}
\]}
{\(x = -4\)}{\(x = 4\)}{\(x =\pm 4\)}{\(x = 0\)}{}{}}
\LANGpl{
\Question{
Znajdź takie \(x\in \mathbb{R}\), 
dla którego podane wyrażenie jest równe zeru.
\[ 
 \frac{x^{2} - 16} 
 {2x - 8}
\]}
{\(x = -4\)}{\(x = 4\)}{\(x =\pm 4\)}{\(x = 0\)}{}{}}
\LANGcs{
\Question{
Uveďte všechny hodnoty \(x\in \mathbb{R}\),
pro které je výraz \(\frac{x^{2}-16}
 {2x-8} \) 
roven \(0\).}
{\(x = -4\)}{\(x = 4\)}{\(x =\pm 4\)}{\(x = 0\)}{}{}}
\LANGsk{
\Question{
Určte všetky hodnoty \(x\in \mathbb{R}\),
pre ktoré je výraz \(\frac{x^{2}-16}
 {2x-8} \) 
rovný \(0\).}
{\(x = -4\)}{\(x = 4\)}{\(x =\pm 4\)}{\(x = 0\)}{}{}}
\LANGes{
\Question{
Encuentra todos los valores de \(x\in \mathbb{R}\) 
para los que la fracción dada es igual a cero.  
\[ 
 \frac{x^{2} - 16} 
 {2x - 8}
\]}
{\(x = -4\)}{\(x = 4\)}{\(x =\pm 4\)}{\(x = 0\)}{}{}}
\End

\Data\ID{9000083702}
\Updated{2020-11-09 07:11:39}
\Level{A}
\SubArea{Rational equations and inequalities}
\LANGen{
\Question{
Find all the values of \(x\in \mathbb{R}\) 
for which the given expression equals zero. 
\[ 
 \frac{x^{2} + 6x + 9} 
 {x^{2} - 9}
\]}
{The expression never equals zero.}{\(x =\pm 3\)}{\(x = 3\)}{\(x = -3\)}{}{}}
\LANGpl{
\Question{
Znajdź takie  \(x\in \mathbb{R}\),
dla którego podane wyrażenie jest równe zeru.
\[ 
 \frac{x^{2} + 6x + 9} 
 {x^{2} - 9}
\]}
{Wyrażenie nigdy nie jest równe zeru.}{\(x =\pm 3\)}{\(x = 3\)}{\(x = -3\)}{}{}}
\LANGcs{
\Question{
Uveďte všechny hodnoty \(x\in \mathbb{R}\),
pro které je výraz \(\frac{x^{2}+6x+9}
 {x^{2}-9} \) 
roven \(0\).}
{Uvedený výraz nenabývá hodnoty
\(0\) pro
žádné reálné číslo.}{\(x =\pm 3\)}{\(x = 3\)}{\(x = -3\)}{}{}}
\LANGsk{
\Question{
Určte všetky hodnoty \(x\in \mathbb{R}\),
pre ktoré je výraz \(\frac{x^{2}+6x+9}
 {x^{2}-9} \) 
rovný \(0\).}
{Daný výraz nenadobúda hodnotu
\(0\) pre
žiadne reálne číslo.}{\(x =\pm 3\)}{\(x = 3\)}{\(x = -3\)}{}{}}
\LANGes{
\Question{
Encuentra todos los valores de \(x\in \mathbb{R}\) 
para los que la fracción dada es igual a cero.  
\[ 
 \frac{x^{2} + 6x + 9} 
 {x^{2} - 9}
\]}
{Esta fracción nunca es igual a cero.}{\(x =\pm 3\)}{\(x = 3\)}{\(x = -3\)}{}{}}
\End

\Data\ID{9000083703}
\Updated{2020-11-09 07:11:23}
\Level{A}
\SubArea{Rational equations and inequalities}
\LANGen{
\Question{
Find all the values of \(x\in \mathbb{R}\) 
for which the given expression equals zero. 
\[ 
 \frac{x^{3} - x} 
 {x - 1}
\]}
{\(x = -1,\ x = 0\)}{\(x = 0\)}{\(x = 1\)}{\(x = -1,\ x = 0,\ x = 1\)}{}{}}
\LANGpl{
\Question{
Znajdź takie  \(x\in \mathbb{R}\), 
dla którego podane wyrażenie jest równe zeru.
\[ 
 \frac{x^{3} - x} 
 {x - 1}
\]}
{\(x = -1,\ x = 0\)}{\(x = 0\)}{\(x = 1\)}{\(x = -1,\ x = 0,\ x = 1\)}{}{}}
\LANGcs{
\Question{
Uveďte všechny hodnoty \(x\in \mathbb{R}\),
pro které je výraz \(\frac{x^{3}-x}
 {x-1} \) 
roven \(0\).}
{\(x = -1,\ x = 0\)}{\(x = 0\)}{\(x = 1\)}{\(x = -1,\ x = 0,\ x = 1\)}{}{}}
\LANGsk{
\Question{
Určte všetky hodnoty \(x\in \mathbb{R}\),
pre ktoré je výraz \(\frac{x^{3}-x}
 {x-1} \) 
rovný \(0\).}
{\(x = -1,\ x = 0\)}{\(x = 0\)}{\(x = 1\)}{\(x = -1,\ x = 0,\ x = 1\)}{}{}}
\LANGes{
\Question{
Halla todos los valores de \(x\in \mathbb{R}\) 
para los que la fracción dada es igual a cero. 
\[ 
 \frac{x^{3} - x} 
 {x - 1}
\]}
{\(x = -1,\ x = 0\)}{\(x = 0\)}{\(x = 1\)}{\(x = -1,\ x = 0,\ x = 1\)}{}{}}
\End

\Data\ID{9000083704}
\Updated{2020-11-09 07:11:48}
\Level{A}
\SubArea{Rational equations and inequalities}
\LANGen{
\Question{
Find all the values of \(x\in \mathbb{R}\) 
for which the given expression equals zero. 
\[ 
 \frac{x^{2} - 4x + 4} 
 {x(x - 2)}
\]}
{The expression never equals zero.}{\(x = 0\)}{\(x = 2\)}{\(x = -2,\ x = 0\)}{}{}}
\LANGpl{
\Question{
Znajdź takie  \(x\in \mathbb{R}\),
dla którego podane wyrażenie jest równe zeru.
\[ 
 \frac{x^{2} - 4x + 4} 
 {x(x - 2)}
\]}
{Wyrażenie nigdy nie jest równe zeru.}{\(x = 0\)}{\(x = 2\)}{\(x = -2,\ x = 0\)}{}{}}
\LANGcs{
\Question{
Uveďte všechny hodnoty \(x\in \mathbb{R}\),
pro které je výraz \(\frac{x^{2}-4x+4}
 {x(x-2)} \) 
roven \(0\).}
{Uvedený výraz nenabývá hodnoty
\(0\) pro
žádné reálné číslo.}{\(x = 0\)}{\(x = 2\)}{\(x = -2,\ x = 0\)}{}{}}
\LANGsk{
\Question{
Určte všetky hodnoty \(x\in \mathbb{R}\),
pre ktoré je výraz \(\frac{x^{2}-4x+4}
 {x(x-2)} \) 
rovný \(0\).}
{Daný výraz nenadobúda hodnotu \(0\) pre žiadne reálne číslo.}{\(x = 0\)}{\(x = 2\)}{\(x = -2,\ x = 0\)}{}{}}
\LANGes{
\Question{
Halla todos los valores de \(x\in \mathbb{R}\) 
para los que la fracción dada es igual a cero. 
\[ 
 \frac{x^{2} - 4x + 4} 
 {x(x - 2)}
\]}
{Esta fracción nunca es igual a cero.}{\(x = 0\)}{\(x = 2\)}{\(x = -2,\ x = 0\)}{}{}}
\End

\Data\ID{9000083705}
\Updated{2020-11-09 07:11:07}
\Level{A}
\SubArea{Rational equations and inequalities}
\LANGen{
\Question{
Find all the values of \(x\in \mathbb{R}\) 
for which the given expression equals zero. 
\[ 
 \frac{2x(x + 2)(x - 3)} 
 {x^{2} - 4}
\]}
{\(x = 0,\ x = 3\)}{\(x = -2,\ x = 0,\ x = 3\)}{\(x = 0\)}{\(x =\pm 2\)}{}{}}
\LANGpl{
\Question{
Znajdź takie \(x\in \mathbb{R}\), 
dla którego podane wyrażenie jest równe zeru.
\[ 
 \frac{2x(x + 2)(x - 3)} 
 {x^{2} - 4}
\]}
{\(x = 0,\ x = 3\)}{\(x = -2,\ x = 0,\ x = 3\)}{\(x = 0\)}{\(x =\pm 2\)}{}{}}
\LANGcs{
\Question{
Uveďte všechny hodnoty \(x\in \mathbb{R}\),
pro které je výraz \(\frac{2x(x+2)(x-3)}
 {x^{2}-4} \) 
roven \(0\).}
{\(x = 0,\ x = 3\)}{\(x = -2,\ x = 0,\ x = 3\)}{\(x = 0\)}{\(x =\pm 2\)}{}{}}
\LANGsk{
\Question{
Určte všetky hodnoty \(x\in \mathbb{R}\),
pre ktoré je výraz \(\frac{2x(x+2)(x-3)}
 {x^{2}-4} \) 
rovný \(0\).}
{\(x = 0,\ x = 3\)}{\(x = -2,\ x = 0,\ x = 3\)}{\(x = 0\)}{\(x =\pm 2\)}{}{}}
\LANGes{
\Question{
Halla todos los valores de \(x\in \mathbb{R}\) 
para los que la fracción dada es igual a cero. 
\[ 
 \frac{2x(x + 2)(x - 3)} 
 {x^{2} - 4}
\]}
{\(x = 0,\ x = 3\)}{\(x = -2,\ x = 0,\ x = 3\)}{\(x = 0\)}{\(x =\pm 2\)}{}{}}
\End

\Data\ID{9000083706}
\Updated{2020-11-09 07:11:22}
\Level{A}
\SubArea{Rational equations and inequalities}
\LANGen{
\Question{
Find all the values of \(x\in \mathbb{R}\) 
for which the given expression equals zero. 
\[ 
 \frac{4x^{2} - 36} 
{4x^{2} + 24x + 36}
\]}
{\(x = 3\)}{\(x = 4\)}{\(x = -3,\ x = 3\)}{The expression never equals zero.}{}{}}
\LANGpl{
\Question{
Znajdź takie \(x\in \mathbb{R}\), dla którego podane wyrażenie jest równe zeru.
\[ 
 \frac{4x^{2} - 36} 
{4x^{2} + 24x + 36}
\]}
{\(x = 3\)}{\(x = 4\)}{\(x = -3,\ x = 3\)}{Wyrażenie nigdy nie jest równe zeru.}{}{}}
\LANGcs{
\Question{
Uveďte všechny hodnoty \(x\in \mathbb{R}\),
pro které je výraz \(\frac{4x^{2}-36}
{4x^{2}+24x+36}\) 
roven \(0\).}
{\(x = 3\)}{\(x = 4\)}{\(x = -3,\ x = 3\)}{Uvedený výraz nenabývá hodnoty
\(0\) pro
žádné reálné číslo.}{}{}}
\LANGsk{
\Question{
Určte všetky hodnoty \(x\in \mathbb{R}\),
pre ktoré je výraz \(\frac{4x^{2}-36}
{4x^{2}+24x+36}\) 
rovný \(0\).}
{\(x = 3\)}{\(x = 4\)}{\(x = -3,\ x = 3\)}{Daný výraz nenadobúda hodnotu \(0\) pre žiadne reálne číslo.}{}{}}
\LANGes{
\Question{
Halla todos los valores de \(x\in \mathbb{R}\) 
para los que la fracción dada es igual a cero. 
\[ 
 \frac{4x^{2} - 36} 
{4x^{2} + 24x + 36}
\]}
{\(x = 3\)}{\(x = 4\)}{\(x = -3,\ x = 3\)}{Esta fracción nunca es igual a cero.}{}{}}
\End

\Data\ID{9000083707}
\Updated{2020-11-09 07:11:38}
\Level{A}
\SubArea{Rational equations and inequalities}
\LANGen{
\Question{
Find all the values of \(x\in \mathbb{R}\) 
for which the given expression equals zero. 
\[ 
 \frac{4x^{3} + 20x^{2} + 25x} 
 {x + 1}
\]}
{\(x = 0,\ x = -\frac{5}
{2}\)}{\(x = 0\)}{\(x = -\frac{5}
{2}\)}{\(x = -1\)}{}{}}
\LANGpl{
\Question{
Znajdź takie \(x\in \mathbb{R}\), 
dla którego podane wyrażenie jest równe zeru.
\[ 
 \frac{4x^{3} + 20x^{2} + 25x} 
 {x + 1}
\]}
{\(x = 0,\ x = -\frac{5}
{2}\)}{\(x = 0\)}{\(x = -\frac{5}
{2}\)}{\(x = -1\)}{}{}}
\LANGcs{
\Question{
Uveďte všechny hodnoty \(x\in \mathbb{R}\),
pro které je výraz \(\frac{4x^{3}+20x^{2}+25x}
 {x+1} \) 
roven \(0\).}
{\(x = 0,\ x = -\frac{5}
{2}\)}{\(x = 0\)}{\(x = -\frac{5}
{2}\)}{\(x = -1\)}{}{}}
\LANGsk{
\Question{
Určte všetky hodnoty \(x\in \mathbb{R}\),
pre ktoré je výraz \(\frac{4x^{3}+20x^{2}+25x}
 {x+1} \) 
rovný \(0\).}
{\(x = 0,\ x = -\frac{5}
{2}\)}{\(x = 0\)}{\(x = -\frac{5}
{2}\)}{\(x = -1\)}{}{}}
\LANGes{
\Question{
Halla todos los valores de \(x\in \mathbb{R}\) 
para los que la fracción dada es igual a cero. 
\[ 
 \frac{4x^{3} + 20x^{2} + 25x} 
 {x + 1}
\]}
{\(x = 0,\ x = -\frac{5}
{2}\)}{\(x = 0\)}{\(x = -\frac{5}
{2}\)}{\(x = -1\)}{}{}}
\End

\Data\ID{9000083708}
\Updated{2020-11-09 07:11:02}
\Level{A}
\SubArea{Rational equations and inequalities}
\LANGen{
\Question{
Find all the values of \(x\in \mathbb{R}\) 
for which the given expression equals zero. 
\[ 
 \frac{x^{2} - (2x - 1)^{2}} 
 {x^{2} - 4}
\]}
{\(x = \frac{1} 
{3},\ x = 1\)}{\(x = -\frac{1}
{3},\ x = 1\)}{\(x =\pm 2\)}{\(x = 1\)}{}{}}
\LANGpl{
\Question{
Znajdź takie  \(x\in \mathbb{R}\), 
dla którego podane wyrażenie jest równe zeru.
\[ 
 \frac{x^{2} - (2x - 1)^{2}} 
 {x^{2} - 4}
\]}
{\(x = \frac{1} 
{3},\ x = 1\)}{\(x = -\frac{1}
{3},\ x = 1\)}{\(x =\pm 2\)}{\(x = 1\)}{}{}}
\LANGcs{
\Question{
Uveďte všechny hodnoty \(x\in \mathbb{R}\),
pro které je výraz \(\frac{x^{2}-(2x-1)^{2}} 
 {x^{2}-4} \) 
roven \(0\).}
{\(x = \frac{1} 
{3},\ x = 1\)}{\(x = -\frac{1}
{3},\ x = 1\)}{\(x =\pm 2\)}{\(x = 1\)}{}{}}
\LANGsk{
\Question{
Určte všetky hodnoty \(x\in \mathbb{R}\),
pre ktoré je výraz \(\frac{x^{2}-(2x-1)^{2}} 
 {x^{2}-4} \) 
rovný \(0\).}
{\(x = \frac{1} 
{3},\ x = 1\)}{\(x = -\frac{1}
{3},\ x = 1\)}{\(x =\pm 2\)}{\(x = 1\)}{}{}}
\LANGes{
\Question{
Halla todos los valores de \(x\in \mathbb{R}\) 
para los que la fracción dada es igual a cero.
\[ 
 \frac{x^{2} - (2x - 1)^{2}} 
 {x^{2} - 4}
\]}
{\(x = \frac{1} 
{3},\ x = 1\)}{\(x = -\frac{1}
{3},\ x = 1\)}{\(x =\pm 2\)}{\(x = 1\)}{}{}}
\End

\Data\ID{9000083709}
\Updated{2020-11-09 07:11:32}
\Level{A}
\SubArea{Rational equations and inequalities}
\LANGen{
\Question{
Find all the values of \(x\in \mathbb{R}\) 
for which the given expression equals zero. 
\[ 
 \frac{(2x + 3)^{2} - (3x - 2)^{2}} 
 {x - 5}
\]}
{\(x = -\frac{1}
{5}\)}{\(x = 5\)}{\(x = -5\)}{\(x = \frac{1} 
{5}\)}{}{}}
\LANGpl{
\Question{
Znajdź takie  \(x\in \mathbb{R}\), 
dla którego podane wyrażenie jest równe zeru.
\[ 
 \frac{(2x + 3)^{2} - (3x - 2)^{2}} 
 {x - 5}
\]}
{\(x = -\frac{1}
{5}\)}{\(x = 5\)}{\(x = -5\)}{\(x = \frac{1} 
{5}\)}{}{}}
\LANGcs{
\Question{
Uveďte všechny hodnoty \(x\in \mathbb{R}\),
pro které je výraz \(\frac{(2x+3)^{2}-(3x-2)^{2}} 
 {x-5} \) 
roven \(0\).}
{\(x = -\frac{1}
{5}\)}{\(x = 5\)}{\(x = -5\)}{\(x = \frac{1} 
{5}\)}{}{}}
\LANGsk{
\Question{
Určte všetky hodnoty \(x\in \mathbb{R}\),
pre ktoré je výraz \(\frac{(2x+3)^{2}-(3x-2)^{2}} 
 {x-5} \) 
rovný \(0\).}
{\(x = -\frac{1}
{5}\)}{\(x = 5\)}{\(x = -5\)}{\(x = \frac{1} 
{5}\)}{}{}}
\LANGes{
\Question{
Halla todos los valores de \(x\in \mathbb{R}\) 
para los que la fracción dada es igual a cero.
\[ 
 \frac{(2x + 3)^{2} - (3x - 2)^{2}} 
 {x - 5}
\]}
{\(x = -\frac{1}
{5}\)}{\(x = 5\)}{\(x = -5\)}{\(x = \frac{1} 
{5}\)}{}{}}
\End

\Data\ID{9000083710}
\Updated{2020-11-09 07:11:57}
\Level{A}
\SubArea{Rational equations and inequalities}
\LANGen{
\Question{
Find all the values of \(x\in \mathbb{R}\) 
for which the given expression equals zero. 
\[ 
 \frac{(4x + 3)^{2} - (5x - 2)^{2}} 
 {5 + x}
\]}
{\(x = 5,\ x = -\frac{1}
{9}\)}{\(x = -5\)}{\(x = -\frac{5}
{9},\ x = 1\)}{\(x = 1,\ x = \frac{5} 
{9}\)}{}{}}
\LANGpl{
\Question{
Znajdź takie \(x\in \mathbb{R}\), 
dla którego podane wyrażenie jest równe zeru.
\[ 
 \frac{(4x + 3)^{2} - (5x - 2)^{2}} 
 {5 + x}
\]}
{\(x = 5,\ x = -\frac{1}
{9}\)}{\(x = -5\)}{\(x = -\frac{5}
{9},\ x = 1\)}{\(x = 1,\ x = \frac{5} 
{9}\)}{}{}}
\LANGcs{
\Question{
Uveďte všechny hodnoty \(x\in \mathbb{R}\),
pro které je výraz \(\frac{(4x+3)^{2}-(5x-2)^{2}} 
 {5+x} \) 
roven \(0\).}
{\(x = 5,\ x = -\frac{1}
{9}\)}{\(x = -5\)}{\(x = -\frac{5}
{9},\ x = 1\)}{\(x = 1,\ x = \frac{5} 
{9}\)}{}{}}
\LANGsk{
\Question{
Určte všetky hodnoty \(x\in \mathbb{R}\),
pre ktoré je výraz \(\frac{(4x+3)^{2}-(5x-2)^{2}} 
 {5+x} \) 
rovný \(0\).}
{\(x = 5,\ x = -\frac{1}
{9}\)}{\(x = -5\)}{\(x = -\frac{5}
{9},\ x = 1\)}{\(x = 1,\ x = \frac{5} 
{9}\)}{}{}}
\LANGes{
\Question{
Halla todos los valores de \(x\in \mathbb{R}\) 
para los que la fracción dada es igual a cero.
\[ 
 \frac{(4x + 3)^{2} - (5x - 2)^{2}} 
 {5 + x}
\]}
{\(x = 5,\ x = -\frac{1}
{9}\)}{\(x = -5\)}{\(x = -\frac{5}
{9},\ x = 1\)}{\(x = 1,\ x = \frac{5} 
{9}\)}{}{}}
\End

\Data\ID{1003029001}
\Updated{2022-08-25 10:08:11}
\Level{B}
\SubArea{Rational equations and inequalities}
\LANGen{
\Question{
Find the solution set of the inequality.
\[ \left(x^2+1\right)\left(x^2+3\right)\geq0 \]}
{\( \mathbb{R} \)}{\( (-\infty;-1]\cup[1;\infty) \)}{\( (-\infty;-1)\cup(1;\infty) \)}{\( (-\infty;-\sqrt3]\cup[\sqrt3;\infty) \)}{\( \emptyset \)}{}}
\LANGpl{
\Question{
Wyznacz zbiór rozwiązań podanej nierówności.
\[ \left(x^2+1\right)\left(x^2+3\right)\geq0 \]}
{\( \mathbb{R} \)}{\( (-\infty;-1\rangle\cup\langle1;\infty) \)}{\( (-\infty;-1)\cup(1;\infty) \)}{\( (-\infty;-\sqrt3\rangle\cup\langle\sqrt3;\infty) \)}{\( \emptyset \)}{}}
\LANGcs{
\Question{
Určete množinu řešení nerovnice.
\[ \left(x^2+1\right)\left(x^2+3\right)\geq0 \]}
{\( \mathbb{R} \)}{\( (-\infty;-1\rangle\cup\langle1;\infty) \)}{\( (-\infty;-1)\cup(1;\infty) \)}{\( (-\infty;-\sqrt3\rangle\cup\langle\sqrt3;\infty) \)}{\( \emptyset \)}{}}
\LANGsk{
\Question{
Určte množinu riešení nerovnice.
\[ \left(x^2+1\right)\left(x^2+3\right)\geq0 \]}
{\( \mathbb{R} \)}{\( (-\infty;-1\rangle\cup\langle1;\infty) \)}{\( (-\infty;-1)\cup(1;\infty) \)}{\( (-\infty;-\sqrt3\rangle\cup\langle\sqrt3;\infty) \)}{\( \emptyset \)}{}}
\LANGes{
\Question{
Encuentra el conjunto de soluciones de la inecuación. 
\[ \left(x^2+1\right)\left(x^2+3\right)\geq0 \]}
{\( \mathbb{R} \)}{\( (-\infty;-1]\cup[1;\infty) \)}{\( (-\infty;-1)\cup(1;\infty) \)}{\( (-\infty;-\sqrt3]\cup[\sqrt3;\infty) \)}{\( \emptyset \)}{}}
\End

\Data\ID{1003029101}
\Updated{2022-08-25 10:08:11}
\Level{B}
\SubArea{Rational equations and inequalities}
\LANGen{
\Question{
Find the solution set of the inequality.
\[ \frac{x^4-1}{x\left(x^2+3\right)}\geq0 \]}
{\( [-1;0)\cup[1;\infty) \)}{\( [-1;0]\cup[1;\infty) \)}{\( (-\infty;0)\cup[1;\infty) \)}{\( (-\infty;0]\cup[1;\infty) \)}{}{}}
\LANGpl{
\Question{
Wyznacz zbiór rozwiązań podanej nierówności.
\[ \frac{x^4-1}{x\left(x^2+3\right)}\geq0 \]}
{\( \langle-1;0)\cup\langle1;\infty) \)}{\( \langle-1;0\rangle\cup\langle1;\infty) \)}{\( (-\infty;0)\cup\langle1;\infty) \)}{\( (-\infty;0\rangle\cup\langle1;\infty) \)}{}{}}
\LANGcs{
\Question{
Určete množinu řešení nerovnice.
\[ \frac{x^4-1}{x\left(x^2+3\right)}\geq0 \]}
{\( \langle-1;0)\cup\langle1;\infty) \)}{\( \langle-1;0\rangle\cup\langle1;\infty) \)}{\( (-\infty;0)\cup\langle1;\infty) \)}{\( (-\infty;0\rangle\cup\langle1;\infty) \)}{}{}}
\LANGsk{
\Question{
Určte množinu riešení nerovnice.
\[ \frac{x^4-1}{x\left(x^2+3\right)}\geq0 \]}
{\( \langle-1;0)\cup\langle1;\infty) \)}{\( \langle-1;0\rangle\cup\langle1;\infty) \)}{\( (-\infty;0)\cup\langle1;\infty) \)}{\( (-\infty;0\rangle\cup\langle1;\infty) \)}{}{}}
\LANGes{
\Question{
Encuentra el conjunto de soluciones de la inecuación. 
\[ \frac{x^4-1}{x\left(x^2+3\right)}\geq0 \]}
{\( [-1;0)\cup[1;\infty) \)}{\( [-1;0]\cup[1;\infty) \)}{\( (-\infty;0)\cup[1;\infty) \)}{\( (-\infty;0]\cup[1;\infty) \)}{}{}}
\End

\Data\ID{1003029102}
\Updated{2022-08-25 10:08:14}
\Level{B}
\SubArea{Rational equations and inequalities}
\LANGen{
\Question{
Find the solution set of the inequality.
\[ \frac{x\left(x^3-8\right)}{-x^4-4}\leq0 \]}
{\( (-\infty;0]\cup[2;\infty) \)}{\( [ 0;2 ] \)}{\( (-\infty;-2]\cup[0;2] \)}{\( (-\infty;-2]\cup[2;\infty) \)}{}{}}
\LANGpl{
\Question{
Wyznacz zbiór rozwiązań podanej nierówności.
\[ \frac{x\left(x^3-8\right)}{-x^4-4}\leq0 \]}
{\( (-\infty;0\rangle\cup\langle2;\infty) \)}{\( \langle 0;2 \rangle \)}{\( (-\infty;-2\rangle\cup\langle0;2\rangle \)}{\( (-\infty;-2\rangle\cup\langle2;\infty) \)}{}{}}
\LANGcs{
\Question{
Určete množinu řešení nerovnice.
\[ \frac{x\left(x^3-8\right)}{-x^4-4}\leq0 \]}
{\( (-\infty;0\rangle\cup\langle2;\infty) \)}{\( \langle 0;2 \rangle \)}{\( (-\infty;-2\rangle\cup\langle0;2\rangle \)}{\( (-\infty;-2\rangle\cup\langle2;\infty) \)}{}{}}
\LANGsk{
\Question{
Určte množinu riešení nerovnice.
\[ \frac{x\left(x^3-8\right)}{-x^4-4}\leq0 \]}
{\( (-\infty;0\rangle\cup\langle2;\infty) \)}{\( \langle 0;2 \rangle \)}{\( (-\infty;-2\rangle\cup\langle0;2\rangle \)}{\( (-\infty;-2\rangle\cup\langle2;\infty) \)}{}{}}
\LANGes{
\Question{
Encuentra el conjunto de soluciones de la inecuación. 
\[ \frac{x\left(x^3-8\right)}{-x^4-4}\leq0 \]}
{\( (-\infty;0]\cup[2;\infty) \)}{\( [ 0;2 ] \)}{\( (-\infty;-2]\cup[0;2] \)}{\( (-\infty;-2]\cup[2;\infty) \)}{}{}}
\End

\Data\ID{1003029103}
\Updated{2022-08-25 10:08:14}
\Level{B}
\SubArea{Rational equations and inequalities}
\LANGen{
\Question{
Find the domain of the expression on the left side of the inequality.
\[\frac{x^4}{x^2\left(x^5-1\right)\left(2x^2-4\right)}\leq0 \]}
{\( \mathbb{R}\setminus\left\{0;1;\pm\sqrt2\right\}\)}{\( \mathbb{R}\setminus\left\{1;\pm\sqrt2\right\}\)}{\( \mathbb{R}\setminus\left\{\pm1;\pm\sqrt2\right\}\)}{\( \mathbb{R}\setminus\left\{0;\pm1;\pm\sqrt2\right\}\)}{}{}}
\LANGpl{
\Question{
Wyznacz dziedzinę wyrażenia po lewej stronie podanej nierówności.
\[\frac{x^4}{x^2\left(x^5-1\right)\left(2x^2-4\right)}\leq0 \]}
{\( \mathbb{R}\setminus\left\{0;1;\pm\sqrt2\right\}\)}{\( \mathbb{R}\setminus\left\{1;\pm\sqrt2\right\}\)}{\( \mathbb{R}\setminus\left\{\pm1;\pm\sqrt2\right\}\)}{\( \mathbb{R}\setminus\left\{0;\pm1;\pm\sqrt2\right\}\)}{}{}}
\LANGcs{
\Question{
Určete definiční obor nerovnice.
\[\frac{x^4}{x^2\left(x^5-1\right)\left(2x^2-4\right)}\leq0 \]}
{\( \mathbb{R}\setminus\left\{0;1;\pm\sqrt2\right\}\)}{\( \mathbb{R}\setminus\left\{1;\pm\sqrt2\right\}\)}{\( \mathbb{R}\setminus\left\{\pm1;\pm\sqrt2\right\}\)}{\( \mathbb{R}\setminus\left\{0;\pm1;\pm\sqrt2\right\}\)}{}{}}
\LANGsk{
\Question{
Určte definičný obor nerovnice.
\[\frac{x^4}{x^2\left(x^5-1\right)\left(2x^2-4\right)}\leq0 \]}
{\( \mathbb{R}\setminus\left\{0;1;\pm\sqrt2\right\}\)}{\( \mathbb{R}\setminus\left\{1;\pm\sqrt2\right\}\)}{\( \mathbb{R}\setminus\left\{\pm1;\pm\sqrt2\right\}\)}{\( \mathbb{R}\setminus\left\{0;\pm1;\pm\sqrt2\right\}\)}{}{}}
\LANGes{
\Question{
Halla el dominio de la inecuación. \[\frac{x^4}{x^2\left(x^5-1\right)\left(2x^2-4\right)}\leq0 \]}
{\( \mathbb{R}\setminus\left\{0;1;\pm\sqrt2\right\}\)}{\( \mathbb{R}\setminus\left\{1;\pm\sqrt2\right\}\)}{\( \mathbb{R}\setminus\left\{\pm1;\pm\sqrt2\right\}\)}{\( \mathbb{R}\setminus\left\{0;\pm1;\pm\sqrt2\right\}\)}{}{}}
\End

\Data\ID{1003029104}
\Updated{2022-08-25 10:08:14}
\Level{B}
\SubArea{Rational equations and inequalities}
\LANGen{
\Question{
Find the domain of the expression on the left side of the inequality.
\[ \frac{x^3-x^2+1}{\left(x^2+9\right)\left(x^3-1\right)}>0 \]}
{\( \mathbb{R}\setminus\left\{1\right\} \)}{\( \mathbb{R} \)}{\( \mathbb{R}\setminus\left\{\pm1\right\} \)}{\( \mathbb{R}\setminus\left\{\pm3;\pm1\right\} \)}{}{}}
\LANGpl{
\Question{
Wyznacz dziedzinę wyrażenia po lewej stronie podanej nierówności.
\[ \frac{x^3-x^2+1}{\left(x^2+9\right)\left(x^3-1\right)}>0 \]}
{\( \mathbb{R}\setminus\left\{1\right\} \)}{\( \mathbb{R} \)}{\( \mathbb{R}\setminus\left\{\pm1\right\} \)}{\( \mathbb{R}\setminus\left\{\pm3;\pm1\right\} \)}{}{}}
\LANGcs{
\Question{
Určete definiční obor nerovnice.
\[ \frac{x^3-x^2+1}{\left(x^2+9\right)\left(x^3-1\right)}>0 \]}
{\( \mathbb{R}\setminus\left\{1\right\} \)}{\( \mathbb{R} \)}{\( \mathbb{R}\setminus\left\{\pm1\right\} \)}{\( \mathbb{R}\setminus\left\{\pm3;\pm1\right\} \)}{}{}}
\LANGsk{
\Question{
Určte definičný obor nerovnice.
\[ \frac{x^3-x^2+1}{\left(x^2+9\right)\left(x^3-1\right)}>0 \]}
{\( \mathbb{R}\setminus\left\{1\right\} \)}{\( \mathbb{R} \)}{\( \mathbb{R}\setminus\left\{\pm1\right\} \)}{\( \mathbb{R}\setminus\left\{\pm3;\pm1\right\} \)}{}{}}
\LANGes{
\Question{
Halla el dominio de la inecuación.
\[ \frac{x^3-x^2+1}{\left(x^2+9\right)\left(x^3-1\right)}>0 \]}
{\( \mathbb{R}\setminus\left\{1\right\} \)}{\( \mathbb{R} \)}{\( \mathbb{R}\setminus\left\{\pm1\right\} \)}{\( \mathbb{R}\setminus\left\{\pm3;\pm1\right\} \)}{}{}}
\End

\Data\ID{1003138301}
\Updated{2020-08-22 11:08:51}
\Level{B}
\SubArea{Rational equations and inequalities}
\LANGen{
\Question{
Choose the resulting form of the given inequality after multiplying both sides by \( x^2-16 \), where \( x\in(4;\infty) \).
\[ \frac1{x^2-16}-\frac x{4-x} < \frac{3+x}{x+4} \]}
{\( 1+x(x+4) < (3+x)(x-4) \)}{\( 1-x(x+4) < (3+x)(x-4) \)}{\( 1+x(x+4) > (3+x)(x-4) \)}{\( 1-x(x-4) > (3+x)(x+4) \)}{}{}}
\LANGpl{
\Question{
Wybierz wzór danej nierówności, który powstał w wyniku  pomnożeniu obu stron przez  \( x^2-16 \), gdzie \( x\in(4;\infty) \).
\[ \frac1{x^2-16}-\frac x{4-x} < \frac{3+x}{x+4} \]}
{\( 1+x(x+4) < (3+x)(x-4) \)}{\( 1-x(x+4) < (3+x)(x-4) \)}{\( 1+x(x+4) > (3+x)(x-4) \)}{\( 1-x(x-4) > (3+x)(x+4) \)}{}{}}
\LANGcs{
\Question{
Vyberte tvar nerovnice, který dostaneme po vynásobení obou stran dané nerovnice výrazem  \( x^2-16 \), kde \( x\in(4;\infty) \).
\[ \frac1{x^2-16}-\frac x{4-x} < \frac{3+x}{x+4} \]}
{\( 1+x(x+4) < (3+x)(x-4) \)}{\( 1-x(x+4) < (3+x)(x-4) \)}{\( 1+x(x+4) > (3+x)(x-4) \)}{\( 1-x(x-4) > (3+x)(x+4) \)}{}{}}
\LANGsk{
\Question{
Vyberte tvar nerovnice, ktorý dostaneme po vynásobení obidvoch strán danej nerovnice výrazom \( x^2-16 \), kde \( x\in(4;\infty) \).
\[ \frac1{x^2-16}-\frac x{4-x} < \frac{3+x}{x+4} \]}
{\( 1+x(x+4) < (3+x)(x-4) \)}{\( 1-x(x+4) < (3+x)(x-4) \)}{\( 1+x(x+4) > (3+x)(x-4) \)}{\( 1-x(x-4) > (3+x)(x+4) \)}{}{}}
\LANGes{
\Question{
Elige la forma final de la inecuación dada después de multiplicar los dos lados por \( x^2-16 \), donde \( x\in(4;\infty) \).
\[ \frac1{x^2-16}-\frac x{4-x} < \frac{3+x}{x+4} \]}
{\( 1+x(x+4) < (3+x)(x-4) \)}{\( 1-x(x+4) < (3+x)(x-4) \)}{\( 1+x(x+4) > (3+x)(x-4) \)}{\( 1-x(x-4) > (3+x)(x+4) \)}{}{}}
\End

\Data\ID{1003138302}
\Updated{2020-08-22 11:08:32}
\Level{B}
\SubArea{Rational equations and inequalities}
\LANGen{
\Question{
Choose the resulting form of the given inequality after multiplying both sides by \( x^2-25 \), where \( x\in(-1;1) \).
\[ \frac{3+x}{x+5}-\frac{x+1}{x-5} < \frac x{x^2-25} \]}
{\( (3+x)(x-5)-(x+1)(x+5) > x \)}{\( (3+x)(x-5)-(x+1)(x+5) < x \)}{\( (3+x)(x-5)+(x+1)(x+5) > x \)}{\( (3+x)(x+5)-(x+1)(x-5) > x \)}{}{}}
\LANGpl{
\Question{
Wybierz wzór danej nierówności, który powstał w wyniku  pomnożeniu obu stron przez  \( x^2-25 \), gdzie \( x\in(-1;1) \).
\[ \frac{3+x}{x+5}-\frac{x+1}{x-5} < \frac x{x^2-25} \]}
{\( (3+x)(x-5)-(x+1)(x+5) > x \)}{\( (3+x)(x-5)-(x+1)(x+5) < x \)}{\( (3+x)(x-5)+(x+1)(x+5) > x \)}{\( (3+x)(x+5)-(x+1)(x-5) > x \)}{}{}}
\LANGcs{
\Question{
Vyberte tvar nerovnice, který dostaneme po vynásobení obou stran dané nerovnice výrazem  \( x^2-25 \), kde \( x\in(-1;1) \).
\[ \frac{3+x}{x+5}-\frac{x+1}{x-5} < \frac x{x^2-25} \]}
{\( (3+x)(x-5)-(x+1)(x+5) > x \)}{\( (3+x)(x-5)-(x+1)(x+5) < x \)}{\( (3+x)(x-5)+(x+1)(x+5) > x \)}{\( (3+x)(x+5)-(x+1)(x-5) > x \)}{}{}}
\LANGsk{
\Question{
Vyberte tvar nerovnice, ktorý dostaneme po vynásobení obidvoch strán danej nerovnice výrazom \( x^2-25 \), kde \( x\in(-1;1) \).
\[ \frac{3+x}{x+5}-\frac{x+1}{x-5} < \frac x{x^2-25} \]}
{\( (3+x)(x-5)-(x+1)(x+5) > x \)}{\( (3+x)(x-5)-(x+1)(x+5) < x \)}{\( (3+x)(x-5)+(x+1)(x+5) > x \)}{\( (3+x)(x+5)-(x+1)(x-5) > x \)}{}{}}
\LANGes{
\Question{
Elige la forma final de la inecuación dada después de multiplicar los dos lados por \( x^2-25 \), donde \( x\in(-1;1) \).
\[ \frac{3+x}{x+5}-\frac{x+1}{x-5} < \frac x{x^2-25} \]}
{\( (3+x)(x-5)-(x+1)(x+5) > x \)}{\( (3+x)(x-5)-(x+1)(x+5) < x \)}{\( (3+x)(x-5)+(x+1)(x+5) > x \)}{\( (3+x)(x+5)-(x+1)(x-5) > x \)}{}{}}
\End

\Data\ID{1003138303}
\Updated{2020-08-22 11:08:05}
\Level{B}
\SubArea{Rational equations and inequalities}
\LANGen{
\Question{
Choose the resulting form of the given inequality after multiplying both sides by \( 4x^2 \), where \( x\neq0 \). 
\[ \frac2{x^2}-\frac x{2x} \geq \frac{2-x}4 \]}
{\( 8-2x^2 \geq (2-x)x^2 \)}{\( 4-2x \geq (2-x)x^2 \)}{\( 8-2x \leq (2-x) x^2 \)}{\( 4-2x^2 \geq (2-x)  x^2 \)}{}{}}
\LANGpl{
\Question{
Wybierz wzór danej nierówności, który powstał w wyniku  pomnożeniu obu stron przez  \( 4x^2 \), gdzie \( x\neq0 \). 
\[ \frac2{x^2}-\frac x{2x} \geq \frac{2-x}4 \]}
{\( 8-2x^2 \geq (2-x)x^2 \)}{\( 4-2x \geq (2-x)x^2 \)}{\( 8-2x \leq (2-x) x^2 \)}{\( 4-2x^2 \geq (2-x)  x^2 \)}{}{}}
\LANGcs{
\Question{
Vyberte tvar nerovnice, který dostaneme po vynásobení obou stran dané nerovnice výrazem  \( 4x^2 \), kde \( x\neq0 \). 
\[ \frac2{x^2}-\frac x{2x} \geq \frac{2-x}4 \]}
{\( 8-2x^2 \geq (2-x)x^2 \)}{\( 4-2x \geq (2-x)x^2 \)}{\( 8-2x \leq (2-x) x^2 \)}{\( 4-2x^2 \geq (2-x)  x^2 \)}{}{}}
\LANGsk{
\Question{
Vyberte tvar nerovnice, ktorý dostaneme po vynásobení obidvoch strán danej nerovnice výrazom \( 4x^2 \), kde \( x\neq0 \). 
\[ \frac2{x^2}-\frac x{2x} \geq \frac{2-x}4 \]}
{\( 8-2x^2 \geq (2-x)x^2 \)}{\( 4-2x \geq (2-x)x^2 \)}{\( 8-2x \leq (2-x) x^2 \)}{\( 4-2x^2 \geq (2-x)  x^2 \)}{}{}}
\LANGes{
\Question{
Elige la forma final de la inecuación dada después de multiplicar los dos lados por \( 4x^2 \), donde \( x\neq0 \). 
\[ \frac2{x^2}-\frac x{2x} \geq \frac{2-x}4 \]}
{\( 8-2x^2 \geq (2-x)x^2 \)}{\( 4-2x \geq (2-x)x^2 \)}{\( 8-2x \leq (2-x) x^2 \)}{\( 4-2x^2 \geq (2-x)  x^2 \)}{}{}}
\End

\Data\ID{1003138304}
\Updated{2020-08-22 11:08:44}
\Level{B}
\SubArea{Rational equations and inequalities}
\LANGen{
\Question{
Choose the resulting form of the given inequality after multiplying both sides by \( 4x-12 \), where \( x\in(-\infty;0) \).
\[ \frac{x+1}{x-3}-\frac x4 < 0 \]}
{\( 4x+4-x(x-3) > 0 \)}{\( 4(x+1)-x(x-3) < 0 \)}{\( 4x+1-x(x-3) > 0 \)}{\( 4x+4-x(x-3) > 4x-12 \)}{}{}}
\LANGpl{
\Question{
Wybierz wzór danej nierówności, który powstał w wyniku  pomnożeniu obu stron przez  \( 4x-12 \), gdzie \( x\in(-\infty;0) \).
\[ \frac{x+1}{x-3}-\frac x4 < 0 \]}
{\( 4x+4-x(x-3) > 0 \)}{\( 4(x+1)-x(x-3) < 0 \)}{\( 4x+1-x(x-3) > 0 \)}{\( 4x+4-x(x-3) > 4x-12 \)}{}{}}
\LANGcs{
\Question{
Vyberte tvar nerovnice, který dostaneme po vynásobení obou stran dané nerovnice výrazem  \( 4x-12 \), kde \( x\in(-\infty;0) \).
\[ \frac{x+1}{x-3}-\frac x4 < 0 \]}
{\( 4x+4-x(x-3) > 0 \)}{\( 4(x+1)-x(x-3) < 0 \)}{\( 4x+1-x(x-3) > 0 \)}{\( 4x+4-x(x-3) > 4x-12 \)}{}{}}
\LANGsk{
\Question{
Vyberte tvar nerovnice, ktorý dostaneme po vynásobení obidvoch strán danej nerovnice výrazom \( 4x-12 \), kde \( x\in(-\infty;0) \).
\[ \frac{x+1}{x-3}-\frac x4 < 0 \]}
{\( 4x+4-x(x-3) > 0 \)}{\( 4(x+1)-x(x-3) < 0 \)}{\( 4x+1-x(x-3) > 0 \)}{\( 4x+4-x(x-3) > 4x-12 \)}{}{}}
\LANGes{
\Question{
Elige la forma final de la inecuación dada después de multiplicar los dos lados por \( 4x-12 \), donde \( x\in(-\infty;0) \).
\[ \frac{x+1}{x-3}-\frac x4 < 0 \]}
{\( 4x+4-x(x-3) > 0 \)}{\( 4(x+1)-x(x-3) < 0 \)}{\( 4x+1-x(x-3) > 0 \)}{\( 4x+4-x(x-3) > 4x-12 \)}{}{}}
\End

\Data\ID{1003138305}
\Updated{2020-08-22 11:08:23}
\Level{B}
\SubArea{Rational equations and inequalities}
\LANGen{
\Question{
Choose the resulting form of the given inequality after multiplying both sides by \( (x-1)(x-2) \), where \( x\in(0;1) \).
\[ 1 \leq \frac{x-3}{1-x}+\frac{x-1}{x-2} \]}
{\( (x-1)(x-2) \leq (3-x)(x-2)+(x-1)(x-1) \)}{\( (x-1)(x-2) \geq (x-3)(2-x)+(x-1)(x-1) \)}{\( (x-1)(x-2) \leq (x-3)(x-2)+(x-1)(x-1) \)}{\( (x-1)(x-2) \leq -x-3(x-2)+(x-1)^2 \)}{}{}}
\LANGpl{
\Question{
Wybierz wzór danej nierówności, który powstał w wyniku  pomnożeniu obu stron przez \( (x-1)(x-2) \), gdzie \( x\in(0;1) \).
\[ 1 \leq \frac{x-3}{1-x}+\frac{x-1}{x-2} \]}
{\( (x-1)(x-2) \leq (3-x)(x-2)+(x-1)(x-1) \)}{\( (x-1)(x-2) \geq (x-3)(2-x)+(x-1)(x-1) \)}{\( (x-1)(x-2) \leq (x-3)(x-2)+(x-1)(x-1) \)}{\( (x-1)(x-2) \leq -x-3(x-2)+(x-1)^2 \)}{}{}}
\LANGcs{
\Question{
Vyberte tvar nerovnice, který dostaneme po vynásobení obou stran dané nerovnice výrazem  \( (x-1)(x-2) \), kde \( x\in(0;1) \).
\[ 1 \leq \frac{x-3}{1-x}+\frac{x-1}{x-2} \]}
{\( (x-1)(x-2) \leq (3-x)(x-2)+(x-1)(x-1) \)}{\( (x-1)(x-2) \geq (x-3)(2-x)+(x-1)(x-1) \)}{\( (x-1)(x-2) \leq (x-3)(x-2)+(x-1)(x-1) \)}{\( (x-1)(x-2) \leq -x-3(x-2)+(x-1)^2 \)}{}{}}
\LANGsk{
\Question{
Vyberte tvar nerovnice, ktorý dostaneme po vynásobení obidvoch strán danej nerovnice výrazom \( (x-1)(x-2) \), kde \( x\in(0;1) \).
\[ 1 \leq \frac{x-3}{1-x}+\frac{x-1}{x-2} \]}
{\( (x-1)(x-2) \leq (3-x)(x-2)+(x-1)(x-1) \)}{\( (x-1)(x-2) \geq (x-3)(2-x)+(x-1)(x-1) \)}{\( (x-1)(x-2) \leq (x-3)(x-2)+(x-1)(x-1) \)}{\( (x-1)(x-2) \leq -x-3(x-2)+(x-1)^2 \)}{}{}}
\LANGes{
\Question{
Elige la forma final de la inecuación dada después de multiplicar los dos lados por \( (x-1)(x-2) \), donde \( x\in(0;1) \).
\[ 1 \leq \frac{x-3}{1-x}+\frac{x-1}{x-2} \]}
{\( (x-1)(x-2) \leq (3-x)(x-2)+(x-1)(x-1) \)}{\( (x-1)(x-2) \geq (x-3)(2-x)+(x-1)(x-1) \)}{\( (x-1)(x-2) \leq (x-3)(x-2)+(x-1)(x-1) \)}{\( (x-1)(x-2) \leq -x-3(x-2)+(x-1)^2 \)}{}{}}
\End

\Data\ID{2000005304}
\Updated{2022-02-15 07:02:33}
\Level{B}
\SubArea{Rational equations and inequalities}
\LANGen{
\Question{
Find all real values of \( x\) such  that the fraction \( \frac{5}{x^2}\) is positive.}
{\( x \in \mathbb{R}\setminus \{0\}\)}{\( x \in \mathbb{R}\)}{\( x \in (0;+\infty)\)}{\( x \in [ 0;+\infty)\)}{}{}}
\LANGpl{
\Question{
Znajdź wszystkie wartości rzeczywiste \( x\) takie, że ułamek \( \frac{5}{x^2}\) jest dodatni.}
{\( x \in \mathbb{R}\setminus \{0\}\)}{\( x \in \mathbb{R}\)}{\( x \in (0;+\infty)\)}{\( x \in \langle 0;+\infty)\)}{}{}}
\LANGcs{
\Question{
Najděte všechny reálné hodnoty \( x\)  tak, aby byl zlomek \( \frac{5}{x^2}\)  kladný.}
{\( x \in \mathbb{R}\setminus \{0\}\)}{\( x \in \mathbb{R}\)}{\( x \in (0;+\infty)\)}{\( x \in \langle 0;+\infty)\)}{}{}}
\LANGsk{
\Question{
Nájdite všetky reálne hodnoty \(x \) tak, aby bol zlomok \( \frac{5}{x^2}\) kladný.}
{\( x \in \mathbb{R}\setminus \{0\}\)}{\( x \in \mathbb{R}\)}{\( x \in (0;+\infty)\)}{\( x \in \langle 0;+\infty)\)}{}{}}
\LANGes{
\Question{
Halla todos los valores reales de \( x\) para los que la fracción \( \frac{5}{x^2}\) es positiva.}
{\( x \in \mathbb{R}\setminus \{0\}\)}{\( x \in \mathbb{R}\)}{\( x \in (0;+\infty)\)}{\( x \in [ 0;+\infty)\)}{}{}}
\End

\Data\ID{2000005305}
\Updated{2022-02-15 07:02:33}
\Level{B}
\SubArea{Rational equations and inequalities}
\LANGen{
\Question{
Find all real  values of \(x\) such that the fraction \( \frac{7}{x^2+1} \) is negative.}
{\( \emptyset \)}{\( x \in \mathbb{R} \)}{\( x \in (-1;+\infty) \)}{\( x \in \mathbb{R}\setminus \{ \pm 1\}\)}{}{}}
\LANGpl{
\Question{
Znajdź wszystkie rzeczywiste wartości \(x\) takie, że ułamek \( \frac{7}{x^2+1} \) jest ujemny.}
{\( \emptyset \)}{\( x \in \mathbb{R} \)}{\( x \in (-1;+\infty) \)}{\( x \in \mathbb{R}\setminus \{ \pm 1\}\)}{}{}}
\LANGcs{
\Question{
Určete všechna reálná čísla  \(x\) tak, aby po jejich dosazení byl zlomek  \( \frac{7}{x^2+1} \) záporný.}
{\( \emptyset \)}{\( x \in \mathbb{R} \)}{\( x \in (-1;+\infty) \)}{\( x \in \mathbb{R}\setminus \{ \pm 1\}\)}{}{}}
\LANGsk{
\Question{
Nájdite všetky reálne čísla  \(x\) tak, aby bol po ich dosadení zlomok  \( \frac{7}{x^2+1} \) záporný.}
{\( \emptyset \)}{\( x \in \mathbb{R} \)}{\( x \in (-1;+\infty) \)}{\( x \in \mathbb{R}\setminus \{ \pm 1\}\)}{}{}}
\LANGes{
\Question{
Halla todos los valores reales de \(x\) para los que la fracción \( \frac{7}{x^2+1} \) es negativa.}
{\( \emptyset \)}{\( x \in \mathbb{R} \)}{\( x \in (-1;+\infty) \)}{\( x \in \mathbb{R}\setminus \{ \pm 1\}\)}{}{}}
\End

\Data\ID{2010015901}
\Updated{2022-11-02 09:11:10}
\Level{B}
\SubArea{Rational equations and inequalities}
\LANGen{
\Question{
Find all real values of \(x\)  for which the fraction \(\frac{2}{x+3}\)  is negative.}
{\( x < -3\)}{\( x>-3\)}{\( x< 3\)}{\(x>3\)}{}{}}
\LANGpl{
\Question{
Znajdź wszystkie rzeczywiste wartości \(x\),  dla których ułamek \(\frac{2}{x+3}\) jest ujemny.}
{\( x < -3\)}{\( x>-3\)}{\( x< 3\)}{\(x>3\)}{}{}}
\LANGcs{
\Question{
Určete všechna reálná čísla  \(x\)  tak, aby po jejich dosazení byl zlomek  \(\frac{2}{x+3}\) záporný.}
{\( x < -3\)}{\( x>-3\)}{\( x< 3\)}{\(x>3\)}{}{}}
\LANGsk{
\Question{
Určte všetky reálne čísla  \(x\)  tak, aby po ich dosadení bol zlomok  \(\frac{2}{x+3}\) záporný.}
{\( x < -3\)}{\( x>-3\)}{\( x< 3\)}{\(x>3\)}{}{}}
\LANGes{
\Question{
Calcula todos los valores reales de \(x\)  tales que la fracción \(\frac{2}{x+3}\)  sea negativa.}
{\( x < -3\)}{\( x>-3\)}{\( x< 3\)}{\(x>3\)}{}{}}
\End

\Data\ID{2010015902}
\Updated{2022-11-02 09:11:10}
\Level{B}
\SubArea{Rational equations and inequalities}
\LANGen{
\Question{
Find the solution set of the following inequality.
\[ 
 \frac{x+3}
 {1-2x} \geq 0
\]}
{\(\left[ -3; \frac12 \right) \)}{\( [ -3;\infty )\)}{\( \left(-\infty;\frac12\right)\)}{\( (-\infty;-3 ] \cup \left(\frac12;\infty\right)\)}{}{}}
\LANGpl{
\Question{
Wskaż zbiór rozwiązań nierówności.
\[ 
 \frac{x+3}
 {1-2x} \geq 0
\]}
{\(\left\langle -3; \frac12 \right) \)}{\( \langle -3;\infty )\)}{\( \left(-\infty;\frac12\right)\)}{\( (-\infty;-3 \rangle \cup \left(\frac12;\infty\right)\)}{}{}}
\LANGcs{
\Question{
Určete množinu řešení dané nerovnice.
\[ 
 \frac{x+3}
 {1-2x} \geq 0
\]}
{\(\left\langle -3; \frac12 \right) \)}{\( \langle -3;\infty )\)}{\( \left(-\infty;\frac12\right)\)}{\( (-\infty;-3 \rangle \cup \left(\frac12;\infty\right)\)}{}{}}
\LANGsk{
\Question{
Určte množinu riešení danej nerovnice.
\[ 
 \frac{x+3}
 {1-2x} \geq 0
\]}
{\(\left\langle -3; \frac12 \right) \)}{\( \langle -3;\infty )\)}{\( \left(-\infty;\frac12\right)\)}{\( (-\infty;-3 \rangle \cup \left(\frac12;\infty\right)\)}{}{}}
\LANGes{
\Question{
Halla el conjunto de soluciones de la siguiente inecuación.
\[ 
 \frac{x+3}
 {1-2x} \geq 0
\]}
{\(\left[ -3; \frac12 \right) \)}{\( [ -3;\infty )\)}{\( \left(-\infty;\frac12\right)\)}{\( (-\infty;-3 ] \cup \left(\frac12;\infty\right)\)}{}{}}
\End

\Data\ID{2010015903}
\Updated{2022-11-02 09:11:10}
\Level{B}
\SubArea{Rational equations and inequalities}
\LANGen{
\Question{
Find the solution set of the following inequality. 
\[ 
 \frac{x -2} 
{x +3} > 1
\]}
{\(( -\infty; -3) \)}{\((-3;\infty) \)}{\(( -\infty; 2) \)}{\((2;\infty) \)}{}{}}
\LANGpl{
\Question{
Wskaż zbiór rozwiązań nierówności.
\[ 
 \frac{x -2} 
{x +3} > 1
\]}
{\(( -\infty; -3) \)}{\((-3;\infty) \)}{\(( -\infty; 2) \)}{\((2;\infty) \)}{}{}}
\LANGcs{
\Question{
Určete množinu řešení dané nerovnice.
\[ 
 \frac{x -2} 
{x +3} > 1
\]}
{\(( -\infty; -3) \)}{\((-3;\infty) \)}{\(( -\infty; 2) \)}{\((2;\infty) \)}{}{}}
\LANGsk{
\Question{
Určte množinu riešení danej nerovnice.
\[ 
 \frac{x -2} 
{x +3} > 1
\]}
{\(( -\infty; -3) \)}{\((-3;\infty) \)}{\(( -\infty; 2) \)}{\((2;\infty) \)}{}{}}
\LANGes{
\Question{
Halla el conjunto de soluciones de la siguiente inecuación. 
\[ 
 \frac{x -2} 
{x +3} > 1
\]}
{\(( -\infty; -3) \)}{\((-3;\infty) \)}{\(( -\infty; 2) \)}{\((2;\infty) \)}{}{}}
\End

\Data\ID{2010015904}
\Updated{2022-11-02 09:11:14}
\Level{B}
\SubArea{Rational equations and inequalities}
\IMG{2010015901-figure0.pdf}
\LANGen{
\Question{
An inequality is solved graphically as shown in the picture. The solution is marked in red. Choose the corresponding inequality.}
{\(3\leq \frac{x-2}
 {x} \)}{\(3\geq -\frac{2}
 {x} \)}{\(3\geq \frac{x-2}
 {x} \)}{\(3\leq -\frac{2}
 {x} \)}{}{}}
\LANGpl{
\Question{
Która z nierówności odpowiada rozwiązaniu graficznemu na rysunku?}
{\(3\leq \frac{x-2}
 {x} \)}{\(3\geq -\frac{2}
 {x} \)}{\(3\geq \frac{x-2}
 {x} \)}{\(3\leq -\frac{2}
 {x} \)}{}{}}
\LANGcs{
\Question{
Které z nerovnic odpovídá grafické řešení na obrázku?}
{\(3\leq \frac{x-2}
 {x} \)}{\(3\geq -\frac{2}
 {x} \)}{\(3\geq \frac{x-2}
 {x} \)}{\(3\leq -\frac{2}
 {x} \)}{}{}}
\LANGsk{
\Question{
Ktorá z nerovníc zodpovedá grafickému riešeniu na obrázku?}
{\(3\leq \frac{x-2}
 {x} \)}{\(3\geq -\frac{2}
 {x} \)}{\(3\geq \frac{x-2}
 {x} \)}{\(3\leq -\frac{2}
 {x} \)}{}{}}
\LANGes{
\Question{
Una inecuación se resuelve gráficamente como se muestra en la imagen. La solución está marcada en rojo. Elige la inecuación correspondiente.}
{\(3\leq \frac{x-2}
 {x} \)}{\(3\geq -\frac{2}
 {x} \)}{\(3\geq \frac{x-2}
 {x} \)}{\(3\leq -\frac{2}
 {x} \)}{}{}}
\End

\Data\ID{2010015905}
\Updated{2022-11-02 09:11:14}
\Level{B}
\SubArea{Rational equations and inequalities}
\LANGen{
\Question{
Find the solution set of the following inequality.
\[ 
 \frac{-1}
 {x^2+x-20} \geq 0
\]}
{\( (-5;4)\)}{\(\emptyset\)}{\((-\infty;-5) \cup (4;\infty) \)}{\((-4;5) \)}{}{}}
\LANGpl{
\Question{
Wyznacz zbiór rozwiązań danej nierówności.
\[ 
 \frac{-1}
 {x^2+x-20} \geq 0
\]}
{\( (-5;4)\)}{\(\emptyset\)}{\((-\infty;-5) \cup (4;\infty) \)}{\((-4;5) \)}{}{}}
\LANGcs{
\Question{
Určete množinu řešení dané nerovnice.
\[ 
 \frac{-1}
 {x^2+x-20} \geq 0
\]}
{\( (-5;4)\)}{\(\emptyset\)}{\((-\infty;-5) \cup (4;\infty) \)}{\((-4;5) \)}{}{}}
\LANGsk{
\Question{
Určte množinu riešení danej nerovnice.
\[ 
 \frac{-1}
 {x^2+x-20} \geq 0
\]}
{\( (-5;4)\)}{\(\emptyset\)}{\((-\infty;-5) \cup (4;\infty) \)}{\((-4;5) \)}{}{}}
\LANGes{
\Question{
Halla el conjunto de soluciones de la siguiente inecuación.
\[ 
 \frac{-1}
 {x^2+x-20} \geq 0
\]}
{\( (-5;4)\)}{\(\emptyset\)}{\((-\infty;-5) \cup (4;\infty) \)}{\((-4;5) \)}{}{}}
\End

\Data\ID{2010015906}
\Updated{2022-11-02 09:11:14}
\Level{B}
\SubArea{Rational equations and inequalities}
\LANGen{
\Question{
Find the solution set of the inequality.
\[ \left(x^2+4\right)\left(x^2+2\right)\leq0 \]}
{\( \emptyset\)}{\( \left(-2;-\sqrt2\right)\)}{\(\mathbb{R}\)}{\( \left(-2;-\sqrt2\right) \cup \left(\sqrt2;2\right)\)}{}{}}
\LANGpl{
\Question{
Wyznacz zbiór rozwiązań danej nierówności.
\[ \left(x^2+4\right)\left(x^2+2\right)\leq0 \]}
{\( \emptyset\)}{\( \left(-2;-\sqrt2\right)\)}{\(\mathbb{R}\)}{\( \left(-2;-\sqrt2\right) \cup \left(\sqrt2;2\right)\)}{}{}}
\LANGcs{
\Question{
Určete množinu řešení dané nerovnice.
\[ \left(x^2+4\right)\left(x^2+2\right)\leq0 \]}
{\( \emptyset\)}{\( \left(-2;-\sqrt2\right)\)}{\(\mathbb{R}\)}{\( \left(-2;-\sqrt2\right) \cup \left(\sqrt2;2\right)\)}{}{}}
\LANGsk{
\Question{
Nájdite riešenie danej nerovnice.
\[ \left(x^2+4\right)\left(x^2+2\right)\leq0 \]}
{\( \emptyset\)}{\( \left(-2;-\sqrt2\right)\)}{\(\mathbb{R}\)}{\( \left(-2;-\sqrt2\right) \cup \left(\sqrt2;2\right)\)}{}{}}
\LANGes{
\Question{
Halla el conjunto de soluciones de la siguiente inecuación.
\[ \left(x^2+4\right)\left(x^2+2\right)\leq0 \]}
{\( \emptyset\)}{\( \left(-2;-\sqrt2\right)\)}{\(\mathbb{R}\)}{\( \left(-2;-\sqrt2\right) \cup \left(\sqrt2;2\right)\)}{}{}}
\End

\Data\ID{2110015907}
\Updated{2022-11-02 09:11:00}
\Level{B}
\SubArea{Rational equations and inequalities}
\LANGen{
\Question{
Identify the picture that shows the correct solution set of the following inequality.  In each picture, the set of points corresponding to the solution set is marked in red.
\[ 
 \frac{-x} 
{x + 1} > 0
\]}
{}{}{}{}{}{}}
\LANGpl{
\Question{
Wskaż rysunek, który pokazuje poprawny zbiór rozwiązań poniższej nierówności. Na każdym rysunku na czerwono zaznaczono zbiór punktów odpowiadających zbiorowi rozwiązań.
\[ 
 \frac{-x} 
{x + 1} > 0
\]}
{}{}{}{}{}{}}
\LANGcs{
\Question{
Na kterém z obrázků je správně znázorněno grafické řešení dané nerovnice? 
Řešení je vyznačeno červenou barvou.

\[ 
 \frac{-x} 
{x + 1} > 0
\]}
{}{}{}{}{}{}}
\LANGsk{
\Question{
Na ktorom z obrázkov je správne znázornené grafické riešenie danej nerovnice? 
Riešenie je vyznačené červenou farbou.

\[ 
 \frac{-x} 
{x + 1} > 0
\]}
{}{}{}{}{}{}}
\LANGes{
\Question{
Identifica la imagen que representa la solución de la siguiente inecuación. En todas las imágenes, el conjunto de soluciones se representa en rojo. 
\[ 
 \frac{-x} 
{x + 1} > 0
\]}
{}{}{}{}{}{}}

\IMGanswer{2010015901-figure1.pdf}
\IMGanswer{2010015901-figure2.pdf}
\IMGanswer{2010015901-figure3.pdf}
\IMGanswer{2010015901-figure4.pdf}\End

\Data\ID{9000018005}
\Updated{2022-04-23 05:04:11}
\Level{B}
\SubArea{Rational equations and inequalities}
\LANGen{
\Question{
Find all real values of \(x\)  for which the fraction \(\frac{3}{2-x}\)  is positive.}
{\(x < 2\)}{\(x < -2\)}{\(x > -2\)}{\(x > 2\)}{}{}}
\LANGpl{
\Question{
Wyznacz wszystkie rzeczywiste wartości dodatnie \(x\), dla których ułamek \(\frac{3}
{2-x}\) 
jest dodatni.}
{\(x < 2\)}{\(x < -2\)}{\(x > -2\)}{\(x > 2\)}{}{}}
\LANGcs{
\Question{
Určete všechna reálná čísla
\(x\) tak, aby po jejich
dosazení byl zlomek \(\frac{3}
{2-x}\) 
kladný.}
{\(x < 2\)}{\(x < -2\)}{\(x > -2\)}{\(x > 2\)}{}{}}
\LANGsk{
\Question{
Nájdite všetky reálne čísla \(x\), aby po ich dosadení bol zlomok 
 \(\frac{3}
{2-x}\) kladný.}
{\(x < 2\)}{\(x < -2\)}{\(x > -2\)}{\(x > 2\)}{}{}}
\LANGes{
\Question{
Encuentra todos los números reales \(x\) 
para los cuales la fracción \(\frac{3}
{2-x}\) 
es positiva.}
{\(x < 2\)}{\(x < -2\)}{\(x > -2\)}{\(x > 2\)}{}{}}
\End

\Data\ID{9000018105}
\Updated{2020-11-09 07:11:34}
\Level{B}
\SubArea{Rational equations and inequalities}
\LANGen{
\Question{
Find all the values of \(x\) 
which ensure that the following fraction is positive. 
\[ 
 - \frac{3}
{5 - 2x}
\]}
{\(x > \frac{5} 
{2}\)}{\(x < \frac{5} 
{2}\)}{\(x < -\frac{5}
{2}\)}{\(x > -\frac{5}
{2}\)}{}{}}
\LANGpl{
\Question{
Znajdź wszystkie wartości  \(x\), dla których następujący ułamek jest dodatni.
\[ 
 - \frac{3}
{5 - 2x}
\]}
{\(x > \frac{5} 
{2}\)}{\(x < \frac{5} 
{2}\)}{\(x < -\frac{5}
{2}\)}{\(x > -\frac{5}
{2}\)}{}{}}
\LANGcs{
\Question{
Určete všechna reálná čísla
\(x\) tak, aby po jejich
dosazení byl daný zlomek kladný.
\[- \frac{3} 
{5-2x}\]}
{\(x > \frac{5} 
{2}\)}{\(x < \frac{5} 
{2}\)}{\(x < -\frac{5}
{2}\)}{\(x > -\frac{5}
{2}\)}{}{}}
\LANGsk{
\Question{
Nájdite všetky reálne čísla \(x\), aby po ich dosadení bol zlomok kladný. 
\[ 
 - \frac{3}
{5 - 2x}
\]}
{\(x > \frac{5} 
{2}\)}{\(x < \frac{5} 
{2}\)}{\(x < -\frac{5}
{2}\)}{\(x > -\frac{5}
{2}\)}{}{}}
\LANGes{
\Question{
Encuentra todos los números reales de \(x\) 
para los cuales la siguiente fracción es positiva.  
\[ 
 - \frac{3}
{5 - 2x}
\]}
{\(x > \frac{5} 
{2}\)}{\(x < \frac{5} 
{2}\)}{\(x < -\frac{5}
{2}\)}{\(x > -\frac{5}
{2}\)}{}{}}
\End

\Data\ID{9000021804}
\Updated{2020-08-22 12:08:17}
\Level{B}
\SubArea{Rational equations and inequalities}
\LANGen{
\Question{
Solve the following inequality. 
\[ 
 \frac{1}
{x - 3}\leq \frac{1}
{2 - x}
\]}
{\(x\in (-\infty ;2)\cup \left [ \frac{5}
{2};3\right )\)}{\(x\in (-\infty ;2)\cup \left [ \frac{5}
{3};2\right ] \)}{\(x\in \left (-\infty ; \frac{5} 
{2}\right ] \cup \left (3;\infty \right )\)}{\(x\in \left [ \frac{5}
{2};\infty \right )\)}{}{}}
\LANGpl{
\Question{
Rozwiąż następującą nierówność.
\[ 
 \frac{1}
{x - 3}\leq \frac{1}
{2 - x}
\]}
{\(x\in (-\infty ;2)\cup \left [ \frac{5}
{2};3\right )\)}{\(x\in (-\infty ;2)\cup \left [ \frac{5}
{3};2\right ] \)}{\(x\in \left (-\infty ; \frac{5} 
{2}\right ] \cup \left (3;\infty \right )\)}{\(x\in \left [ \frac{5}
{2};\infty \right )\)}{}{}}
\LANGcs{
\Question{
Určete množinu řešení dané nerovnice. \[\frac{1}
{x-3}\leq \frac{1}
{2-x}\]}
{\((-\infty ;2)\cup \left \langle \frac{5}
{2};3\right )\)}{\((-\infty ;2)\cup \left \langle \frac{5}
{3};2\right \rangle \)}{\(\left (-\infty ; \frac{5} 
{2}\right \rangle \cup \left (3;\infty \right )\)}{\(\left \langle \frac{5}
{2};\infty \right )\)}{}{}}
\LANGsk{
\Question{
Vyriešte danú nerovnicu.
\[ 
 \frac{1}
{x - 3}\leq \frac{1}
{2 - x}
\]}
{\(x\in (-\infty ;2)\cup \left [ \frac{5}
{2};3\right )\)}{\(x\in (-\infty ;2)\cup \left [ \frac{5}
{3};2\right ] \)}{\(x\in \left (-\infty ; \frac{5} 
{2}\right ] \cup \left (3;\infty \right )\)}{\(x\in \left [ \frac{5}
{2};\infty \right )\)}{}{}}
\LANGes{
\Question{
Halla el conjunto de soluciones de la siguiente inecuación.  
\[ 
 \frac{1}
{x - 3}\leq \frac{1}
{2 - x}
\]}
{\(x\in (-\infty ;2)\cup \left [ \frac{5}
{2};3\right )\)}{\(x\in (-\infty ;2)\cup \left [ \frac{5}
{3};2\right ] \)}{\(x\in \left (-\infty ; \frac{5} 
{2}\right ] \cup \left (3;\infty \right )\)}{\(x\in \left [ \frac{5}
{2};\infty \right )\)}{}{}}
\End

\Data\ID{9000021806}
\Updated{2022-04-23 05:04:11}
\Level{B}
\SubArea{Rational equations and inequalities}
\LANGen{
\Question{
Solve the following inequality. 
\[ 
 \frac{1 - 3x}
 {x + 2} \geq 0
\]}
{\(x\in \left (-2; \frac{1} 
{3}\right ] \)}{\(x\in \left [ \frac{1}
{3};\infty \right )\)}{\(x\in \left (\frac{1}
{3};\infty \right )\)}{\(x\in (-\infty ;-2)\cup \left [ \frac{1}
{3};\infty \right )\)}{}{}}
\LANGpl{
\Question{
Rozwiąż następującą nierówność.
\[ 
 \frac{1 - 3x}
 {x + 2} \geq 0
\]}
{\(x\in \left (-2; \frac{1} 
{3}\right ] \)}{\(x\in \left [ \frac{1}
{3};\infty \right )\)}{\(x\in \left (\frac{1}
{3};\infty \right )\)}{\(x\in (-\infty ;-2)\cup \left [ \frac{1}
{3};\infty \right )\)}{}{}}
\LANGcs{
\Question{
Určete množinu řešení dané nerovnice.
\[\frac{1-3x}
 {x+2} \geq 0\]}
{\(\left (-2; \frac{1} 
{3}\right \rangle \)}{\(\left \langle \frac{1}
{3};\infty \right )\)}{\(\left (\frac{1}
{3};\infty \right )\)}{\((-\infty ;-2)\cup \left \langle \frac{1}
{3};\infty \right )\)}{}{}}
\LANGsk{
\Question{
Vyriešte danú nerovnicu. 
\[ 
 \frac{1 - 3x}
 {x + 2} \geq 0
\]}
{\(x\in \left (-2; \frac{1} 
{3}\right ] \)}{\(x\in \left [ \frac{1}
{3};\infty \right )\)}{\(x\in \left (\frac{1}
{3};\infty \right )\)}{\(x\in (-\infty ;-2)\cup \left [ \frac{1}
{3};\infty \right )\)}{}{}}
\LANGes{
\Question{
Halla el conjunto de soluciones de la siguiente inecuación. 
\[ 
 \frac{1 - 3x}
 {x + 2} \geq 0
\]}
{\(x\in \left (-2; \frac{1} 
{3}\right ] \)}{\(x\in \left [ \frac{1}
{3};\infty \right )\)}{\(x\in \left (\frac{1}
{3};\infty \right )\)}{\(x\in (-\infty ;-2)\cup \left [ \frac{1}
{3};\infty \right )\)}{}{}}
\End

\Data\ID{9000021809}
\Updated{2020-08-22 12:08:00}
\Level{B}
\SubArea{Rational equations and inequalities}
\LANGen{
\Question{
Solve the following inequality. 
\[ 
 \frac{2x + 4}
 {x - 1} < 1
\]}
{\(x\in (-5;1)\)}{\(x\in (-\infty ;5)\)}{\(x\in (1;\infty )\)}{\(x\in (-\infty ;-3)\cup (1;\infty )\)}{}{}}
\LANGpl{
\Question{
Rozwiąż następującą nierówność.
\[ 
 \frac{2x + 4}
 {x - 1} < 1
\]}
{\(x\in (-5;1)\)}{\(x\in (-\infty ;5)\)}{\(x\in (1;\infty )\)}{\(x\in (-\infty ;-3)\cup (1;\infty )\)}{}{}}
\LANGcs{
\Question{
Určete množinu řešení dané nerovnice.
\[\frac{2x+4}
 {x-1} < 1\]}
{\((-5;1)\)}{\((-\infty ;5)\)}{\((1;\infty )\)}{\((-\infty ;-3)\cup (1;\infty )\)}{}{}}
\LANGsk{
\Question{
Vyriešte danú nerovnicu. 
\[ 
 \frac{2x + 4}
 {x - 1} < 1
\]}
{\(x\in (-5;1)\)}{\(x\in (-\infty ;5)\)}{\(x\in (1;\infty )\)}{\(x\in (-\infty ;-3)\cup (1;\infty )\)}{}{}}
\LANGes{
\Question{
Halla el conjunto de soluciones de la siguiente inecuación.  
\[ 
 \frac{2x + 4}
 {x - 1} < 1
\]}
{\(x\in (-5;1)\)}{\(x\in (-\infty ;5)\)}{\(x\in (1;\infty )\)}{\(x\in (-\infty ;-3)\cup (1;\infty )\)}{}{}}
\End

\Data\ID{9000021810}
\Updated{2020-11-09 07:11:08}
\Level{B}
\SubArea{Rational equations and inequalities}
\LANGen{
\Question{
Find all the values of \(x\) 
for which the following expression takes on values smaller than or equal to
\(1\).
\[ 
 \frac{x + 1}
{x - 1} - \frac{1} 
{x + 1}
\]}
{\(x\in (-\infty ;-3] \cup (-1;1)\)}{\(x\in (-\infty ;-3] \)}{\(x\in (-\infty ;-1)\cup (-1;1)\cup (1;\infty )\)}{\(x\in [ - 3;-1)\)}{}{}}
\LANGpl{
\Question{
Znajdź wszystkie wartości  \(x\), dla których podane wyrażenie przyjmuje wartości mniejsze lub równe liczbie
\(1\).
\[ 
 \frac{x + 1}
{x - 1} - \frac{1} 
{x + 1}
\]}
{\(x\in (-\infty ;-3] \cup (-1;1)\)}{\(x\in (-\infty ;-3] \)}{\(x\in (-\infty ;-1)\cup (-1;1)\cup (1;\infty )\)}{\(x\in [ - 3;-1)\)}{}{}}
\LANGcs{
\Question{
Hodnota výrazu \(\frac{x+1}
{x-1} - \frac{1} 
{x+1}\) je
menší nebo rovna 1 pro \(x\) 
z množiny:}
{\((-\infty ;-3\rangle \cup (-1;1)\)}{\((-\infty ;-3\rangle \)}{\((-\infty ;-1)\cup (-1;1)\cup (1;\infty )\)}{\(\langle - 3;-1)\)}{}{}}
\LANGsk{
\Question{
Nájdite všetky hodnoty  \(x\),
pre ktoré nasledujúci výraz nadobúda hodnoty menšie alebo rovné
\(1\).
\[ 
 \frac{x + 1}
{x - 1} - \frac{1} 
{x + 1}
\]}
{\(x\in (-\infty ;-3] \cup (-1;1)\)}{\(x\in (-\infty ;-3] \)}{\(x\in (-\infty ;-1)\cup (-1;1)\cup (1;\infty )\)}{\(x\in [ - 3;-1)\)}{}{}}
\LANGes{
\Question{
Encuentra todos los números \(x\) 
para los que la siguiente expresión es menor o igual que 
\(1\).
\[ 
 \frac{x + 1}
{x - 1} - \frac{1} 
{x + 1}
\]}
{\(x\in (-\infty ;-3] \cup (-1;1)\)}{\(x\in (-\infty ;-3] \)}{\(x\in (-\infty ;-1)\cup (-1;1)\cup (1;\infty )\)}{\(x\in [ - 3;-1)\)}{}{}}
\End

\Data\ID{9000022804}
\Updated{2020-11-09 07:11:38}
\Level{B}
\SubArea{Rational equations and inequalities}
\LANGen{
\Question{
Establish the values of the real parameter
\(t\) which
ensure that the following expression is nonpositive. 
\[ 
 \frac{2} 
{2t^{2} + t - 1}
\]}
{\(\left (-1; \frac{1} 
{2}\right )\)}{\(\left [ -\frac{1}
{2};1\right ] \)}{\(\left [ -1; \frac{1} 
{2}\right ] \)}{\(\left (-\frac{1}
{2};1\right )\)}{}{}}
\LANGpl{
\Question{
Wyznacz wartości rzeczywistego parametru
\(t\),
podane wyrażenie jest niedodatnie.
\[ 
 \frac{2} 
{2t^{2} + t - 1}
\]}
{\(\left (-1; \frac{1} 
{2}\right )\)}{\(\left [ -\frac{1}
{2};1\right ] \)}{\(\left [ -1; \frac{1} 
{2}\right ] \)}{\(\left (-\frac{1}
{2};1\right )\)}{}{}}
\LANGcs{
\Question{
Množina všech takových \(t\),
pro která není zlomek \(\frac{2}
{2t^{2}+t-1}\) 
kladný, je:}
{\(\left (-1; \frac{1} 
{2}\right )\)}{\(\left \langle -\frac{1}
{2};1\right \rangle \)}{\(\left \langle -1; \frac{1} 
{2}\right \rangle \)}{\(\left (-\frac{1}
{2};1\right )\)}{}{}}
\LANGsk{
\Question{
Množina všetkých hodnôt 
\(t\), pre ktoré
daný výraz 
\( 
 \frac{2} 
{2t^{2} + t - 1}
\) nie je kladný, je:}
{\(\left (-1; \frac{1} 
{2}\right )\)}{\(\left [ -\frac{1}
{2};1\right ] \)}{\(\left [ -1; \frac{1} 
{2}\right ] \)}{\(\left (-\frac{1}
{2};1\right )\)}{}{}}
\LANGes{
\Question{
Identifica todos los valores del parámetro real \(t\) para los que la siguiente fracción no es positiva.  
\[ 
 \frac{2} 
{2t^{2} + t - 1}
\]}
{\(\left (-1; \frac{1} 
{2}\right )\)}{\(\left [ -\frac{1}
{2};1\right ] \)}{\(\left [ -1; \frac{1} 
{2}\right ] \)}{\(\left (-\frac{1}
{2};1\right )\)}{}{}}
\End

\Data\ID{9000026106}
\Updated{2022-04-23 05:04:11}
\Level{B}
\SubArea{Rational equations and inequalities}
\LANGen{
\Question{
Find the solution set of the following inequality. 
\[ 
 \frac{x + 3} 
{x - 1} > 1
\]}
{\(\left (1;\infty \right )\)}{\(\left (-\infty ;1\right )\)}{\(\left (-\infty ;3\right ] \)}{\(\left [ 3;\infty \right )\)}{}{}}
\LANGpl{
\Question{
Wyznacz zbiór rozwiązań podanej nierówności
\[ 
 \frac{x + 3} 
{x - 1} > 1
\]}
{\(\left (1;\infty \right )\)}{\(\left (-\infty ;1\right )\)}{\(\left (-\infty ;3\right ] \)}{\(\left [ 3;\infty \right )\)}{}{}}
\LANGcs{
\Question{
Určete množinu řešení dané nerovnice. \[\frac{x+3}
{x-1} > 1\]}
{\(\left (1;\infty \right )\)}{\(\left (-\infty ;1\right )\)}{\(\left (-\infty ;3\right \rangle \)}{\(\left \langle 3;\infty \right )\)}{}{}}
\LANGsk{
\Question{
Určete množinu riešení danej nerovnice.
\[ 
 \frac{x + 3} 
{x - 1} > 1
\]}
{\(\left (1;\infty \right )\)}{\(\left (-\infty ;1\right )\)}{\(\left (-\infty ;3\right ] \)}{\(\left [ 3;\infty \right )\)}{}{}}
\LANGes{
\Question{
Halla el conjunto de soluciones de la siguiente inecuación. 
\[ 
 \frac{x + 3} 
{x - 1} > 1
\]}
{\(\left (1;\infty \right )\)}{\(\left (-\infty ;1\right )\)}{\(\left (-\infty ;3\right ] \)}{\(\left [ 3;\infty \right )\)}{}{}}
\End

\Data\ID{9000026108}
\Updated{2022-04-23 05:04:11}
\Level{B}
\SubArea{Rational equations and inequalities}
\IMG{000261_08-figure0.pdf}
\LANGen{
\Question{
An inequality is solved graphically as shown in the picture. The solution is marked in red. Choose the corresponding inequality.}
{\(2\leq \frac{x+3}
 {x} \)}{\(2\geq \frac{3}
{x}\)}{\(2\geq \frac{x+3}
 {x} \)}{\(2\leq \frac{3}
{x}\)}{}{}}
\LANGpl{
\Question{
Która z podanych nierówności nawiązuje do poniższego rysunku?}
{\(2\leq \frac{x+3}
 {x} \)}{\(2\geq \frac{3}
{x}\)}{\(2\geq \frac{x+3}
 {x} \)}{\(2\leq \frac{3}
{x}\)}{}{}}
\LANGcs{
\Question{
Které z nerovnic odpovídá grafické řešení na obrázku?}
{\(2\leq \frac{x+3}
 {x} \)}{\(2\geq \frac{3}
{x}\)}{\(2\geq \frac{x+3}
 {x} \)}{\(2\leq \frac{3}
{x}\)}{}{}}
\LANGsk{
\Question{
Ktorej z nerovníc zodpovedá grafické riešenie na obrázku?}
{\(2\leq \frac{x+3}
 {x} \)}{\(2\geq \frac{3}
{x}\)}{\(2\geq \frac{x+3}
 {x} \)}{\(2\leq \frac{3}
{x}\)}{}{}}
\LANGes{
\Question{
Identifica la inecuación que corresponde a la imagen.}
{\(2\leq \frac{x+3}
 {x} \)}{\(2\geq \frac{3}
{x}\)}{\(2\geq \frac{x+3}
 {x} \)}{\(2\leq \frac{3}
{x}\)}{}{}}
\End

\Data\ID{9000033304}
\Updated{2020-08-22 08:08:06}
\Level{B}
\SubArea{Rational equations and inequalities}
\LANGen{
\Question{
Find the solution set of the following inequality. 
\[ 
 \frac{x + 4} 
{x + 2}\leq 0
\]}
{\([ - 4;-2)\)}{\((-\infty ;-4] \cup [ 2;\infty )\)}{\((-\infty ;-4)\cup (-2;\infty )\)}{\((-4;-2] \)}{}{}}
\LANGpl{
\Question{
Wyznacz zbiór rozwiązań podanej nierówności.
\[ 
 \frac{x + 4} 
{x + 2}\leq 0
\]}
{\([ - 4;-2)\)}{\((-\infty ;-4] \cup [ 2;\infty )\)}{\((-\infty ;-4)\cup (-2;\infty )\)}{\((-4;-2] \)}{}{}}
\LANGcs{
\Question{
Určete množinu řešení dané nerovnice. \[\frac{x+4}
{x+2}\leq 0\]}
{\(\langle - 4;-2)\)}{\((-\infty ;-4\rangle \cup \langle 2;\infty )\)}{\((-\infty ;-4)\cup (-2;\infty )\)}{\((-4;-2\rangle \)}{}{}}
\LANGsk{
\Question{
Určte množinu riešení danej nerovnice.
\[ 
 \frac{x + 4} 
{x + 2}\leq 0
\]}
{\([ - 4;-2)\)}{\((-\infty ;-4] \cup [ 2;\infty )\)}{\((-\infty ;-4)\cup (-2;\infty )\)}{\((-4;-2] \)}{}{}}
\LANGes{
\Question{
Halla el conjunto de soluciones de la siguiente inecuación.  
\[ 
 \frac{x + 4} 
{x + 2}\leq 0
\]}
{\([ - 4;-2)\)}{\((-\infty ;-4] \cup [ 2;\infty )\)}{\((-\infty ;-4)\cup (-2;\infty )\)}{\((-4;-2] \)}{}{}}
\End

\Data\ID{9000033305}
\Updated{2020-08-22 08:08:18}
\Level{B}
\SubArea{Rational equations and inequalities}
\LANGen{
\Question{
Find the solution set of the following inequality. 
\[ 
 \frac{2} 
{x + 1}\geq 1
\]}
{\((-1;1] \)}{\([ - 1;1)\)}{\((-\infty ;-1)\cup [ 1;\infty )\)}{\((-\infty ;1] \)}{}{}}
\LANGpl{
\Question{
Wyznacz zbiór rozwiązań podanej nierówności.
\[ 
 \frac{2} 
{x + 1}\geq 1
\]}
{\((-1;1] \)}{\([ - 1;1)\)}{\((-\infty ;-1)\cup [ 1;\infty )\)}{\((-\infty ;1] \)}{}{}}
\LANGcs{
\Question{
Určete množinu řešení dané nerovnice. \[\frac{2}
{x+1}\geq 1\]}
{\((-1;1\rangle \)}{\(\langle - 1;1)\)}{\((-\infty ;-1)\cup \langle 1;\infty )\)}{\((-\infty ;1\rangle \)}{}{}}
\LANGsk{
\Question{
Určte množinu riešení danej nerovnice.
\[ 
 \frac{2} 
{x + 1}\geq 1
\]}
{\((-1;1] \)}{\([ - 1;1)\)}{\((-\infty ;-1)\cup [ 1;\infty )\)}{\((-\infty ;1] \)}{}{}}
\LANGes{
\Question{
Halla el conjunto de soluciones de la siguiente inecuación. 
\[ 
 \frac{2} 
{x + 1}\geq 1
\]}
{\((-1;1] \)}{\([ - 1;1)\)}{\((-\infty ;-1)\cup [ 1;\infty )\)}{\((-\infty ;1] \)}{}{}}
\End

\Data\ID{9000033306}
\Updated{2020-08-22 08:08:48}
\Level{B}
\SubArea{Rational equations and inequalities}
\LANGen{
\Question{
Find the solution set of the following inequality. 
\[ 
 \frac{2} 
{3} < \frac{2 + x} 
{3 + x}
\]}
{\((-\infty ;-3)\cup (0;\infty )\)}{\(\mathbb{R}\)}{\((-3;\infty )\)}{\((-3;0)\)}{}{}}
\LANGpl{
\Question{
Wyznacz zbiór rozwiązań podanej nierówności.
\[ 
 \frac{2} 
{3} < \frac{2 + x} 
{3 + x}
\]}
{\((-\infty ;-3)\cup (0;\infty )\)}{\(\mathbb{R}\)}{\((-3;\infty )\)}{\((-3;0)\)}{}{}}
\LANGcs{
\Question{
Určete množinu řešení dané nerovnice. \[\frac{2}
{3} < \frac{2+x} 
{3+x}\]}
{\((-\infty ;-3)\cup (0;\infty )\)}{\(\mathbb{R}\)}{\((-3;\infty )\)}{\((-3;0)\)}{}{}}
\LANGsk{
\Question{
Určte množinu riešení danej nerovnice.
\[ 
 \frac{2} 
{3} < \frac{2 + x} 
{3 + x}
\]}
{\((-\infty ;-3)\cup (0;\infty )\)}{\(\mathbb{R}\)}{\((-3;\infty )\)}{\((-3;0)\)}{}{}}
\LANGes{
\Question{
Halla el conjunto de soluciones de la siguiente inecuación. 
\[ 
 \frac{2} 
{3} < \frac{2 + x} 
{3 + x}
\]}
{\((-\infty ;-3)\cup (0;\infty )\)}{\(\mathbb{R}\)}{\((-3;\infty )\)}{\((-3;0)\)}{}{}}
\End

\Data\ID{9000033307}
\Updated{2020-08-22 08:08:03}
\Level{B}
\SubArea{Rational equations and inequalities}
\LANGen{
\Question{
Find the solution set of the following inequality. 
\[ 
 \frac{4} 
{x^{2} - x - 6}\leq 0
\]}
{\((-2;3)\)}{\(\mathbb{R}\)}{\((-\infty ;-2)\cup (3;\infty )\)}{\((-3;2)\)}{}{}}
\LANGpl{
\Question{
Wyznacz zbiór rozwiązań podanej nierówności.
\[ 
 \frac{4} 
{x^{2} - x - 6}\leq 0
\]}
{\((-2;3)\)}{\(\mathbb{R}\)}{\((-\infty ;-2)\cup (3;\infty )\)}{\((-3;2)\)}{}{}}
\LANGcs{
\Question{
Určete množinu řešení dané nerovnice. \[\frac{4}
{x^{2}-x-6}\leq 0\]}
{\((-2;3)\)}{\(\mathbb{R}\)}{\((-\infty ;-2)\cup (3;\infty )\)}{\((-3;2)\)}{}{}}
\LANGsk{
\Question{
Určte množinu riešení danej nerovnice.
\[ 
 \frac{4} 
{x^{2} - x - 6}\leq 0
\]}
{\((-2;3)\)}{\(\mathbb{R}\)}{\((-\infty ;-2)\cup (3;\infty )\)}{\((-3;2)\)}{}{}}
\LANGes{
\Question{
Halla el conjunto de soluciones de la siguiente inecuación. 
\[ 
 \frac{4} 
{x^{2} - x - 6}\leq 0
\]}
{\((-2;3)\)}{\(\mathbb{R}\)}{\((-\infty ;-2)\cup (3;\infty )\)}{\((-3;2)\)}{}{}}
\End

\Data\ID{9000039005}
\Updated{2020-08-22 12:08:12}
\Level{B}
\SubArea{Rational equations and inequalities}
\LANGen{
\Question{
Find all the values of \(x\) 
for which the following expression is positive. 
\[ 
 \frac{2x - 3} 
{7 - 3x}
\]}
{\(x\in \left (\frac{3}
{2}; \frac{7} 
{3}\right )\)}{\(x\in \left (\frac{3}
{2};+\infty \right )\)}{\(x\in \left (\frac{7}
{3};+\infty \right )\)}{\(x\in (0;+\infty )\)}{}{}}
\LANGpl{
\Question{
Znajdź wszystkie wartości \(x\), dla których następujące wyrażenie jest dodatnie.
\[ 
 \frac{2x - 3} 
{7 - 3x}
\]}
{\(x\in \left (\frac{3}
{2}; \frac{7} 
{3}\right )\)}{\(x\in \left (\frac{3}
{2};+\infty \right )\)}{\(x\in \left (\frac{7}
{3};+\infty \right )\)}{\(x\in (0;+\infty )\)}{}{}}
\LANGcs{
\Question{
Pro která \(x\) 
nabývá daný zlomek  
kladných hodnot?
\[\frac{2x-3}
{7-3x}\]}
{\(x\in \left (\frac{3}
{2}; \frac{7} 
{3}\right )\)}{\(x\in \left (\frac{3}
{2};+\infty \right )\)}{\(x\in \left (\frac{7}
{3};+\infty \right )\)}{\(x\in (0;+\infty )\)}{}{}}
\LANGsk{
\Question{
Pre aké  \(x\)  nadobúda daný zlomok kladné hodnoty?
\[ 
 \frac{2x - 3} 
{7 - 3x}
\]}
{\(x\in \left (\frac{3}
{2}; \frac{7} 
{3}\right )\)}{\(x\in \left (\frac{3}
{2};+\infty \right )\)}{\(x\in \left (\frac{7}
{3};+\infty \right )\)}{\(x\in (0;+\infty )\)}{}{}}
\LANGes{
\Question{
Encuentra todos los \(x\) 
para los que la siguiente fracción es positiva. 
\[ 
 \frac{2x - 3} 
{7 - 3x}
\]}
{\(x\in \left (\frac{3}
{2}; \frac{7} 
{3}\right )\)}{\(x\in \left (\frac{3}
{2};+\infty \right )\)}{\(x\in \left (\frac{7}
{3};+\infty \right )\)}{\(x\in (0;+\infty )\)}{}{}}
\End

\Data\ID{9100033310}
\Updated{2022-07-06 12:07:35}
\Level{B}
\SubArea{Rational equations and inequalities}
\LANGen{
\Question{
Identify the picture that shows the correct solution set of the following inequality.  In each picture, the set of points corresponding to the solution set is marked in red.
\[ 
 \frac{x} 
{x - 1} < 0
\]}
{}{}{}{}{}{}}
\LANGpl{
\Question{
Który z poniższych rysunków przedstawia rozwiązanie podanej nierówności?
\[ 
 \frac{x} 
{x - 1} < 0
\] Zbiór rozwiązań oznaczono kolorem czerwonym.}
{}{}{}{}{}{}}
\LANGcs{
\Question{
Na kterém z nákresů je znázorněno grafické řešení dané nerovnice?
\[\frac{x}
{x-1} < 0\] Řešení je vyznačeno červenou barvou.}
{}{}{}{}{}{}}
\LANGsk{
\Question{
Na ktorom obrázku je znázornené grafické riešenie danej nerovnice?
\[ 
 \frac{x} 
{x - 1} < 0
\] Riešení je vyznačené červenou farbou.}
{}{}{}{}{}{}}
\LANGes{
\Question{
Identifica la imagen que representa la solución de la siguiente inecuación. 
\[ 
 \frac{x} 
{x - 1} < 0
\] En todas las imágenes, el conjunto de soluciones se representa en rojo.}
{}{}{}{}{}{}}

\IMGanswer{000333_10-figure0.pdf}
\IMGanswer{000333_10-figure1.pdf}
\IMGanswer{000333_10-figure2.pdf}
\IMGanswer{000333_10-figure3.pdf}\End

\Data\ID{1003181001}
\Updated{2020-08-22 11:08:17}
\Level{C}
\SubArea{Rational equations and inequalities}
\LANGen{
\Question{
Find the solution set of the given inequality.
\[ \frac{\left(3x^2+2\right)|x+1|}{(x-5)(1-2x)} < 0 \]}
{\( (-\infty;-1)\cup\left(-1;\frac12\right)\cup(5;\infty) \)}{\( \left(-\infty;\frac12\right)\cup(5;\infty) \)}{\( (5;\infty) \)}{\(  \left(\frac12;5\right) \)}{}{}}
\LANGpl{
\Question{
Wyznacz zbiór rozwiązań podanej nierówności.
\[ \frac{\left(3x^2+2\right)|x+1|}{(x-5)(1-2x)} < 0 \]}
{\( (-\infty;-1)\cup\left(-1;\frac12\right)\cup(5;\infty) \)}{\( \left(-\infty;\frac12\right)\cup(5;\infty) \)}{\( (5;\infty) \)}{\(  \left(\frac12;5\right) \)}{}{}}
\LANGcs{
\Question{
Určete množinu řešení dané nerovnice.
\[ \frac{\left(3x^2+2\right)|x+1|}{(x-5)(1-2x)} < 0 \]}
{\( (-\infty;-1)\cup\left(-1;\frac12\right)\cup(5;\infty) \)}{\( \left(-\infty;\frac12\right)\cup(5;\infty) \)}{\( (5;\infty) \)}{\(  \left(\frac12;5\right) \)}{}{}}
\LANGsk{
\Question{
Určte množinu riešení danej nerovnice.
\[ \frac{\left(3x^2+2\right)|x+1|}{(x-5)(1-2x)} < 0 \]}
{\( (-\infty;-1)\cup\left(-1;\frac12\right)\cup(5;\infty) \)}{\( \left(-\infty;\frac12\right)\cup(5;\infty) \)}{\( (5;\infty) \)}{\(  \left(\frac12;5\right) \)}{}{}}
\LANGes{
\Question{
Halla el conjunto de soluciones de la siguiente inecuación. 
\[ \frac{\left(3x^2+2\right)|x+1|}{(x-5)(1-2x)} < 0 \]}
{\( (-\infty;-1)\cup\left(-1;\frac12\right)\cup(5;\infty) \)}{\( \left(-\infty;\frac12\right)\cup(5;\infty) \)}{\( (5;\infty) \)}{\(  \left(\frac12;5\right) \)}{}{}}
\End

\Data\ID{1003181002}
\Updated{2022-11-17 01:11:30}
\Level{C}
\SubArea{Rational equations and inequalities}
\LANGen{
\Question{
Find the solution set of the given inequality.
\[ \frac{|13-5x|}{x(x-2)} \geq 0 \]}
{\( (-\infty;0)\cup(2;\infty) \)}{\( (-\infty;0)\cup\left(\frac35;2\right)\cup(2;\infty) \)}{\( (-\infty;0]\cup[2;\infty) \)}{\( (2;\infty) \)}{}{}}
\LANGpl{
\Question{
Wyznacz zbiór rozwiązań podanej nierówności.
\[ \frac{|13-5x|}{x(x-2)} \geq 0 \]}
{\( (-\infty;0)\cup(2;\infty) \)}{\( (-\infty;0)\cup\left(\frac35;2\right)\cup(2;\infty) \)}{\( (-\infty;0\rangle\cup\langle2;\infty) \)}{\( (2;\infty) \)}{}{}}
\LANGcs{
\Question{
Určete množinu řešení dané nerovnice.
\[ \frac{|13-5x|}{x(x-2)} \geq 0 \]}
{\( (-\infty;0)\cup(2;\infty) \)}{\( (-\infty;0)\cup\left(\frac35;2\right)\cup(2;\infty) \)}{\( (-\infty;0\rangle\cup\langle2;\infty) \)}{\( (2;\infty) \)}{}{}}
\LANGsk{
\Question{
Určte množinu riešení danej nerovnice.
\[ \frac{|13-5x|}{x(x-2)} \geq 0 \]}
{\( (-\infty;0)\cup(2;\infty) \)}{\( (-\infty;0)\cup\left(\frac35;2\right)\cup(2;\infty) \)}{\( (-\infty;0\rangle\cup\langle2;\infty) \)}{\( (2;\infty) \)}{}{}}
\LANGes{
\Question{
Halla el conjunto de soluciones de la siguiente inecuación. \[ \frac{|13-5x|}{x(x-2)} \geq 0 \]}
{\( (-\infty;0)\cup(2;\infty) \)}{\( (-\infty;0)\cup\left(\frac35;2\right)\cup(2;\infty) \)}{\( (-\infty;0]\cup[2;\infty) \)}{\( (2;\infty) \)}{}{}}
\End

\Data\ID{1003181003}
\Updated{2020-08-22 11:08:33}
\Level{C}
\SubArea{Rational equations and inequalities}
\LANGen{
\Question{
Find the solution set of the given inequality.
\[ \frac{\left(x^2-1\right)(x+3)}{|x+1|2x} \leq 0 \]}
{\( [-3;-1)\cup(0;1] \)}{\( (-3; -1)\cup(0; 1) \)}{\( [-3;0)\cup[1;\infty) \)}{\( (-\infty; -3] \)}{}{}}
\LANGpl{
\Question{
Wyznacz zbiór rozwiązań podanej nierówności.
\[ \frac{\left(x^2-1\right)(x+3)}{|x+1|2x} \leq 0 \]}
{\( \langle-3;-1)\cup(0;1\rangle \)}{\( (-3; -1)\cup(0; 1) \)}{\( \langle-3;0)\cup\langle1;\infty) \)}{\( (-\infty; -3\rangle \)}{}{}}
\LANGcs{
\Question{
Určete množinu řešení dané nerovnice.
\[ \frac{\left(x^2-1\right)(x+3)}{|x+1|2x} \leq 0 \]}
{\( \langle-3;-1)\cup(0;1\rangle \)}{\( (-3; -1)\cup(0; 1) \)}{\( \langle-3;0)\cup\langle1;\infty) \)}{\( (-\infty; -3\rangle \)}{}{}}
\LANGsk{
\Question{
Určte množinu riešení danej nerovnice.
\[ \frac{\left(x^2-1\right)(x+3)}{|x+1|2x} \leq 0 \]}
{\( \langle-3;-1)\cup(0;1\rangle \)}{\( (-3; -1)\cup(0; 1) \)}{\( \langle-3;0)\cup\langle1;\infty) \)}{\( (-\infty; -3\rangle \)}{}{}}
\LANGes{
\Question{
Halla el conjunto de soluciones de la siguiente inecuación. 
\[ \frac{\left(x^2-1\right)(x+3)}{|x+1|2x} \leq 0 \]}
{\( [-3;-1)\cup(0;1] \)}{\( (-3; -1)\cup(0; 1) \)}{\( [-3;0)\cup[1;\infty) \)}{\( (-\infty; -3] \)}{}{}}
\End

\Data\ID{1003181004}
\Updated{2022-04-23 05:04:48}
\Level{C}
\SubArea{Rational equations and inequalities}
\LANGen{
\Question{
Find the solution set of the given inequality.
\[ \frac{|5-2x|}{1-x} > 1 \]}
{\( (-\infty;1) \)}{\( (-\infty;1)\cup\left(\frac52;\infty\right) \)}{\( \left(-\infty;\frac52\right) \)}{\( \left(\frac52;\infty\right) \)}{}{}}
\LANGpl{
\Question{
Wyznacz zbiór rozwiązań podanej nierówności.
\[ \frac{|5-2x|}{1-x} > 1 \]}
{\( (-\infty;1) \)}{\( (-\infty;1)\cup\left(\frac52;\infty\right) \)}{\( \left(-\infty;\frac52\right) \)}{\( \left(\frac52;\infty\right) \)}{}{}}
\LANGcs{
\Question{
Určete množinu řešení dané nerovnice.
\[ \frac{|5-2x|}{1-x} > 1 \]}
{\( (-\infty;1) \)}{\( (-\infty;1)\cup\left(\frac52;\infty\right) \)}{\( \left(-\infty;\frac52\right) \)}{\( \left(\frac52;\infty\right) \)}{}{}}
\LANGsk{
\Question{
Určte množinu riešení danej nerovnice.
\[ \frac{|5-2x|}{1-x} > 1 \]}
{\( (-\infty;1) \)}{\( (-\infty;1)\cup\left(\frac52;\infty\right) \)}{\( \left(-\infty;\frac52\right) \)}{\( \left(\frac52;\infty\right) \)}{}{}}
\LANGes{
\Question{
Halla el conjunto de soluciones de la siguiente inecuación. 
\[ \frac{|5-2x|}{1-x} > 1 \]}
{\( (-\infty;1) \)}{\( (-\infty;1)\cup\left(\frac52;\infty\right) \)}{\( \left(-\infty;\frac52\right) \)}{\( \left(\frac52;\infty\right) \)}{}{}}
\End

\Data\ID{1003181005}
\Updated{2022-04-23 05:04:48}
\Level{C}
\SubArea{Rational equations and inequalities}
\LANGen{
\Question{
Find the solution set of the given inequality.
\[ \frac{3-x}{|x-2|}\leq 1 \]}
{\( \left[\frac52;\infty\right) \)}{\( (-\infty;2)\cup\left[\frac52;\infty\right) \)}{\( (-\infty;0]\cup(3;\infty) \)}{\( (2;\infty) \)}{}{}}
\LANGpl{
\Question{
Wyznacz zbiór rozwiązań podanej nierówności.
\[ \frac{3-x}{|x-2|}\leq 1 \]}
{\( \left\langle\frac52;\infty\right) \)}{\( (-\infty;2)\cup\left\langle\frac52;\infty\right) \)}{\( (-\infty;0\rangle\cup(3;\infty) \)}{\( (2;\infty) \)}{}{}}
\LANGcs{
\Question{
Určete množinu řešení dané nerovnice.
\[ \frac{3-x}{|x-2|}\leq 1 \]}
{\( \left\langle\frac52;\infty\right) \)}{\( (-\infty;2)\cup\left\langle\frac52;\infty\right) \)}{\( (-\infty;0\rangle\cup(3;\infty) \)}{\( (2;\infty) \)}{}{}}
\LANGsk{
\Question{
Určte množinu riešení danej nerovnice.
\[ \frac{3-x}{|x-2|}\leq 1 \]}
{\( \left\langle\frac52;\infty\right) \)}{\( (-\infty;2)\cup\left\langle\frac52;\infty\right) \)}{\( (-\infty;0\rangle\cup(3;\infty) \)}{\( (2;\infty) \)}{}{}}
\LANGes{
\Question{
Halla el conjunto de soluciones de la siguiente inecuación. 
\[ \frac{3-x}{|x-2|}\leq 1 \]}
{\( \left[\frac52;\infty\right) \)}{\( (-\infty;2)\cup\left[\frac52;\infty\right) \)}{\( (-\infty;0]\cup(3;\infty) \)}{\( (2;\infty) \)}{}{}}
\End

\Data\ID{1003181006}
\Updated{2020-08-22 11:08:21}
\Level{C}
\SubArea{Rational equations and inequalities}
\LANGen{
\Question{
Find the solution set of the given inequality.
\[ \frac{4x-1}{|2x+1|} \leq -(2x+1) \]}
{\( \left(-\infty;-\frac12\right)\cup\left(-\frac12;0\right] \)}{\( \left(-\infty; 0\right] \)}{\( \left(-\frac12; 0\right] \)}{\( \left(-\infty;-\frac12\right) \)}{}{}}
\LANGpl{
\Question{
Wyznacz zbiór rozwiązań podanej nierówności.
\[ \frac{4x-1}{|2x+1|} \leq -(2x+1) \]}
{\( \left(-\infty;-\frac12\right)\cup\left(-\frac12;0\right\rangle \)}{\( \left(-\infty; 0\right\rangle \)}{\( \left(-\frac12; 0\right\rangle \)}{\( \left(-\infty;-\frac12\right) \)}{}{}}
\LANGcs{
\Question{
Určete množinu řešení dané nerovnice.
\[ \frac{4x-1}{|2x+1|} \leq -(2x+1) \]}
{\( \left(-\infty;-\frac12\right)\cup\left(-\frac12;0\right\rangle \)}{\( \left(-\infty; 0\right\rangle \)}{\( \left(-\frac12; 0\right\rangle \)}{\( \left(-\infty;-\frac12\right) \)}{}{}}
\LANGsk{
\Question{
Určte množinu riešení danej nerovnice.
\[ \frac{4x-1}{|2x+1|} \leq -(2x+1) \]}
{\( \left(-\infty;-\frac12\right)\cup\left(-\frac12;0\right\rangle \)}{\( \left(-\infty; 0\right\rangle \)}{\( \left(-\frac12; 0\right\rangle \)}{\( \left(-\infty;-\frac12\right) \)}{}{}}
\LANGes{
\Question{
Halla el conjunto de soluciones de la siguiente inecuación. 
\[ \frac{4x-1}{|2x+1|} \leq -(2x+1) \]}
{\( \left(-\infty;-\frac12\right)\cup\left(-\frac12;0\right] \)}{\( \left(-\infty; 0\right] \)}{\( \left(-\frac12; 0\right] \)}{\( \left(-\infty;-\frac12\right) \)}{}{}}
\End

\Data\ID{1003187305}
\Updated{2020-11-09 07:11:28}
\Level{C}
\SubArea{Rational equations and inequalities}
\LANGen{
\Question{
Which of the given numbers are the solutions of the equation \( \frac{|2x|-4}{1-|x|}+\frac{12}7=0 \)?}
{\( -8;8 \)}{\( -\frac8{13};\frac8{13} \)}{\( -\frac{20}{13};\frac{20}{13} \)}{\( -20; 20 \)}{}{}}
\LANGpl{
\Question{
Które z podanych liczb są rozwiązaniem równania \( \frac{|2x|-4}{1-|x|}+\frac{12}7=0 \)?}
{\( -8;8 \)}{\( -\frac8{13};\frac8{13} \)}{\( -\frac{20}{13};\frac{20}{13} \)}{\( -20; 20 \)}{}{}}
\LANGcs{
\Question{
Které z daných čísel jsou řešením rovnice \( \frac{|2x|-4}{1-|x|}+\frac{12}7=0 \)?}
{\( -8;8 \)}{\( -\frac8{13};\frac8{13} \)}{\( -\frac{20}{13};\frac{20}{13} \)}{\( -20; 20 \)}{}{}}
\LANGsk{
\Question{
Ktoré z daných čísel sú riešením rovnice \( \frac{|2x|-4}{1-|x|}+\frac{12}7=0 \)?}
{\( -8;8 \)}{\( -\frac8{13};\frac8{13} \)}{\( -\frac{20}{13};\frac{20}{13} \)}{\( -20; 20 \)}{}{}}
\LANGes{
\Question{
¿Cuáles de los siguientes números son soluciones de la ecuación \( \frac{|2x|-4}{1-|x|}+\frac{12}7=0 \)?}
{\( -8;8 \)}{\( -\frac8{13};\frac8{13} \)}{\( -\frac{20}{13};\frac{20}{13} \)}{\( -20; 20 \)}{}{}}
\End

\Data\ID{2000005302}
\Updated{2022-02-15 07:02:33}
\Level{C}
\SubArea{Rational equations and inequalities}
\IMG{2100005301_8-figure0_0.pdf}
\LANGen{
\Question{
Use the graphs of the functions  \(f(x)=x^2-x-6\) and \(g(x)=x-3\) to find the solution set of the following equation.
\[\frac{x^2-x-6}{x-3}=1\]}
{\( \{ -1\} \)}{\( \{ 3\} \)}{\( \{ -1;3\} \)}{\( \{ -2;3\} \)}{}{}}
\LANGpl{
\Question{
Dane są wykresy funkcji   \(f(x)=x^2-x-6\) i \(g(x)=x-3\). Wyznacz zbiór rozwiązań równania. 
\[\frac{x^2-x-6}{x-3}=1\]}
{\( \{ -1\} \)}{\( \{ 3\} \)}{\( \{ -1;3\} \)}{\( \{ -2;3\} \)}{}{}}
\LANGcs{
\Question{
Pomocí grafů funkcí  \(f(x)=x^2-x-6\) a \(g(x)=x-3\) určete množinu řešení dané rovnice.
\[\frac{x^2-x-6}{x-3}=1\]}
{\( \{ -1\} \)}{\( \{ 3\} \)}{\( \{ -1;3\} \)}{\( \{ -2;3\} \)}{}{}}
\LANGsk{
\Question{
Pomocou grafov funkcií  \(f(x)=x^2-x-6\) a \(g(x)=x-3\) určte množinu riešení danej rovnice.
\[\frac{x^2-x-6}{x-3}=1\]}
{\( \{ -1\} \)}{\( \{ 3\} \)}{\( \{ -1;3\} \)}{\( \{ -2;3\} \)}{}{}}
\LANGes{
\Question{
Usa las gráficas de las funciones \(f(x)=x^2-x-6\) y \(g(x)=x-3\) para encontrar el conjunto de soluciones de la siguiente ecuación. 
\[\frac{x^2-x-6}{x-3}=1\]}
{\( \{ -1\} \)}{\( \{ 3\} \)}{\( \{ -1;3\} \)}{\( \{ -2;3\} \)}{}{}}
\End

\Data\ID{2000005303}
\Updated{2022-02-15 07:02:33}
\Level{C}
\SubArea{Rational equations and inequalities}
\IMG{2100005301_8-figure1.pdf}
\LANGen{
\Question{
Use the graphs of the functions  \( f(x)=x^2-5x+4\) and \(g(x)=x^2+2x-3\) to find the solution set of the following equation.
\[ \frac{x^2-5x+4}{x^2+2x-3}=1\]}
{\( \emptyset\)}{\( \{1\}\)}{\( \{-3;4\}\)}{\( \{-3;1;4\}\)}{}{}}
\LANGpl{
\Question{
Dane są wykresy funkcji   \( f(x)=x^2-5x+4\) i \(g(x)=x^2+2x-3\). Wyznacz zbiór rozwiązań równania. 
\[ \frac{x^2-5x+4}{x^2+2x-3}=1\]}
{\( \emptyset\)}{\( \{1\}\)}{\( \{-3;4\}\)}{\( \{-3;1;4\}\)}{}{}}
\LANGcs{
\Question{
Pomocí grafů funkcí  \( f(x)=x^2-5x+4\) a \(g(x)=x^2+2x-3\) určete množinu řešení dané rovnice.
\[ \frac{x^2-5x+4}{x^2+2x-3}=1\]}
{\( \emptyset\)}{\( \{1\}\)}{\( \{-3;4\}\)}{\( \{-3;1;4\}\)}{}{}}
\LANGsk{
\Question{
Pomocou grafov funkcií  \( f(x)=x^2-5x+4\) a \(g(x)=x^2+2x-3\) určte množinu riešení danej rovnice.
\[ \frac{x^2-5x+4}{x^2+2x-3}=1\]}
{\( \emptyset\)}{\( \{1\}\)}{\( \{-3;4\}\)}{\( \{-3;1;4\}\)}{}{}}
\LANGes{
\Question{
Usa las gráficas de las funciones \( f(x)=x^2-5x+4\) y \(g(x)=x^2+2x-3\) para encontrar el conjunto de soluciones de la siguiente ecuación. \[\frac{x^2-5x+4}{x^2+2x-3}=1\]}
{\( \emptyset\)}{\( \{1\}\)}{\( \{-3;4\}\)}{\( \{-3;1;4\}\)}{}{}}
\End

\Data\ID{2000005306}
\Updated{2022-02-15 07:02:30}
\Level{C}
\SubArea{Rational equations and inequalities}
\IMG{2100005301_8-figure2.pdf}
\LANGen{
\Question{
Use the graphs of the functions \(f(x)=x^2-4\) and \(g(x)=x+2\) to find the solution set of the following inequality.
\[\frac{x^2-4}{x+2} \geq 0\]}
{\( x \in [ 2;+\infty) \)}{\( x \in ( 2;+\infty) \)}{\( x \in (-\infty;-2) \cup ( 2;+\infty) \)}{\( x \in (-\infty;-2] \cup [ 2;+\infty) \)}{}{}}
\LANGpl{
\Question{
Dane są wykresy funkcji \(f(x)=x^2-4\) i  \(g(x)=x+2\). Wyznacz zbiór rozwiązań nierówności. 
\[\frac{x^2-4}{x+2} \geq 0\]}
{\( x \in \langle 2;+\infty) \)}{\( x \in ( 2;+\infty) \)}{\( x \in (-\infty;-2) \cup ( 2;+\infty) \)}{\( x \in (-\infty;-2\rangle \cup \langle 2;+\infty) \)}{}{}}
\LANGcs{
\Question{
Pomocí grafů funkcí \(f(x)=x^2-4\) a \(g(x)=x+2\) určete množinu řešení nerovnice.
\[\frac{x^2-4}{x+2} \geq 0\]}
{\( x \in \langle 2;+\infty) \)}{\( x \in ( 2;+\infty) \)}{\( x \in (-\infty;-2) \cup ( 2;+\infty) \)}{\( x \in (-\infty;-2\rangle \cup \langle 2;+\infty) \)}{}{}}
\LANGsk{
\Question{
Pomocou grafov funkcií \(f(x)=x^2-4\) a \(g(x)=x+2\) určte množinu riešení nerovnice.
\[\frac{x^2-4}{x+2} \geq 0\]}
{\( x \in \langle 2;+\infty) \)}{\( x \in ( 2;+\infty) \)}{\( x \in (-\infty;-2) \cup ( 2;+\infty) \)}{\( x \in (-\infty;-2\rangle \cup \langle 2;+\infty) \)}{}{}}
\LANGes{
\Question{
Usa las gráficas de las funciones \(f(x)=x^2-4\) y \(g(x)=x+2\) para encontrar el conjunto de soluciones de la siguiente desigualdad. 
\[\frac{x^2-4}{x+2} \geq 0\]}
{\( x \in [ 2;+\infty) \)}{\( x \in ( 2;+\infty) \)}{\( x \in (-\infty;-2) \cup ( 2;+\infty) \)}{\( x \in (-\infty;-2] \cup [ 2;+\infty) \)}{}{}}
\End

\Data\ID{2000005307}
\Updated{2022-02-15 07:02:30}
\Level{C}
\SubArea{Rational equations and inequalities}
\IMG{2100005301_8-figure3.pdf}
\LANGen{
\Question{
Use the graphs of the functions \(f(x)=x+3\) and \(g(x)=x-1\) to find the solution set of the following inequality.
\[\frac{x+3}{x-1} \leq 0\]}
{\( x \in [ -3;1) \)}{\( x \in [ -3;1] \)}{\( x \in (1;+\infty) \)}{\( x \in (-1;+\infty) \)}{}{}}
\LANGpl{
\Question{
Dane są wykresy funkcji  \(f(x)=x+3\) i \(g(x)=x-1\). Wyznacz zbiór rozwiązań nierówności. 
\[\frac{x+3}{x-1} \leq 0\]}
{\( x \in \langle -3;1) \)}{\( x \in \langle -3;1\rangle \)}{\( x \in (1;+\infty) \)}{\( x \in (-1;+\infty) \)}{}{}}
\LANGcs{
\Question{
Pomocí grafů funkcí  \(f(x)=x+3\) a \(g(x)=x-1\) určete množinu řešení nerovnice.
\[\frac{x+3}{x-1} \leq 0\]}
{\( x \in \langle -3;1) \)}{\( x \in \langle -3;1\rangle \)}{\( x \in (1;+\infty) \)}{\( x \in (-1;+\infty) \)}{}{}}
\LANGsk{
\Question{
Pomocou grafov funkcií  \(f(x)=x+3\) a \(g(x)=x-1\) určte množinu riešení nerovnice.
\[\frac{x+3}{x-1} \leq 0\]}
{\( x \in \langle -3;1) \)}{\( x \in \langle -3;1\rangle \)}{\( x \in (1;+\infty) \)}{\( x \in (-1;+\infty) \)}{}{}}
\LANGes{
\Question{
Usa las gráficas de las funciones \(f(x)=x+3\) y \(g(x)=x-1\) para encontrar el conjunto de soluciones de la siguiente inecuación. 
\[\frac{x+3}{x-1} \leq 0\]}
{\( x \in [ -3;1) \)}{\( x \in [ -3;1] \)}{\( x \in (1;+\infty) \)}{\( x \in (-1;+\infty) \)}{}{}}
\End

\Data\ID{2000005308}
\Updated{2022-02-15 07:02:30}
\Level{C}
\SubArea{Rational equations and inequalities}
\IMG{2100005301_8-figure4.pdf}
\LANGen{
\Question{
Use the graphs of the functions \(f(x)=x-3\) and \(g(x)=4-x\) to find the solution set of the following inequality.
\[\frac{x-3}{4-x} < 0\]}
{\( x \in (-\infty;3) \cup (4;+\infty) \)}{\( x \in (3;4)  \)}{\( x \in (-\infty;3)  \)}{\( x \in (-\infty;0)  \)}{}{}}
\LANGpl{
\Question{
Dane są wykresy funkcji  \(f(x)=x-3\) i \(g(x)=4-x\). Wyznacz zbiór rozwiązań nierówności. 
\[\frac{x-3}{4-x} < 0\]}
{\( x \in (-\infty;3) \cup (4;+\infty) \)}{\( x \in (3;4)  \)}{\( x \in (-\infty;3)  \)}{\( x \in (-\infty;0)  \)}{}{}}
\LANGcs{
\Question{
Pomocí grafů funkcí  \(f(x)=x-3\) a \(g(x)=4-x\) určete množinu řešení nerovnice.
\[\frac{x-3}{4-x} < 0\]}
{\( x \in (-\infty;3) \cup (4;+\infty) \)}{\( x \in (3;4)  \)}{\( x \in (-\infty;3)  \)}{\( x \in (-\infty;0)  \)}{}{}}
\LANGsk{
\Question{
Pomocou grafov funkcií  \(f(x)=x-3\) a \(g(x)=4-x\) určte množinu riešení nerovnice.
\[\frac{x-3}{4-x} < 0\]}
{\( x \in (-\infty;3) \cup (4;+\infty) \)}{\( x \in (3;4)  \)}{\( x \in (-\infty;3)  \)}{\( x \in (-\infty;0)  \)}{}{}}
\LANGes{
\Question{
Usa las gráficas de las funciones \(f(x)=x-3\) y \(g(x)=4-x\) para encontrar el conjunto de soluciones de la siguiente inecuación.
\[\frac{x-3}{4-x} < 0\]}
{\( x \in (-\infty;3) \cup (4;+\infty) \)}{\( x \in (3;4)  \)}{\( x \in (-\infty;3)  \)}{\( x \in (-\infty;0)  \)}{}{}}
\End

\Data\ID{2010011901}
\Updated{2022-11-02 09:11:14}
\Level{C}
\SubArea{Rational equations and inequalities}
\LANGen{
\Question{
Find the solution set of the given inequality.
\[ \frac{|4x+1|}{x(x+3)} \leq 0 \]}
{\( (-3;0) \)}{\( [ -3;0 ] \)}{\( (-\infty;-3)\cup(0;\infty) \)}{\(\left(-3;-\frac14 \right]\)}{}{}}
\LANGpl{
\Question{
Wskaż zbiór rozwiązań nierówności.
\[ \frac{|4x+1|}{x(x+3)} \leq 0 \]}
{\( (-3;0) \)}{\( \langle -3;0 \rangle \)}{\( (-\infty;-3)\cup(0;\infty) \)}{\(\left(-3;-\frac14 \right\rangle\)}{}{}}
\LANGcs{
\Question{
Určete množinu řešení dané nerovnice.
\[ \frac{|4x+1|}{x(x+3)} \leq 0 \]}
{\( (-3;0) \)}{\( \langle -3;0 \rangle \)}{\( (-\infty;-3)\cup(0;\infty) \)}{\(\left(-3;-\frac14 \right\rangle\)}{}{}}
\LANGsk{
\Question{
Určte množinu riešení danej nerovnice.
\[ \frac{|4x+1|}{x(x+3)} \leq 0 \]}
{\( (-3;0) \)}{\( \langle -3;0 \rangle \)}{\( (-\infty;-3)\cup(0;\infty) \)}{\(\left(-3;-\frac14 \right\rangle\)}{}{}}
\LANGes{
\Question{
Halla el conjunto de soluciones de la siguiente inecuación.
\[ \frac{|4x+1|}{x(x+3)} \leq 0 \]}
{\( (-3;0) \)}{\( [ -3;0 ] \)}{\( (-\infty;-3)\cup(0;\infty) \)}{\(\left(-3;-\frac14 \right]\)}{}{}}
\End

\Data\ID{2010011902}
\Updated{2022-11-02 09:11:14}
\Level{C}
\SubArea{Rational equations and inequalities}
\LANGen{
\Question{
Find the solution set of the given inequality.
\[ \frac{|3x+4|}{2-x} > 1 \]}
{\( (-\infty;-3)\cup\left(-\frac12;2\right) \)}{\( (2;\infty)\)}{\( (-\infty;-3)\cup(2;\infty) \)}{\((-3;-\frac12 ]\)}{}{}}
\LANGpl{
\Question{
Wskaż zbiór rozwiązań nierówności.
\[ \frac{|3x+4|}{2-x} > 1 \]}
{\( (-\infty;-3)\cup\left(-\frac12;2\right) \)}{\( (2;\infty)\)}{\( (-\infty;-3)\cup(2;\infty) \)}{\((-3;-\frac12 \rangle\)}{}{}}
\LANGcs{
\Question{
Určete množinu řešení dané nerovnice.
\[ \frac{|3x+4|}{2-x} > 1 \]}
{\( (-\infty;-3)\cup\left(-\frac12;2\right) \)}{\( (2;\infty)\)}{\( (-\infty;-3)\cup(2;\infty) \)}{\((-3;-\frac12 \rangle\)}{}{}}
\LANGsk{
\Question{
Určte riešenie danej nerovnice.
\[ \frac{|3x+4|}{2-x} > 1 \]}
{\( (-\infty;-3)\cup\left(-\frac12;2\right) \)}{\( (2;\infty)\)}{\( (-\infty;-3)\cup(2;\infty) \)}{\((-3;-\frac12 \rangle\)}{}{}}
\LANGes{
\Question{
Halla el conjunto de soluciones de la siguiente inecuación.
\[ \frac{|3x+4|}{2-x} > 1 \]}
{\( (-\infty;-3)\cup\left(-\frac12;2\right) \)}{\( (2;\infty)\)}{\( (-\infty;-3)\cup(2;\infty) \)}{\((-3;-\frac12 ]\)}{}{}}
\End

\Data\ID{2010011903}
\Updated{2022-11-02 09:11:14}
\Level{C}
\SubArea{Rational equations and inequalities}
\LANGen{
\Question{
Find the solution set of the given inequality.
\[ \frac{x+1}{|3-x|} < 1 \]}
{\( (-\infty;1) \)}{\( [ -1;3)\)}{\([ 1,3)\cup(3;\infty) \)}{\((-\infty;-1 ]\)}{}{}}
\LANGpl{
\Question{
Wskaż zbiór rozwiązań nierówności.
\[ \frac{x+1}{|3-x|} < 1 \]}
{\( (-\infty;1) \)}{\( \langle -1;3)\)}{\(\langle 1,3)\cup(3;\infty) \)}{\((-\infty;-1 \rangle\)}{}{}}
\LANGcs{
\Question{
Určete množinu řešení dané nerovnice.
\[ \frac{x+1}{|3-x|} < 1 \]}
{\( (-\infty;1) \)}{\( \langle -1;3)\)}{\(\langle 1,3)\cup(3;\infty) \)}{\((-\infty;-1 \rangle\)}{}{}}
\LANGsk{
\Question{
Určte riešenie danej nerovnice.
\[ \frac{x+1}{|3-x|} < 1 \]}
{\( (-\infty;1) \)}{\( \langle -1;3)\)}{\(\langle 1,3)\cup(3;\infty) \)}{\((-\infty;-1 \rangle\)}{}{}}
\LANGes{
\Question{
Halla el conjunto de soluciones de la siguiente inecuación.
\[ \frac{x+1}{|3-x|} < 1 \]}
{\( (-\infty;1) \)}{\( [ -1;3)\)}{\([ 1,3)\cup(3;\infty) \)}{\((-\infty;-1 ]\)}{}{}}
\End

\Data\ID{2110011904}
\Updated{2022-11-02 09:11:56}
\Level{C}
\SubArea{Rational equations and inequalities}
\LANGen{
\Question{
Identify the picture that shows the correct solution set of the following inequality.  In each picture, the set of points corresponding to the solution set is marked in red.
\[ 
 \frac{2} 
{x}\geq x-1
\]}
{}{}{}{}{}{}}
\LANGpl{
\Question{
Wskaż rysunek, który pokazuje poprawny zbiór rozwiązań poniższej nierówności. Na każdym rysunku na czerwono zaznaczono zbiór punktów odpowiadających zbiorowi rozwiązań.
\[ 
 \frac{2} 
{x}\geq x-1
\]}
{}{}{}{}{}{}}
\LANGcs{
\Question{
Na kterém obrázku je znázorněno grafické řešení dané nerovnice?
\[ 
 \frac{2} 
{x}\geq x-1
\]
Řešení je vyznačeno červenou barvou.}
{}{}{}{}{}{}}
\LANGsk{
\Question{
Na ktorom obrázku je znázornené grafické riešenie danej nerovnice?
\[ 
 \frac{2} 
{x}\geq x-1
\]
Riešenie je vyznačené červenou farbou.}
{}{}{}{}{}{}}
\LANGes{
\Question{
Identifica la imagen que muestra el conjunto de soluciones correcto de la siguiente inecuación.  En cada imagen, este conjunto de puntos está marcado en rojo.
\[ 
 \frac{2} 
{x}\geq x-1
\]}
{}{}{}{}{}{}}

\IMGanswer{2010011901-figure0.pdf}
\IMGanswer{2010011901-figure1_0.pdf}
\IMGanswer{2010011901-figure2.pdf}
\IMGanswer{2010011901-figure3.pdf}\End

\Data\ID{9000025807}
\Updated{2020-08-22 12:08:09}
\Level{C}
\SubArea{Rational equations and inequalities}
\LANGen{
\Question{
In the following list identify a true statement on the function
\(f\).
\[ 
 f(x) = \frac{-2(3x + 1)} 
{(2x + 3)(2 - x)}
\]}
{\(f(x) > 0 \iff x\in \left (-\frac{3}
{2};-\frac{1}
{3}\right )\cup (2;\infty )\)}{\(f(x) > 0 \iff x\in \left (-\infty ;-\frac{3}
{2}\right )\cup \left (-\frac{1}
{3};2\right )\)}{\(f(x) > 0 \iff x\in \left (-\frac{3}
{2};2\right )\)}{\(f(x) > 0 \iff x\in \left (-\infty ;-\frac{3}
{2}\right )\cup (2;\infty )\)}{}{}}
\LANGpl{
\Question{
Które z poniższych stwierdzeń dotyczących funkcji 
\(f\) jest prawdziwe?
\[ 
 f\colon y = \frac{-2(3x + 1)} 
{(2x + 3)(2 - x)}
\]}
{\(f(x) > 0 \iff x\in \left (-\frac{3}
{2};-\frac{1}
{3}\right )\cup (2;\infty )\)}{\(f(x) > 0 \iff x\in \left (-\infty ;-\frac{3}
{2}\right )\cup \left (-\frac{1}
{3};2\right )\)}{\(f(x) > 0 \iff x\in \left (-\frac{3}
{2};2\right )\)}{\(f(x) > 0 \iff x\in \left (-\infty ;-\frac{3}
{2}\right )\cup (2;\infty )\)}{}{}}
\LANGcs{
\Question{
Který z následujících výroků o funkci
\(f\colon y = \frac{-2(3x+1)} 
{(2x+3)(2-x)}\) je
pravdivý?}
{\(f(x) > 0 \iff x\in \left (-\frac{3}
{2};-\frac{1}
{3}\right )\cup (2;\infty )\)}{\(f(x) > 0 \iff x\in \left (-\infty ;-\frac{3}
{2}\right )\cup \left (-\frac{1}
{3};2\right )\)}{\(f(x) > 0 \iff x\in \left (-\frac{3}
{2};2\right )\)}{\(f(x) > 0 \iff x\in \left (-\infty ;-\frac{3}
{2}\right )\cup (2;\infty )\)}{}{}}
\LANGsk{
\Question{
Ktorý z nasledujúcich výrokov o funkcii
\(f\) je pravdivý?
\[ 
 f\colon y = \frac{-2(3x + 1)} 
{(2x + 3)(2 - x)}
\]}
{\(f(x) > 0 \iff x\in \left (-\frac{3}
{2};-\frac{1}
{3}\right )\cup (2;\infty )\)}{\(f(x) > 0 \iff x\in \left (-\infty ;-\frac{3}
{2}\right )\cup \left (-\frac{1}
{3};2\right )\)}{\(f(x) > 0 \iff x\in \left (-\frac{3}
{2};2\right )\)}{\(f(x) > 0 \iff x\in \left (-\infty ;-\frac{3}
{2}\right )\cup (2;\infty )\)}{}{}}
\LANGes{
\Question{
Elige el enunciado verdadero sobre la función 
\(f\).
\[ 
 f(x) = \frac{-2(3x + 1)} 
{(2x + 3)(2 - x)}
\]}
{\(f(x) > 0 \iff x\in \left (-\frac{3}
{2};-\frac{1}
{3}\right )\cup (2;\infty )\)}{\(f(x) > 0 \iff x\in \left (-\infty ;-\frac{3}
{2}\right )\cup \left (-\frac{1}
{3};2\right )\)}{\(f(x) > 0 \iff x\in \left (-\frac{3}
{2};2\right )\)}{\(f(x) > 0 \iff x\in \left (-\infty ;-\frac{3}
{2}\right )\cup (2;\infty )\)}{}{}}
\End

\Data\ID{9000033308}
\Updated{2020-08-22 08:08:32}
\Level{C}
\SubArea{Rational equations and inequalities}
\LANGen{
\Question{
Find the solution set of the following inequality. 
\[ 
 \frac{x^{2} + x + 2} 
{x^{2} + 4x + 3}\geq 0
\]}
{\((-\infty ;-3)\cup (-1;\infty )\)}{\((1;3)\)}{\((-\infty ;1)\cup (3;\infty )\)}{\((-3;-1)\)}{}{}}
\LANGpl{
\Question{
Wyznacz zbiór rozwiązań podanej nierówności.
\[ 
 \frac{x^{2} + x + 2} 
{x^{2} + 4x + 3}\geq 0
\]}
{\((-\infty ;-3)\cup (-1;\infty )\)}{\((1;3)\)}{\((-\infty ;1)\cup (3;\infty )\)}{\((-3;-1)\)}{}{}}
\LANGcs{
\Question{
Určete množinu řešení dané nerovnice. \[\frac{x^{2}+x+2}
{x^{2}+4x+3}\geq 0\]}
{\((-\infty ;-3)\cup (-1;\infty )\)}{\((1;3)\)}{\((-\infty ;1)\cup (3;\infty )\)}{\((-3;-1)\)}{}{}}
\LANGsk{
\Question{
Určte množinu riešení danej nerovnice.
\[ 
 \frac{x^{2} + x + 2} 
{x^{2} + 4x + 3}\geq 0
\]}
{\((-\infty ;-3)\cup (-1;\infty )\)}{\((1;3)\)}{\((-\infty ;1)\cup (3;\infty )\)}{\((-3;-1)\)}{}{}}
\LANGes{
\Question{
Halla el conjunto de soluciones de la siguiente inecuación. 
\[ 
 \frac{x^{2} + x + 2} 
{x^{2} + 4x + 3}\geq 0
\]}
{\((-\infty ;-3)\cup (-1;\infty )\)}{\((1;3)\)}{\((-\infty ;1)\cup (3;\infty )\)}{\((-3;-1)\)}{}{}}
\End

\Data\ID{9100026109}
\Updated{2022-05-29 12:05:45}
\Level{C}
\SubArea{Rational equations and inequalities}
\LANGen{
\Question{
Identify the picture that shows the correct solution set of the following inequality.  In each picture, the set of points corresponding to the solution set is marked in red.
\[ 
 \frac{4} 
{x}\leq x
\]}
{}{}{}{}{}{}}
\LANGpl{
\Question{
Który z poniższych rysunków przedstawia rozwiązanie podanej nierówności?
\[ 
 \frac{4} 
{x}\leq x
\]
Zbiór rozwiązań oznaczono kolorem czerwonym.}
{}{}{}{}{}{}}
\LANGcs{
\Question{
Na kterém obrázku je znázorněno grafické řešení nerovnice
\(\frac{4}
{x}\leq x\)? Řešení je vyznačeno červenou barvou.}
{}{}{}{}{}{}}
\LANGsk{
\Question{
Na ktorom obrázku je znázornené grafické riešenie danej nerovnice?
\[ 
 \frac{4} 
{x}\leq x
\]
Riešenie je vyznačené červenou farbou.}
{}{}{}{}{}{}}
\LANGes{
\Question{
Identifica la imagen que representa la gráfica de la siguiente inecuación. 
\[ 
 \frac{4} 
{x}\leq x
\]
El conjunto de soluciones está en rojo.}
{}{}{}{}{}{}}

\IMGanswer{000261_09-figure3.pdf}
\IMGanswer{000261_09-figure0.pdf}
\IMGanswer{000261_09-figure1.pdf}
\IMGanswer{000261_09-figure2.pdf}\End

\Data\ID{9100026110}
\Updated{2021-01-27 05:01:45}
\Level{C}
\SubArea{Rational equations and inequalities}
\LANGen{
\Question{
Identify the picture which shows the solution of the following equation.
\[ 
 \frac{1} 
{x - 1} = x
\]
The solution set is drawn in a red color.}
{}{}{}{}{}{}}
\LANGpl{
\Question{
Który z poniższych rysunków przedstawia rozwiązanie podanego równania?
\[ 
 \frac{1} 
{x - 1} = x
\] Zbiór rozwiązań oznaczono kolorem czerwonym.}
{}{}{}{}{}{}}
\LANGcs{
\Question{
Na kterém obrázku je znázorněno grafické řešení rovnice
\(\frac{1}
{x-1} = x\)? Řešení je vyznačeno červenou barvou.}
{}{}{}{}{}{}}
\LANGsk{
\Question{
Na ktorom obrázku je znázornené grafické riešenie danej rovnice?
\[ 
 \frac{1} 
{x - 1} = x
\]
Riešenie je vyznačené červenou farbou.}
{}{}{}{}{}{}}
\LANGes{
\Question{
Identifica la imagen que representa la gráfica de la siguiente ecuación.
\[ 
 \frac{1} 
{x - 1} = x
\]
El conjunto de soluciones está en rojo.}
{}{}{}{}{}{}}

\IMGanswer{000261_10-figure0.pdf}
\IMGanswer{000261_10-figure1.pdf}
\IMGanswer{000261_10-figure2.pdf}
\IMGanswer{000261_10-figure3.pdf}\End

\Data\ID{9100033309}
\Updated{2021-01-27 05:01:31}
\Level{C}
\SubArea{Rational equations and inequalities}
\LANGen{
\Question{
Identify the picture which shows the solution of the following inequality.
\[ 
 \frac{1} 
{x} > 1
\] The solution set is drawn in a red color.}
{}{}{}{}{}{}}
\LANGpl{
\Question{
Który z poniższych rysunków przedstawia rozwiązanie podanej nierówności?
\[ 
 \frac{1} 
{x} > 1
\] Zbiór rozwiązań oznaczono kolorem czerwonym.}
{}{}{}{}{}{}}
\LANGcs{
\Question{
Na kterém z nákresů je znázorněno grafické řešení dané nerovnice?
\[\frac{1}
{x} > 1\] Řešení je vyznačeno červenou barvou.}
{}{}{}{}{}{}}
\LANGsk{
\Question{
Na ktorom obrázku je znázornené grafické riešenie danej nerovnice? 
\[ 
 \frac{1} 
{x} > 1
\] Riešenie je vyznačené červenou farbou.}
{}{}{}{}{}{}}
\LANGes{
\Question{
Identifica la imagen que representa la gráfica de la siguiente inecuación.
\[ 
 \frac{1} 
{x} > 1
\] El conjunto de soluciones está en rojo.}
{}{}{}{}{}{}}

\IMGanswer{000333_09-figure2.pdf}
\IMGanswer{000333_09-figure0.pdf}
\IMGanswer{000333_09-figure1.pdf}
\IMGanswer{000333_09-figure3.pdf}\End
